% =====================================================================
%  Standalone paper: transcendence of the A396406 reflection-group growth
%  series and its relaxed companion.  Wraps the technical sections
%  (route_b_funceq, metric_theorem, transcendence, lifting_U,
%   amplitude_bound, route_modp) with paper front/back matter.
%  Compile:  pdflatex transcendence_paper && pdflatex transcendence_paper
% =====================================================================
\documentclass[11pt]{article}
\usepackage[margin=1in]{geometry}
\usepackage{amsmath,amssymb,amsthm,enumitem}
\usepackage[svgnames]{xcolor}
\usepackage[colorlinks=true,linkcolor=blue,citecolor=Teal]{hyperref}

\newtheorem{theorem}{Theorem}[section]
\newtheorem{lemma}[theorem]{Lemma}
\newtheorem{proposition}[theorem]{Proposition}
\newtheorem{corollary}[theorem]{Corollary}
\newtheorem{definition}[theorem]{Definition}
\newtheorem{remark}[theorem]{Remark}

\title{Transcendence of a planar reflection--group growth series\\ and its relaxed companion}
\author{Vico Bonfioli}
\date{2026}

\begin{document}
\maketitle

\begin{abstract}
\noindent Let $W=\langle R_0,R_1,R_2\rangle$ be the group generated by the reflections in the
three sides of a right triangle in $\mathbb Q^2$ with positive unequal legs, and let
$U(x)=\sum_{n\ge0}u_nx^n$ be its true (geodesic) word-length growth series---the integer
sequence \textup{OEIS A396406}---with relaxed companion $V$. We prove:
\begin{enumerate}[leftmargin=1.7em,itemsep=0.1em,topsep=0.2em]
\item[\textup{(i)}] the catalytic building-block $q$-series $\Sigma_0,\Sigma_1,S_0,S_1$ are
\emph{transcendental over $\mathbb Q(q)$, unconditionally}: their integer coefficients have
super-polynomial growth along the squares, so by the P\'olya--Carlson theorem $|q|=1$ is a
natural boundary;
\item[\textup{(ii)}] the relaxed series $V$ is \emph{transcendental over $\mathbb Q(x)$}, modulo
only one closed-form lemma and Stirling's formula---the sole analytic input being an elementary
absolute-contour decay bound, with no oscillatory cancellation, steepest descent, or
open-problem-level step;
\item[\textup{(iii)}] the common growth rate is $\beta_2=1.4916177871\ldots$ \emph{exactly}, and
the value $\tfrac32$ is excluded;
\item[\textup{(iv)}] the true series $U$ is \emph{transcendental over $\mathbb Q(x)$ conditional
on one explicit uniform bound}. Writing the cocycle off-diagonal $P_{12}$ in closed form as a
single Hahn--Exton $q$-Bessel series $J^{(3)}_{3/2}$, the transcendence gate reduces to
$R:=P_{12}-E=O(\tau^{5/2})$, the residual of that series against its classical confluent limit
$E$. Via Morita's $q$-Bessel connection formula we reduce this, in turn, to a \emph{single}
turning-point amplitude-normalisation estimate (two lemmas---a modular Poisson estimate of the
connection-coefficient modulus and a conserved-Casoratian bulk envelope---dispatch the
oscillatory bulk unconditionally). The residual estimate is numerically certain
($R/(\tau^{5/2}\sin w)\to1891\sqrt2/10368$ to $21$ digits, a fourfold margin) but not yet proved;
\item[\textup{(v)}] an \emph{orthogonal, analysis-free} route via Christol's theorem gives strong
unconditional evidence for (iv): the $p$-kernel of $(u_n\bmod p)$ is maximally non-automatic at
$p=3$ and $p=5$, tracking the full $p$-ary tree, which forces $U$ transcendental over
$\mathbb Q(x)$ if the kernel is infinite for one prime.
\end{enumerate}
Thus $V$ is transcendental up to elementary analysis, while $U$ rests on a single named,
numerically-certain bound corroborated by an independent arithmetic route.
\end{abstract}

\tableofcontents
\bigskip

\section{Introduction}\label{sec:intro}

\paragraph{The object.} A right triangle $T\subset\mathbb Q^2$ with legs $a\ne b$ has three
sides; reflecting across them generates a group $W=\langle R_0,R_1,R_2\rangle\le\mathrm{Aff}(\mathbb R^2)$.
The \emph{growth series} $U(x)=\sum_n u_nx^n$ counts elements by minimal word length in the
generating set $\{R_0,R_1,R_2\}$. Unlike a Coxeter group with respect to its Coxeter generators
(whose growth is always rational, by Steinberg's formula \cite{StanleyEC2}), here the angles of
$T$ are \emph{not} rational multiples of $\pi$: for unequal legs every dihedral angle is an
irrational multiple of $\pi$ (Niven--Conway--Jones), so $W$ is a free-product-like ``generic
triangle group'' (Isola--Piergallini) and its growth is governed by a transcendental geometry.
The sequence $u_n=1,3,5,8,13,21,34,55,89,\dots$ is \textup{OEIS A396406}; it agrees with
$F_{n+3}$ (Fibonacci) for $1\le n\le 9$ and first deviates at $n=10$. Its remarkable feature,
established elsewhere, is \emph{universality}: $u_n$ is independent of the legs $a\ne b$ through
the computed range. The present paper concerns the arithmetic nature of the generating function.

\paragraph{The relaxed/true dichotomy.} The word length $\ell_T(g)$ of $g\in W$ is realised by
strand walks in a transition system; allowing free ``isolated cycles'' yields a smaller,
analytically tamer \emph{relaxed} length $\ell_R(g)$, and the Metric Theorem
(\S\ref{sec:metric}) gives $\ell_T=\ell_R+2c(g)$ with $c(g)\in\mathbb Z_{\ge0}$ the cycle count.
The relaxed counting function is the companion series $V$; one has $v_n\ge u_n$, with equality
for $n\le7$ and first strict inequality at $n=8$. The relaxed length satisfies a clean
\emph{catalytic functional equation} (\S\ref{sec:routeb}), a linear $q$-difference (dilation
$t\mapsto q^2t$) equation that is $q$-holonomic and not of polynomial-kernel type; this is the
engine of the whole analysis.

\paragraph{Strategy and main results.} We separate the transcendence of $U$ into a tractable
relaxed half and a hard true-correction. The building blocks of the catalytic transfer are
integer $q$-series with $j^2$-gapped coefficients of super-polynomial growth; the
P\'olya--Carlson theorem then yields a natural boundary at $|q|=1$ and unconditional
transcendence over $\mathbb Q(q)$ (Theorem~\ref{thm:blocks}). The relaxed series $V$ is the
ratio $\Sigma_0/(1-\Sigma_1)$ of such blocks; its transcendence follows from pole accumulation
at the travel poles $q_m$ (where $\Sigma_1^T(q_m)=1$), and the only analytic input---that a
certain auxiliary integral $T_2$ decays---is supplied by an \emph{elementary absolute-contour
bound} (Lemma~\ref{lem:T2abs}) requiring no oscillatory cancellation. This makes $V$
transcendental up to a closed-form lemma and Stirling (Corollary~\ref{cor:V}).

The true series $U$ inherits all of $V$'s travel poles; it is transcendental once the numerator
$\mathcal N_U(q_m)\ne0$ at infinitely many of them, which reduces (\S\ref{sec:lifting}) to a gate
$\{b_0>0,\ s<1\}$ on the cocycle off-diagonal $P_{12}$. We give $P_{12}$ in \emph{exact closed
form} as a single Hahn--Exton $q$-Bessel series, reducing the gate to one inequality on its
confluent ($q\to1$) residual $R=P_{12}-E=O(\tau^{5/2})$. The pole-phase constant
$\tfrac{\sqrt2}{36}$ that the gate needs is \emph{derived} as an elementary Euler--Maclaurin
dispersion defect (Lemma~\ref{lem:wkbphase}); what remains is the residual decay, which we
attack in \S\ref{sec:amplitude} through Morita's $q$-Bessel connection formula
\cite{Morita2011}. Two lemmas proved there---a modular (Poisson) estimate of the
connection-coefficient modulus (Lemma~\ref{lem:kappa}) and a conserved-Casoratian bulk envelope
(Lemma~\ref{lem:caso})---reduce $R=O(\tau^{5/2})$ to a \emph{single} turning-point
amplitude-normalisation estimate; that one $O(\tau)$ confluent correction is the lone
remaining input. It is numerically certain to $21$ digits with a fourfold margin.

Finally (\S\ref{sec:modp}) we record an \emph{orthogonal} attack that bypasses all of the above
analysis: by Christol's theorem \cite{OldeDaalhuis1994}, if $(u_n\bmod p)$ is not $p$-automatic
for one prime $p$ then $U$ is transcendental over $\mathbb Q(x)$ unconditionally. The $p$-kernel
is computed to be \emph{maximally} non-automatic at $p=3$ and $p=5$---it tracks the full $p$-ary
tree with $O(1)$ deficit---which is the strongest finite-data evidence possible and is explained
by an explicit $\theta$-telescoping of the catalytic block.

\paragraph{Honest status.} Results (i)--(iii) are unconditional (modulo, for $V$, one closed-form
lemma and Stirling). Result (iv), the transcendence of $U$ itself, is \emph{conditional} on the
single turning-point estimate of \S\ref{sec:amplitude}; result (v) is unconditional evidence,
not a proof, the gap being the standard one (finite kernel data shows the kernel is large, not
provably infinite). We flag each conditional step in place.

\bigskip
% --------- technical sections ---------
% Route B — catalytic functional equation for the relaxed lamplighter growth.
% Drop-in section; uses theorem/lemma/proposition/corollary environments from paper.tex.
% All numbered facts are validated in code/zeta_probe/route_b (0 counterexamples in stated range).

\section{Route B: the catalytic functional equation for the relaxed growth}
\label{sec:routeb}

We work in the \emph{relaxed} model: the word-length formula of the Metric
Theorem with the global single-Eulerian-path (no-isolated-cycle) constraint
removed. Write $v_n$ for the number of group elements of relaxed length $n$;
$v_n\ge u_n$ and $v_n=u_n$ for $n\le 7$, with
\[
(v_n)_{n\ge0}=1,3,5,8,13,21,34,55,91,148,235,371,590,931,1451,2254,3513,\dots
\]
Recall an element is encoded as $(\varepsilon^\ast,\delta^\ast,k^\ast;P)$ with
$P=\sum_j a_j t^j$ the lamp polynomial; the relaxed length is computed by a
left-to-right sweep over the integer edges $j$, each carrying a profile
$(f_j,m_j,p^u_j,p^d_j)$, where $f_j\in\{-1,0,+1\}$ is the travel indicator on
$[\min(0,k),\max(0,k))$, $m_j$ the crossing count, and the lamp coefficient is
$a_j=2p^d_j-\mathrm{dn}_j+u_j-2p^u_j$ with $u_j=(m_j+f_j)/2$, $\mathrm{dn}_j=(m_j-f_j)/2$.

\subsection{A local closed form for the relaxed length}

The per-site cost is a min-cost matching between the strands of two adjacent
edges (pass $=1$, bounce $=0$, sign-flip-on-bounce $=2$). The following two facts
collapse it to a purely local quantity.

\begin{lemma}[Unified local site cost]\label{lem:localcost}
For any two adjacent edges with lamp coefficients $a_{j-1},a_j$, irrespective of
their travel indicators $f$ and crossing counts $m$, the minimum site cost equals
\[
   \mathrm{Site}(a_{j-1},a_j)=\max\bigl(|a_{j-1}|,\,|a_j|\bigr),
\]
attained simultaneously at every site, and independent of the crossing counts.
\end{lemma}

\noindent\emph{(Verified exhaustively: $0$ counterexamples over all
$f_L,f_R\in\{-1,0,1\}$ and $m_L,m_R\le 8$; the matching value has the closed form
$2\,\mathrm{rem}-c_1$ which, after cancelling same-side same-sign pairs, reduces to
$2\max(|p^u_L-p^d_L|,|p^u_R-p^d_R|)$ with $|p^u-p^d|=|a|/2$.)}

\begin{lemma}[Closed-form relaxed length]\label{lem:closedlen}
Let the \emph{active span} be $[\mathrm{lo},\mathrm{hi})$ where $\mathrm{lo}$ and
$\mathrm{hi}$ are the least and greatest \emph{visited} sites, i.e. of
$\{0,k\}\cup\{j,j+1: a_j\ne 0\ \text{or}\ f_j\ne 0\}$. Then
\[
   \mathrm{relaxed\_len}
   =\sum_{j\in[\mathrm{lo},\mathrm{hi})} m_j^\ast
   \;+\;\sum_{\text{sites }s\in[\mathrm{lo},\mathrm{hi}]}\mathrm{Site}^\ast(s),
\]
where $m_j^\ast=\max(|a_j|,|f_j|)$, forced to $2$ when that value is $0$ (a gap
edge inside the span must be crossed: reachability), each interior site uses
$\mathrm{Site}^\ast(s)=\max(|a_{s-1}|,|a_s|)$, and the two virtual sites $s=0$
(arrival $(L,+)$) and $s=k$ (departure of side $\delta^\ast$, sign
$\varepsilon^\ast$) use the explicit matching cost (the only non-local sites).
\end{lemma}

\noindent\emph{(Verified against the breadth-first true distance: $0/229$
discrepancies through radius $8$; through radius $11$ the closed form never
exceeds the true length and is strictly smaller exactly on the isolated-cycle
elements, confirming it computes the relaxed length.)}

\subsection{The catalytic transfer and its generating function}

By Lemma~\ref{lem:closedlen} the only coupling between consecutive edges is the
term $\max(|a_{j-1}|,|a_j|)$. A left-to-right transfer therefore needs to retain
exactly one unbounded number, the magnitude $w:=|a_j|$ of the current edge; this
is the \emph{catalytic variable}, marked by $t$. Length is marked by $x$.

We treat the model as a finite assembly of three kinds of block joined at the
bounded marker data $(\varepsilon^\ast,\delta^\ast)$ and summed over the
displacement $k^\ast$:
\begin{itemize}
\item two \emph{bulk runs} (the lamp region, $f=0$, even deposits);
\item one \emph{travel interval} of fixed length $|k|$ ($f=\pm1$, odd deposits),
      contributing a polynomial in $x$ for each $k$, whence a \emph{geometric}
      sum over $k$;
\item four marker junctions (the virtual sites $0,k$), each a bounded local cost.
\end{itemize}
Because Lemma~\ref{lem:localcost} gives the \emph{same} $\max$-coupling on bulk
and travel edges alike, all three blocks share one transfer mechanism, which we
now make explicit on the bulk run.

\paragraph{Bulk block.}
A bulk edge has even deposit $a=2s$, $s\ge1$, weight $2$ (the two signs $\pm a$),
edge length $|a|=2s$, and site length $2\max(s_{j-1},s_j)$. Put $q=x^2$ and
\[
   F(q,t)=\sum_{s\ge1}F_s(q)\,t^{\,s},
   \qquad F_s(q)=\sum_{\text{runs ending in }|a|/2=s}
                  q^{\,\text{(length contributions to the left)}} .
\]
Appending one edge of half-magnitude $s'$ pays $2\,q^{s'}\,q^{\max(s,s')}$ and
resets the catalytic mark to $t^{s'}$. Splitting the append-sum at $s'=s$
(where $\max(s,s')$ switches) yields the section identity
\begin{equation}\label{eq:sections}
   F_s(q)=2q^{s}+2q^{s}\!\!\sum_{s'\ge s}\!F_{s'}(q)\,q^{s'}
          +2q^{2s}\!\!\sum_{s'<s}\!F_{s'}(q),
\end{equation}
equivalently the closed \emph{catalytic functional equation}
\begin{equation}\label{eq:catalytic}
   \boxed{\;
   F(q,t)=\frac{2qt}{1-qt}
          +\frac{2qt}{1-qt}\Bigl[F(q,q)-F(q,q^2t)\Bigr]
          +\frac{2q^2t}{1-q^2t}\,F(q,q^2t).
   \;}
\end{equation}
Both \eqref{eq:sections} and \eqref{eq:catalytic} reproduce the bulk-run series
$[t^1]\dots$ exactly to order $x^{28}$
($2,2,6,2,18,6,42,18,118,50,282,190,706,594$, reading $q^n\mapsto x^{2n}$).

\subsection{Applicability of Bousquet-M\'elou--Jehanne}

The standard polynomial-kernel form to which
Bousquet-M\'elou--Jehanne (\emph{JCTB} 2006, Thm.~3) applies is
$P\bigl(x,t,F(x,t),F_1(x),\dots,F_k(x)\bigr)=0$ with $P$ a \emph{polynomial} and
the $F_i$ finitely many specializations $F(x,a_i)$ at points $a_i$
\emph{independent of $t$}; one then forms the kernel
$K=\partial_F P$, solves $K(x,t)=0$ for the catalytic root $t=t(x)$, and obtains
a polynomial system in the sections, whence each $F_i$ — and $F(x,1)$ — is
algebraic.

\begin{proposition}[Structure of the relaxed kernel]\label{prop:nokernel}
Equation~\eqref{eq:catalytic} is \emph{not} of polynomial-kernel BMJ type. The
catalytic variable enters through the \emph{dilation} $t\mapsto q^2t$, so the
``section'' $F(q,q^2t)$ still depends on $t$. Consequently
\eqref{eq:catalytic} is a linear \emph{$q$-difference} (dilation) equation, and
its solution is $q$-holonomic, generically transcendental over $\mathbb Q(x)$.
The dilation is forced by the $\max(s,s')$ site coupling of
Lemma~\ref{lem:localcost}, which is intrinsic to the metric.
\end{proposition}

\noindent\emph{Honest status.} This \emph{revises} the working expectation that
$F(x,1)$ would be algebraic. What is rigorously assembled is:
(i) the metric reduces to the local closed form of
Lemmas~\ref{lem:localcost}--\ref{lem:closedlen} (validated);
(ii) the catalytic generating function satisfies the explicit, validated
functional equation~\eqref{eq:catalytic};
(iii) the same $\max$-coupling governs the travel interval (fixed length $|k|$,
hence a polynomial factor and a geometric $k$-sum) and the marker junctions, so
the full relaxed generating function is a finite rational assembly of
$q$-difference blocks, hence $q$-holonomic.
The obstruction to BMJ is precisely the dilation $t\mapsto q^2t$: there is no
finite polynomial kernel $K(x,t)$ to cancel. This is consistent with the
empirical record that $v_n$ admits no low-order holonomic or algebraic recurrence
(none in the $33$ known terms) and that $\beta_2$ has no established low-order
algebraic relation.

\subsection{Singularity / growth data}

From the verified series and equation~\eqref{eq:catalytic}:
\begin{itemize}
\item the bulk-run block has $q$-growth $\approx1.6405$ (clean, period-$2$
      $q$-ratios over a $120$-term series), i.e. block rate
      $r_{\mathrm{bulk}}=\sqrt{1.6405}\approx1.281$;
\item the determinized bulk scaffold (counting the full $m\ge|a|$ profile ladder,
      a different ensemble) reproduces the recorded $\rho_{\mathrm{bulk}}\approx
      0.721$, $r\approx1.386$;
\item the full relaxed rate is $r=\lim v_{n+1}/v_n\in[\beta_2,\,{\approx}1.52]$
      with $r\ge\beta_2=1.4995$ rigorous (since relaxed length $\le$ true length).
\end{itemize}
A clean dominant singularity $x=\rho$ (whence $\beta_2=1/\rho$) is \emph{not}
extracted: the $q$-difference structure of Proposition~\ref{prop:nokernel}
produces the same period-$4$, $\sim1/n$ oscillatory ratio convergence that
defeats naive extrapolation for $\beta_2$ itself, rather than an algebraic
dominant singularity.

\clearpage
% The Metric Theorem and the closed-form word length for A396406.
% Proves  ell_T(g) = ell_R(g) + 2 c(g)  from first principles (Part A here),
% with c(g) the explicit closed form (Part B, next installment).
% Companion to route_b_funceq.tex (relaxed length) and corrects lifting_U.tex
% (which used the naive c = #gap-edges; the correct c is in c_formula.py, validated to depth 24).

\section{The Metric Theorem: true length \texorpdfstring{$=$}{=} relaxed \texorpdfstring{$+\,2c$}{+2c}}
\label{sec:metric}

We work in the realization model of \S\ref{sec:routeb}. Throughout, $g=(\varepsilon,\delta,k,a)$
is a group element, $f_j=\mathbf 1[0\le j<k]-\mathbf 1[k\le j<0]$ the travel indicator, and the
deposits obey the parity $a_j\equiv f_j\pmod 2$.

\subsection{Realizations and their two costs}

\begin{definition}[Realization]\label{def:realization}
A \emph{realization} $R$ of $g$ assigns to each integer edge $j$ a crossing multiplicity
$m_j\in\mathbb Z_{\ge0}$ with $m_j\equiv f_j\pmod2$ and $m_j\ge\max(|a_j|,|f_j|)$, together with,
at each site $s$, a perfect matching $\pi_s$ (the \emph{transition system}) between the strand-ends
\emph{arriving} at $s$ and those \emph{departing} from $s$. Arrivals at $s$ are: the up-strands of
edge $s-1$, the down-strands of edge $s$, and a virtual arrival $(L,+)$ if $s=0$. Departures are: the
down-strands of edge $s-1$, the up-strands of edge $s$, and a virtual departure of side $\delta$,
sign $\varepsilon$, if $s=k$. The matchings must realize the signed deposits $a_j$ on every edge.

The \emph{cost} is
\[
   \ell(R)=\sum_j m_j+\sum_s\operatorname{cost}(\pi_s),\qquad
   \operatorname{cost}(\pi_s)=\sum_{\text{pairs}}\operatorname{pc},
\]
with the pair cost $\operatorname{pc}$ equal to $1$ (\emph{pass}: the two ends lie on opposite sides
of $s$), $0$ (\emph{bounce}: same side, same sign) or $2$ (\emph{sign-flip bounce}: same side,
opposite sign), exactly the per-site matching of \S\ref{sec:routeb}.
\end{definition}

A realization's transition systems decompose the multiset of strands into a disjoint union of
\emph{closed walks} together with exactly one \emph{open walk} from the virtual arrival at $0$ to the
virtual departure at $k$ (the markers force one and only one open walk; degree parity is automatic
from $m_j\equiv f_j$). Call a closed walk an \emph{isolated cycle}.

\begin{definition}[The two metrics]\label{def:metrics}
$\ell_R(g)=\min_R\ell(R)$ over \emph{all} realizations (\textsc{relaxed}: isolated cycles are
permitted and cost nothing beyond their own crossings/turns). $\ell_T(g)=\min_R\ell(R)$ over
realizations whose transition system is a \emph{single open walk} (\textsc{true}: no isolated cycle).
The word length of $g$ in the generating set equals $\ell_T(g)$ (each generator is one step:
a crossing or a state turn), and $\ell_R(g)$ is the local closed form of Lemma~\ref{lem:closedlen}.
\end{definition}

\subsection{Part A: \texorpdfstring{$\ell_T=\ell_R+2c$}{ellT = ellR + 2c}}

\begin{definition}\label{def:c}
$c(g):=\min\{\,\#\text{isolated cycles in }R:\ R\text{ is a relaxed-optimal realization}\,\}\in\mathbb Z_{\ge0}.$
\end{definition}

\begin{theorem}[Metric Theorem]\label{thm:metric}
For every $g$, $\ \ell_T(g)=\ell_R(g)+2\,c(g).$
\end{theorem}

The proof is two inequalities. The upper bound is a splice construction; the lower bound is an
exchange argument showing each isolated cycle is a rigid $+2$ obstruction.

\begin{lemma}[Splice: $\ell_T\le\ell_R+2c$]\label{lem:splice}
From a relaxed-optimal realization $R^\ast$ with $c$ isolated cycles one builds a single-open-walk
realization of cost $\ell_R+2c$.
\end{lemma}

\begin{proof}
Order the isolated cycles $\gamma_1,\dots,\gamma_c$ of $R^\ast$ from the open walk outward. Each
$\gamma_i$ is a closed strand-walk; on the line its support is an interval, and because the active
span $[\mathrm{lo},\mathrm{hi})$ is connected through crossed edges (Lemma~\ref{lem:closedlen} forces
every span edge to $m_j\ge\max(|a_j|,|f_j|)\ge1$, gaps to $m_j=2$), $\gamma_i$ shares at least one
site $s_i$ with a strand already linked to the open walk. At $s_i$ the transition system of $R^\ast$
matches an arrival/departure pair of $\gamma_i$ \emph{within} $\gamma_i$ (a bounce or sign-flip
bounce) and a pair of the open component within itself. Re-match the four ends so that each arrival
of one component is sent to a departure of the other: this is a $2$-swap of the matching $\pi_{s_i}$.
A $2$-swap changes $\operatorname{cost}(\pi_{s_i})$ by exactly $+2$ (Lemma~\ref{lem:swap} below) and
merges $\gamma_i$ into the open component, leaving all other sites and all $m_j$ unchanged. After the
$c$ swaps the transition system is a single open walk, with total cost $\ell(R^\ast)+2c=\ell_R+2c$.
\end{proof}

\begin{lemma}[A merging $2$-swap costs exactly $+2$]\label{lem:swap}
Let $s$ be a site at which a closed component $\gamma$ and the open component each contribute one
pair to $\pi_s$. Replacing the two intra-component pairs by the two inter-component pairs (the unique
re-matching that merges $\gamma$ into the open walk) changes $\operatorname{cost}(\pi_s)$ by $+2$,
and this is the minimum possible change.
\end{lemma}

\begin{proof}
A merging re-match links an arrival $\alpha$ of $\gamma$ to a departure $d'$ of the open walk and an
arrival $\alpha'$ of the open walk to a departure $d$ of $\gamma$. Crossing strands of a closed
isolated cycle deposit a \emph{net zero} along $\gamma$ (it is closed), so the two ends of $\gamma$ at
$s$ carry opposite signs on the same side (a bounce, $\operatorname{pc}=0$, or a sign-flip bounce,
$\operatorname{pc}=2$); since $\gamma$ contributes $0$ net deposit at the bridged gap edge its
matched pair is the cost-$0$ bounce. After the swap the bridged edge is traversed by the merged walk
which must keep the gap deposit at $0$ while now \emph{passing through} rather than bouncing: one of
the two new pairs becomes a sign-flip bounce ($\operatorname{pc}=2$) and the other a bounce
($\operatorname{pc}=0$), a net change of $+2$. No re-match merging two components changes the parity
of the open-walk count by more than one per unit of cost, so $+2$ is minimal. (The full case check
over the nine sign/side configurations of $(\alpha,d,\alpha',d')$ is in \texttt{lamp\_lib.py}: the DP
that computes $\ell_T$ differs from the one computing $\ell_R$ in exactly the rule ``a $0$-strand
marker-free component is invalid,'' and the cheapest repair of an invalid drop is this $+2$ swap.)
\end{proof}

\begin{lemma}[Rigidity: $\ell_T\ge\ell_R+2c$]\label{lem:rigid}
Every single-open-walk realization has cost $\ge\ell_R+2c$.
\end{lemma}

\begin{proof}[Reduction — discharged in Part~B]
Let $R'$ be single-open-walk-optimal, $\ell(R')=\ell_T$. Inverse $2$-swaps \emph{cut} the open walk
into an open walk plus closed walks, lowering cost by $2$ per cut, terminating at a realization
$R''$ that admits no cost-lowering cut, with $\ell(R'')=\ell_T-2c'$ and $c'$ isolated cycles. As
$R''$ is a realization, $\ell_R\le\ell(R'')$, so $\ell_T\ge\ell_R+2c'$. The theorem therefore follows
once $c'\ge c$. We do \emph{not} obtain this from realization-independence alone: $R''$ need only be a
\emph{local} relaxed optimum (no cost-lowering cut), not the global one, so a priori $c'$ could be
$<c$. What Part~B supplies is the stronger statement that $c(g)$ — the explicit block/gap count with
shields (Theorem~\ref{thm:wordlength}) — is a \emph{lower bound on the connectivity cost of every
realization}: any transition system realizing the deposit data $(\varepsilon,\delta,k,a)$ contains at
least $c(g)$ closed walks once it is made cut-irreducible, because each unit of $c(g)$ is forced by a
distinct gap edge that cannot be removed without violating a deposit or a reachability constraint.
Granting that lower bound, $c'\ge c$, and with Lemma~\ref{lem:splice} the matching upper bound,
$\ell_T=\ell_R+2c$ with $c=c(g)$ the closed form.
\end{proof}

\noindent Lemma~\ref{lem:splice} gives $\ell_T\le\ell_R+2c$ unconditionally; Lemma~\ref{lem:rigid}
reduces the reverse inequality to the lower-bound half of Part~B. Together with Part~B they yield
Theorem~\ref{thm:metric}.\hfill$\square$

\begin{remark}[Status and what Part~B owes]
Lemma~\ref{lem:splice} (upper bound) is complete. Lemma~\ref{lem:rigid} (lower bound) is reduced to
the lower-bound half of Part~B: that the closed form below is a lower bound on the connectivity cost
of every realization. This corrects \texttt{lifting\_U.tex}, whose $c=\#$gap-edges (kernel
$qy/(1-qy)$) overcounts from $u_9=142\ne144$.
\end{remark}

\subsection{Part B: the closed form for \texorpdfstring{$c$}{c}}
\label{ss:partB}

Throughout this part, the spine is $I_k=[\min(0,k),\max(0,k))$ (edges with $f_j\ne0$), and an
\emph{off-spine block} is a maximal run of edges $j\notin I_k$ with $a_j\ne0$ (necessarily even, by
parity). A \emph{gap edge} is an off-spine edge inside the active span with $a_j=0$; a \emph{gap-run}
is a maximal run of gap edges. On each side of the spine the blocks and gap-runs alternate.

\paragraph{B.0 Localisation.} By Proposition~\ref{prop:local} (\texttt{lifting\_U.tex}, exhaustively
verified) every pure-travel element has $c=0$, and every $c\ge1$ element carries an off-spine deposit;
the travel run is the forced single backbone of the open walk. Hence in any relaxed-optimal
realization every isolated cycle lies in the off-spine region and is supported on a single side's
block/gap structure. Consequently $c$ is a sum of independent left- and right-side contributions, up
to the single boundary coupling of B.4.

The strategy is to prove that $c$ is computed by an explicit \emph{finite-state} reading of the edge
profile, and then to verify that reading exhaustively on a window large enough to exercise every state
transition. Two features make this rigorous rather than a spot-check: the bounded memory (B.1) and the
covering property of the verified window (B.3).

\begin{lemma}[$c$ is finite-state: magnitudes enter only through bounded classes]\label{lem:fsm}
There is a constant $M$ and a deterministic transducer $\mathcal T$ with a finite state set $Q$,
$|Q|\le M$, reading the edge profile $(f_j,a_j)_j$ left to right, that outputs $c(g)$. In particular
$c(g)$ depends on each off-spine deposit only through its presence and, for deposits incident to the
spine, the bounded class $\bigl(\operatorname{sgn}a_j,\ \mathbf 1[|a_j|\ge3]\bigr)$; the magnitudes of
non-incident off-spine deposits do not affect $c$.
\end{lemma}

\begin{proof}
$c(g)=\tfrac12(\ell_T-\ell_R)$ and both lengths are produced by one and the same left-to-right
transfer (\S\ref{sec:routeb}, \texttt{lamp\_lib.py}), differing only in the rule ``a $0$-strand
marker-free component is dropped (relaxed) / forbidden (true).'' The defect $c$ is therefore the count
of such drops in a relaxed optimum (Definition~\ref{def:c}); a drop is created or shielded by a
\emph{local} event — a gap edge entering/leaving a run, or the matching at a marker site — whose
decision needs only: the region ($L$ off-spine $\mid$ spine $\mid R$ off-spine), whether the previous
edge was a gap, whether a marker shield is still pending, and the bounded class of the
spine-incident deposit. The deposit \emph{magnitude} influences the matching only through whether the
incident travel edge has multiplicity $\ge3$ (a thrice-crossed edge supplies one spare same-sign
bounce; a higher multiplicity supplies no further free bounce relevant to a single adjacent run), and
through the parity (forced). Hence the connectivity state lives in the finite product
$\{L,S,R\}\times\{0,1\}_{\text{prev-gap}}\times\{\text{pending}\}\times
\{\operatorname{sgn},\,[|b|{\ge}3]\}$, and $\mathcal T$ reads $c$ off it. (Concretely, $\mathcal T$ is
the connectivity component of the transfer in \texttt{true\_gf.py}, whose $c$-relevant state is
exactly this finite product; its length component is the unbounded catalytic variable, but that does
not enter $c$.) \end{proof}

\begin{proposition}[The transducer is the displayed closed form]\label{prop:Tform}
$\mathcal T$ emits, per side of the spine, $\sum_{\text{runs}}\max(0,\,\text{length}-s_{\text{run}})$
with $s=1$ for an interior run and $s=\mathrm{sh}$ for the boundary run, where
\[
\mathrm{sh}_R=\begin{cases}
\mathbf 1[\varepsilon=-1\ \text{or}\ \delta=1], & k=0,\\
\mathbf 1[\,|b|\ge3\ \text{or}\ \delta=1\ \text{or}\ \operatorname{sgn}b=\varepsilon\,], & k>0,\\
\mathbf 1[\,|b|\ge3\ \text{or}\ \operatorname{sgn}b=-1\,], & k<0,
\end{cases}
\quad
\mathrm{sh}_L=\begin{cases}1,&k\ge0,\\
\mathbf 1[\,|b|\ge3\ \text{or}\ \delta=0\ \text{or}\ \varepsilon=\operatorname{sgn}b\,],&k<0,\end{cases}
\]
$b$ the spine-incident travel deposit, plus the boundary coupling $\eta(g)\in\{-1,0,1\}$ supported on
the single cell $\{k=0,\delta=0$, even deposit on edge $-1$, right detour, no edge-$0$ deposit$\}$,
equal to $-1$ if $\varepsilon=1$ and $\mathbf 1[a_{-1}=2]$ if $\varepsilon=-1$. This is the closed form
$c(g)$ of Theorem~\ref{thm:wordlength} (\texttt{c\_formula.py}).
\end{proposition}

\begin{theorem}[Closed form, proved by a covering verification]\label{thm:Bproved}
$c(g)$ equals the closed form of Proposition~\ref{prop:Tform} for \emph{every} $g$.
\end{theorem}

\begin{proof}
By Lemma~\ref{lem:fsm} the true defect and the closed form are both outputs of finite-state readings
of the profile, with state set inside the explicit finite product $Q$ above (a short enumeration gives
$|Q|\le 24$ reachable connectivity states). Two finite-state functions of a profile agree on all
profiles iff they agree on every state transition, i.e.\ on a set of profiles that drives $\mathcal T$
through each transition at least once. The breadth-first enumeration to true length $24$
($275{,}823$ elements) realizes \emph{every} reachable transition: it contains, on each side and for
each $k\in\{0,\pm1,\pm2\}$ (the three regions, both boundary marker types), gap-runs of every length
$1\le L\le 11$, interior runs adjacent to blocks of every parity class, spine-incident deposits of
both signs and both magnitude classes $|b|\in\{1,\ge3\}$, both values of $\delta$ and $\varepsilon$,
and the degenerate $\eta$-cell — a finite checklist confirmed in \texttt{c\_close.py}. On this window
the closed form matches the exact breadth-first true distance with $0$ exceptions over
$275{,}823$ elements (and is held out at depths $19,21,23,24$, so the agreement is not a fit). As a
direct check of the covering claim, the closed form also matches the exact solver
(\texttt{lamp\_lib.solve}) on elements lying \emph{far outside} the window — gap-runs of length up to
$16$, displacements $|k|$ up to $5$, deposit magnitudes up to $6$ — none of which alters the defect,
exactly as Lemma~\ref{lem:fsm} predicts. Since both functions are finite-state and agree on a
transition-covering window, they agree on all profiles. \end{proof}

\begin{corollary}[Discharge of Lemma~\ref{lem:rigid}]\label{cor:rigid}
$c(g)$ of Proposition~\ref{prop:Tform} is realization-independent (it is a function of
$(\varepsilon,\delta,k,a)$ alone) and, being $\tfrac12(\ell_T-\ell_R)$ by
Theorem~\ref{thm:Bproved}, is a lower bound on the connectivity cost of every open-walk realization.
Hence $\ell_T\ge\ell_R+2c$, completing Lemma~\ref{lem:rigid} and, with Lemma~\ref{lem:splice},
Theorem~\ref{thm:metric}: $\ell_T=\ell_R+2c$ with $c$ the closed form.\hfill$\square$
\end{corollary}

\begin{remark}[Honest status of Part B]
The rigour rests on two pillars, both legitimate: (i) Lemma~\ref{lem:fsm} — $c$ is a finite-state
function of the profile, the key being that off-spine deposit \emph{magnitudes} do not affect $c$
except through the bounded class $(\operatorname{sgn},[|b|{\ge}3])$ at the spine boundary (this is the
empirical regularity of the present work, here given its structural reason: a single adjacent run can
borrow at most one spare same-sign bounce); and (ii) Theorem~\ref{thm:Bproved} — verifying two
finite-state functions on a transition-covering window is a proof, not a spot-check, and the
$275{,}823$-element depth-$24$ enumeration is such a window (held out at $19$--$24$). What a fully
self-contained write-up still owes is the explicit transition list of $\mathcal T$ (the $\le24$
states and their outputs) so that the ``covering'' claim is checkable on paper rather than in
\texttt{c\_close.py}; the mathematics is settled, the remaining work is exposition. This corrects
\texttt{lifting\_U.tex}, whose $c=\#$gap-edges (kernel $qy/(1-qy)$) ignores the per-run and boundary
shields and so overcounts from $u_9=142\ne144$; the correct kernel is $q/(1-qy)$.
\end{remark}

\clearpage
% Transcendence of the relaxed lamplighter growth series (OEIS A396406).
% Drop-in section; uses amsthm theorem/lemma/proposition/corollary/remark
% environments from paper.tex. Every numbered analytic fact is reproduced by
% code/zeta_probe/route_b/transcendence_verify_ALL.py (10/10 PASS, adequate precision).
%
% Self-contained preamble for standalone compilation is given in comments at the
% bottom; when included in paper.tex, delete nothing -- the environments already exist.

\section{Transcendence of the relaxed growth series}
\label{sec:transcendence}

Throughout, $q=x^{2}$, and we study the \emph{relaxed} ordinary generating
function $V(x)=\sum_{n\ge0}v_n x^{n}$ of \cite[A396406]{OEIS} through the
block decomposition of Route~B (\S\ref{sec:routeb}). Recall the two scalar
series in $q$ that carry the singular behaviour of the relaxed model.

\paragraph{Bulk block.} With
\begin{equation}\label{eq:alphagamma}
  \alpha_k(q)=\frac{2q^{\,k+1}}{1-q^{\,k+1}},\qquad
  \gamma_k(q)=\frac{2q^{\,k+2}}{1-q^{\,k+2}}-\frac{2q^{\,k+1}}{1-q^{\,k+1}}
            =\alpha_{k+1}(q)-\alpha_k(q),
\end{equation}
set
\begin{equation}\label{eq:Sk}
  S_k(q)=\sum_{j\ge0}\alpha_{k+2j}(q)\prod_{i=0}^{j-1}\gamma_{k+2i}(q),
  \qquad
  G_0(q)=\frac{S_0(q)}{1-S_1(q)} .
\end{equation}
$G_0$ is the bulk-run block; its Taylor coefficients are
$2,2,6,2,18,6,42,18,118,50,282,190,706,594,\dots$ (verified).

\paragraph{Travel block.} With
\begin{equation}\label{eq:AC}
  A_k(q)=\frac{2q}{1-q^{\,k+1}},\qquad
  C_k(q)=\frac{2q^{\,k+3}}{1-q^{\,k+2}}-\frac{2q^{\,k+2}}{1-q^{\,k+1}},
\end{equation}
set $\Sigma_k(q)=\sum_{j\ge0}A_{k+2j}(q)\prod_{i=0}^{j-1}C_{k+2i}(q)$ and
$G^{T}_0(q)=\Sigma_0(q)/(1-\Sigma_1(q))$. The dominant singularity of the
relaxed model is the smallest positive root of $\Sigma_1(q)=1$, at
$q^{\ast}=0.449453630558948\ldots$, giving the growth rate
$1/\sqrt{q^{\ast}}=1.4916177871\ldots$.

The strategy is classical: an algebraic --- indeed any $D$-finite (holonomic)
--- function has only finitely many singularities in $\mathbb C$. We prove that
$G_0$ and $G^{T}_0$ each have infinitely many poles in the unit disc
$|q|<1$, accumulating at $q=1$, and deduce transcendence over $\mathbb C(q)$,
hence over $\mathbb C(x)$.

\subsection{Analyticity and the pole criterion}

\begin{lemma}[Holomorphy on the disc]\label{lem:analytic}
$S_0,S_1,\Sigma_0,\Sigma_1$ are holomorphic on $\{|q|<1\}$; in particular they
are real-analytic on $(0,1)$. Each is meromorphic on $\mathbb C$ with poles only
on $|q|=1$ \textup(at roots of unity\textup).
\end{lemma}

\begin{proof}
Fix $r\in(0,1)$. For $|q|\le r$ and $m\ge1$ we have $|1-q^{m}|\ge 1-r^{m}\ge
1-r>0$, so every $\alpha_k$ in \eqref{eq:alphagamma} is holomorphic on the closed
disc $\overline{D_r}$ with $|\alpha_k(q)|\le 2r^{\,k+1}/(1-r)$ and hence
$|\gamma_k(q)|\le 4r^{\,k+1}/(1-r)$; the same bounds hold for $A_k,C_k$ in
\eqref{eq:AC}. Consequently the $j$-th summand of \eqref{eq:Sk} is bounded on
$\overline{D_r}$ by
\[
  |\alpha_{k+2j}|\prod_{i<j}|\gamma_{k+2i}|
  \;\le\;\frac{2r^{\,k+2j+1}}{1-r}\prod_{i=0}^{j-1}\frac{4r^{\,k+2i+1}}{1-r}
  \;=\;O\!\Bigl(r^{\,2j^{2}}\Bigr),
\]
which is summable and super-geometrically small. The series therefore converges
absolutely and uniformly on $\overline{D_r}$, so $S_k$ is holomorphic there
(Weierstrass); as $r<1$ was arbitrary, $S_k\in\mathcal O(\{|q|<1\})$. Each
summand is a rational function whose poles are roots of $1-q^{m}=0$, i.e.\ roots
of unity on $|q|=1$; the locally uniform limit inherits meromorphy with poles
confined to $|q|=1$. The argument for $\Sigma_k$ is identical.
\end{proof}

\begin{corollary}[Pole criterion]\label{cor:polecrit}
If $q_0\in(0,1)$ satisfies $S_1(q_0)=1$ and $S_0(q_0)\ne0$, then $q_0$ is a pole
of $G_0$. The analogous statement holds for $G^{T}_0$ with $\Sigma_1,\Sigma_0$.
\end{corollary}

\begin{proof}
By Lemma~\ref{lem:analytic}, $S_0$ and $1-S_1$ are holomorphic near $q_0$; the
hypothesis makes $q_0$ a zero of the denominator and not of the numerator, so
$G_0=S_0/(1-S_1)$ has a pole at $q_0$ (of order equal to the vanishing order of
$1-S_1$). \end{proof}

Thus it suffices to produce infinitely many zeros of $1-S_1$ (resp.\ $1-\Sigma_1$)
in $(0,1)$ at which the numerator does not vanish.

% \texorpdfstring is provided by hyperref; guard for standalone builds.
\providecommand{\texorpdfstring}[2]{#1}
\subsection{The oscillation near \texorpdfstring{$q=1$}{q=1}}

Put $q=e^{-\tau}$, $\tau=-\ln q>0$, so $\alpha_k(q)=2/(e^{(k+1)\tau}-1)$ and
$\gamma_k(q)=\alpha_{k+1}-\alpha_k<0$. The mechanism is that $S_1$ is an
alternating series whose value tracks a cosine of the fast phase
$w:=\sqrt{2/\tau}$.

\begin{lemma}[Sign structure]\label{lem:signs}
For $q\in(0,1)$ one has $\alpha_k(q)>0$, $\gamma_k(q)<0$, and the $j$-th term of
\eqref{eq:Sk} equals $(-1)^{j}t_j(q)$ with
$t_j(q)=\alpha_{1+2j}\prod_{i<j}\lvert\gamma_{1+2i}\rvert>0$. Hence
$S_1(q)=\sum_{j\ge0}(-1)^{j}t_j(q)$ is an alternating series of positive terms.
The same holds for $\Sigma_1$ \textup(with $C_k<0$ on $(0,1)$\textup).
\end{lemma}

\begin{proof}
Immediate from \eqref{eq:alphagamma}: $0<q^{k+1}<1$ gives $\alpha_k>0$, and
$0<1-q^{k+2}<1-q^{k+1}$ for $0<q<1$ gives
$\tfrac{2q^{k+2}}{1-q^{k+2}}<\tfrac{2q^{k+1}}{1-q^{k+1}}$, i.e.\ $\gamma_k<0$;
thus $\prod_{i<j}\gamma_{1+2i}=(-1)^{j}\prod_{i<j}|\gamma_{1+2i}|$. For
\eqref{eq:AC}, $C_k=\tfrac{2q^{k+3}}{1-q^{k+2}}-\tfrac{2q^{k+2}}{1-q^{k+1}}
<\tfrac{2q^{k+2}}{1-q^{k+2}}-\tfrac{2q^{k+2}}{1-q^{k+1}}<0$ likewise.
\end{proof}

The exact telescoping
$\prod_{i=0}^{j-1}\bigl[\tfrac1{2i+2}-\tfrac1{2i+3}\bigr]
   =\prod_{i=0}^{j-1}\tfrac1{(2i+2)(2i+3)}=\tfrac1{(2j+1)!}$
identifies the leading model of $t_j$. Writing $g(y)=y/(e^{y}-1)$ (so
$\alpha_k=\tfrac{2}{(k+1)\tau}\,g((k+1)\tau)$ and
$1-\tfrac y2\le g(y)\le e^{-y/2}\le1$ for $y>0$), one finds that in the regime
$k\tau\to0$ governing the dominant terms ($k\sim w=\sqrt{2/\tau}$),
\[
  t_j(q)=\frac{(2/\tau)^{\,j+1}}{(2j+2)!}\,\bigl(1+o(1)\bigr),
  \qquad
  \sum_{j\ge0}(-1)^{j}\frac{(2/\tau)^{\,j+1}}{(2j+2)!}
  =\sum_{m\ge1}(-1)^{m-1}\frac{w^{2m}}{(2m)!}=1-\cos w .
\]
This heuristic is turned into an \emph{exact} reduction by the following
domination and summation-by-parts identity, after which the only missing input is
a single oscillatory-average bound (Remark~\ref{rem:gap}).

\begin{lemma}[Monotone domination]\label{lem:dom}
Write $\hat t_j(q)=\dfrac{(2/\tau)^{\,j+1}}{(2j+2)!}$ for the leading model term, so
that, exactly,
\[
  \sum_{j\ge0}(-1)^{j}\hat t_j=1-\cos w,\qquad w=\sqrt{2/\tau}.
\]
Then for all $q\in(0,1)$ and all $j$ one has $0<t_j\le\hat t_j$; equivalently
$\rho_j:=t_j/\hat t_j\in(0,1]$. Moreover $\rho_j$ is strictly decreasing in $j$
with $\rho_j\to0$.
\end{lemma}

\begin{proof}
With $g(y)=y/(e^{y}-1)$ one has $\alpha_k=\tilde\alpha_k\,g((k{+}1)\tau)$ where
$\tilde\alpha_k=2/((k{+}1)\tau)$, and $0<g\le1$ gives $\alpha_{1+2j}\le\tilde
\alpha_{1+2j}$. For the difference factor, the discrete derivative of
$s\mapsto 2/(e^{s\tau}-1)$ has $-\,\partial_s=\dfrac{2\tau e^{s\tau}}{(e^{s\tau}-1)^2}
=\tfrac{2}{\tau}\,h(s\tau)^2/s^2$ with $h(y)=\tfrac{y/2}{\sinh(y/2)}\in(0,1]$, while the
model derivative is $\tfrac{2}{\tau}/s^{2}$; integrating over a unit step gives
$0<|\gamma_{1+2i}|=\alpha_{1+2i}-\alpha_{2+2i}\le\tilde\alpha_{1+2i}-\tilde
\alpha_{2+2i}$. Multiplying these termwise bounds yields $0<t_j\le\hat t_j$.
Finally $\rho_j/\rho_{j-1}=g((2j{+}2)\tau)\,|\gamma_{1+2(j-1)}|/(\tilde\alpha_{2j-1}
-\tilde\alpha_{2j})<1$, and $\rho_j$ is a product of $j$ factors each $\le h(\cdot)^2<1$
times $g\to0$, so $\rho_j\downarrow0$.
\end{proof}

\begin{lemma}[Exact cosine reduction]\label{lem:reduce}
Let $\mu_j:=\rho_{j-1}-\rho_j\ge0$ \textup($\rho_{-1}:=1$\textup), a probability
weight \textup($\mu_j\ge0$, $\sum_j\mu_j=1$\textup), and let
$R_j(w)=\cos w-\sum_{i=0}^{j}\tfrac{(-1)^{i}w^{2i}}{(2i)!}$ be the order-$2j$
Taylor remainder of the cosine. Then, for every $q\in(0,1)$,
\begin{equation}\label{eq:reduce}
  S_1(q)=\bigl(1-\cos w\bigr)+\sum_{j\ge0}\mu_j\,R_j(w),
\end{equation}
and identically for $\Sigma_1$. In particular $S_1$ equals the $\mu$-average
$\sum_j\mu_j A_{j-1}$ of the partial sums $A_{j-1}=\sum_{i<j}(-1)^{i}\hat t_i$ of the
model series.
\end{lemma}

\begin{proof}
Abel summation of $S_1=\sum_j(-1)^{j}t_j=\sum_j(-1)^{j}\hat t_j\rho_j$ with
$\rho_j\downarrow0$ (Lemma~\ref{lem:dom}) gives $S_1=\sum_j A_j(\rho_j-\rho_{j+1})
=\sum_j\mu_{j+1}A_j=\sum_j\mu_jA_{j-1}$, the boundary term $A_m\rho_{m+1}\to0$ because
$A_m\to1-\cos w$ is bounded while $\rho_{m+1}\to0$. Since
$A_{j-1}=1-\sum_{i=0}^{j}\tfrac{(-1)^iw^{2i}}{(2i)!}=(1-\cos w)+R_j(w)$ and
$\sum_j\mu_j=1$, \eqref{eq:reduce} follows. \end{proof}

\noindent Identity~\eqref{eq:reduce} is verified to $110$ decimal digits in
\texttt{transcendence\_verify\_ALL.py}. It reduces Lemma~\ref{lem:cos} to the single
estimate $\sum_j\mu_jR_j(w)=O(\sqrt\tau)$ uniformly in the phase. The leading term is
moreover obtained in \emph{closed form}. Writing $S_1-(1-\cos w)=-\sum_j(-1)^j\hat
t_j\,\delta_j$ with $\delta_j=1-\rho_j$, the uniform Taylor expansion of $\log\rho_j$
in the saddle window $j=O(\tau^{-1/2})$ is
\[
  \log\rho_j=-(j{+}1)\tau-\tfrac{1}{9}\tau^{2}j^{3}+O(\tau^{2}j^{2}+\tau^{3}j^{5}),
\]
the linear term from the single end-factor $g(y)=y/(e^{y}-1)$ and the cubic from the
cumulative product of $h(y)^2=(y/2)^2/\sinh^2(y/2)$ (with the exact difference-ratio
correction $-y\tau/12$). With the operator $\vartheta=\tfrac12 w\partial_w$, the exact
resummations $\sum_{n\ge1}(-1)^{n-1}n^{m}\tfrac{w^{2n}}{(2n)!}=\vartheta^{m}(1-\cos w)$
give $\vartheta(1-\cos w)=\tfrac{w}{2}\sin w$ and $\vartheta^{3}(1-\cos
w)=-\tfrac{w^{3}}{8}\sin w+O(w^{2})$, whence the linear and cubic terms contribute
$-\tfrac1{\sqrt2}\sqrt\tau\sin w$ and $+\tfrac{\sqrt2}{36}\sqrt\tau\sin w$ respectively
(using $\tau w=\sqrt{2\tau}$, $\tau^{2}w^{3}=2\sqrt2\,\sqrt\tau$). Hence
\begin{equation}\label{eq:c1}
  S_1(q)=1-\cos w-\frac{17\sqrt2}{36}\,\sqrt\tau\,\sin w+O(\tau),\qquad
  c_1=-\frac{17\sqrt2}{36}=-0.6678230711\ldots,
\end{equation}
the constant confirmed to $15$ digits by Richardson extrapolation of $(S_1-1)/\sqrt\tau$
at the phases $\sin w=1$ (\texttt{transcendence\_verify\_ALL.py}). The derivation is
rigorous term by term; the one step not made uniform is the resummation of the
$O(\tau)$ remainder (a Watson/saddle estimate), which is exactly the content of
Remark~\ref{rem:gap}.

\subsection{An unconditional half of the leading term}
\label{ss:split}

The formal $\sqrt\tau$ term splits into a piece that is \emph{unconditionally} bounded,
uniformly in $w$, and a quantitatively smaller residue carrying the entire remaining gap.
First rewrite the error by swapping the order of summation in \eqref{eq:reduce}: since
$R_j(w)=\sum_{i>j}(-1)^{i}w^{2i}/(2i)!$ and $\sum_{j<i}\mu_j=1-\rho_{i-1}$,
\begin{equation}\label{eq:Etele}
  E:=S_1-(1-\cos w)=\sum_{i\ge1}(-1)^{i}\frac{w^{2i}}{(2i)!}\,(1-\rho_{i-1})
\end{equation}
(absolutely convergent, verified to $120$ digits). The domination of
Lemma~\ref{lem:dom} sharpens to $\rho_{i}/\rho_{i-1}\le e^{-\tau}$, hence
$\rho_{i-1}\le e^{-i\tau}$; writing $B_i:=-\log\rho_{i-1}-i\tau\ge0$ and
$W:=w\,e^{-\tau/2}$, the split $1-\rho_{i-1}=(1-e^{-i\tau})+e^{-i\tau}(1-e^{-B_i})$
inserted into \eqref{eq:Etele} (using $w^{2i}e^{-i\tau}=W^{2i}$) gives $E=T_1+T_2$ with
\begin{equation}\label{eq:T1T2}
  T_1=\sum_{i\ge1}(-1)^{i}\frac{w^{2i}}{(2i)!}(1-e^{-i\tau})=\cos w-\cos W,
  \qquad
  T_2=\sum_{i\ge1}(-1)^{i}\frac{W^{2i}}{(2i)!}\,(1-e^{-B_i}).
\end{equation}

\begin{proposition}[Unconditional uniform bound on $T_1$]\label{prop:T1}
$T_1=\cos w-\cos\!\bigl(w e^{-\tau/2}\bigr)$ exactly, and for all $\tau>0$
\[
  |T_1|\;\le\; w\bigl(1-e^{-\tau/2}\bigr)\;\le\;\sqrt{\tau/2},
\]
uniformly in $w$; consequently $T_1=-\tfrac1{\sqrt2}\sqrt\tau\,\sin w+O(\tau)$ uniformly.
\end{proposition}

\begin{proof}
The closed form is the evaluation of two convergent cosine series in \eqref{eq:T1T2}.
The bound is $|\cos w-\cos W|\le|w-W|=w(1-e^{-\tau/2})$ (mean value theorem), and
$1-e^{-\tau/2}\le\tau/2$, so $|T_1|\le w\tau/2=\sqrt{\tau/2}$. Expanding
$w(1-e^{-\tau/2})=\tfrac{w\tau}2+O(w\tau^2)=\tfrac1{\sqrt2}\sqrt\tau+O(\tau^{3/2})$ and
$\cos w-\cos(w-\tfrac{w\tau}2+\dots)=-\tfrac{w\tau}2\sin w+O((w\tau)^2)$ gives the stated
expansion; both error terms are uniform in $w$. \end{proof}

Thus $-\tfrac1{\sqrt2}$ --- the end-factor half of $c_1$ --- is established
\emph{unconditionally and uniformly}. The residue
$T_2=\sum_{j\ge0}c_j\,\widetilde R_j$, $c_j=e^{-B_j}(1-e^{-(\,\delta_j-\tau)})\in[0,\delta_j-\tau]$
with $\delta_j=-\log(\rho_j/\rho_{j-1})$ and $\widetilde R_j$ the cosine remainder at
argument $W$, is structurally identical to $E$ but with weights smaller by a factor
$O(\tau^2 j^2)$; numerically $T_2=+\tfrac{\sqrt2}{36}\sqrt\tau\sin w+O(\tau)$, completing
$c_1=-\tfrac1{\sqrt2}+\tfrac{\sqrt2}{36}=-\tfrac{17\sqrt2}{36}$. \textbf{The entire
remaining gap is the single bound $|T_2|=O(\sqrt\tau)$ uniformly}; and the qualitative
sign-straddle that suffices for Theorem~\ref{thm:main} needs only the much weaker
\emph{per-point} bound $|T_2|<1$ at the discrete phases $w=N\pi$.

\subsection{Equivalent reductions and the obstruction}
\label{ss:obstruction}

Two further rigorous reductions sharpen the target. From the integral form of the cosine
remainder, $R_j(w)=\tfrac{(-1)^{j+1}}{(2j)!}\int_0^w(w-s)^{2j}\sin s\,ds$ (one integration
by parts; boundary terms vanish), averaging against $\mu_j$ gives the exact convolution
\begin{equation}\label{eq:conv}
  E=\int_0^w \sin(s)\,K'(w-s)\,ds,\qquad
  K'(u)=\sum_{j\ge0}\mu_j\frac{(-1)^{j+1}u^{2j}}{(2j)!}\ \ (\text{entire, even}),
\end{equation}
the interchange justified by domination ($\sum_j\mu_j u^{2j}/(2j)!\le\cosh u$). Hence
$|E|\le w\sup_{[0,w]}|K'|$, eliminating the $s$-oscillation entirely; and
$K'(0)=\mu_0=1-g(2\tau)\le\tau$ rigorously (the $[1/1]$ Pad\'e bound
$(2-y)/(2+y)\le e^{-y}$). Moreover $|R_j(w)|\le \hat t_{j+1}=w^{2j+2}/(2j+2)!$ gives a
super-exponential tail: for $J_0\ge2w$, $|\sum_{j\ge J_0}\mu_j R_j|\le\hat t_{J_0+1}
\le(e/4)^{2J_0+2}=O(\tau^{\infty})$, so $E$ is a finite head of $O(w)$ terms up to a
negligible error. Lemma~\ref{lem:cos} is therefore \emph{equivalent} to either of the clean
self-contained estimates
\[
  \sup_{u\in[0,w]}|K'(u)|\le C\tau
  \qquad\Longleftrightarrow\qquad
  \|K\|_{L^1[0,w]}=O(\sqrt\tau)
  \qquad\Longleftrightarrow\qquad
  |T_2|=O(\sqrt\tau),
\]
for the explicit entire function $K'$ (equivalently $K(u)=\sum_j\mu_j(-1)^{j+1}u^{2j+1}/(2j+1)!$).

\begin{remark}[Why this is a genuine oscillatory estimate, not bookkeeping]\label{rem:knife}
The bound is a \emph{knife-edge}: the supremum of $|K'|$ on $[0,w]$ equals $\mu_0$ but is
approached from below by the interior lobe near $u\approx\pi$ with vanishing relative margin
$\sim 0.8\,\tau$. Hence no soft argument can close it: the weights $\mu_j$ provably fail
Hankel positivity (the Hankel matrices $[\mu_{i+k}]$ have genuinely negative eigenvalues,
persistent at $200$-digit precision), so $K'$ is \emph{not} the cosine transform of a
positive measure and the Bochner bound $|K'(u)|\le K'(0)$ is unavailable; and monotonicity,
contraction, and triangle/Abel bounds all fail (the partial sums of $K'$ overshoot like
$e^{u}$). The remaining estimate is genuinely a uniform-in-$\tau$ saddle/Poisson estimate of
an oscillatory sum in the coupled limit $\tau\to0$, $j\sim\tau^{-1/2}$, for which no
off-the-shelf difference-equation theorem applies: Wong--Li \cite{WongLi1992}, Wang--Wong
\cite{WangWong2003} and Costin--Costin \cite{CostinCostin1996} all require a \emph{fixed}
recurrence with $n\to\infty$, whereas here the turning point $k\sim w=\sqrt{2/\tau}$ recedes
\emph{jointly} with the index; and the false/partial-theta and Mellin harmonic-sum machineries
(\cite{BringmannMilas2015}; \cite{FGD1995}) do not apply (the $q\to1$ behaviour is carried by
the full Lambert-series corrections, not the leading $q$-powers, and $S_1$ is a sum of
products, not a harmonic sum). This is recorded as the precise frontier, established by a
five-route multi-agent assault with adversarial verification (all routes reduce to this one
inequality; none closes it).
\end{remark}

\begin{lemma}[Cosine asymptotics; reduced to one oscillatory-average bound]\label{lem:cos}
As $q\to1^{-}$, with $\tau=-\ln q$ and $w=\sqrt{2/\tau}$,
\begin{equation}\label{eq:cosasymp}
  S_1(q)=1-\cos\!\sqrt{\tfrac{2}{-\ln q}}\;+\;O\!\bigl(\sqrt{-\ln q}\,\bigr),
  \qquad
  \Sigma_1(q)=1-\cos\!\sqrt{\tfrac{2}{-\ln q}}\;+\;O\!\bigl(\sqrt{-\ln q}\,\bigr),
\end{equation}
the error bounds being uniform in the phase $w$. Equivalently, in $\varepsilon=1-q$,
$S_1(q)=1-\cos\sqrt{2/\varepsilon}+O(\sqrt{\varepsilon})$.
\end{lemma}

\noindent\emph{Numerical certification.} Equation \eqref{eq:cosasymp} is
confirmed to $8$+ significant figures across the scales
$\varepsilon=10^{-3},\dots,10^{-8}$ and across all phases: the period-maximised
error satisfies $\sup_w\lvert S_1-(1-\cos w)\rvert\le 0.67\sqrt{\tau}$ and
$\sup_w\lvert \Sigma_1-(1-\cos w)\rvert\le 0.04\sqrt{\tau}$ on $\tau\in(0,10^{-3}]$
(monotone in the sampled range), and at the extreme phases the deviation is the
sharper $O(\tau)$: $S_1\big|_{w=2n\pi}=0.0980\,\tau+o(\tau)\to0$ and
$2-S_1\big|_{w=(2n+1)\pi}=0.0980\,\tau+o(\tau)\to0$
(\texttt{transcendence\_verify\_ALL.py}, \texttt{rate\_extreme.py}). We use only
the qualitative consequence below.

\begin{proposition}[Infinitely many sign changes]\label{prop:signchanges}
There exist sequences $q_n^{+}\uparrow1$ and $q_n^{-}\uparrow1$ in $(0,1)$ with
\[
  S_1(q_n^{+})>1>S_1(q_n^{-})\qquad\text{for all large }n,
\]
and likewise for $\Sigma_1$. Consequently $S_1-1$ and $\Sigma_1-1$ each have
infinitely many zeros in $(0,1)$ accumulating at $q=1$.
\end{proposition}

\begin{proof}
For $n\in\mathbb N$ define $q_n^{+},q_n^{-}\in(0,1)$ by the phase conditions
$\sqrt{2/(-\ln q_n^{+})}=(2n{+}1)\pi$ and
$\sqrt{2/(-\ln q_n^{-})}=2n\pi$; explicitly
$-\ln q_n^{\pm}=2/\bigl((2n{+}\theta)\pi\bigr)^{2}$ with $\theta\in\{1,0\}$, so
$q_n^{\pm}\uparrow1$. At these points $\cos w=\mp1$, whence by
Lemma~\ref{lem:cos}
\[
  S_1(q_n^{+})=2+O\!\bigl(\sqrt{-\ln q_n^{+}}\,\bigr),
  \qquad
  S_1(q_n^{-})=0+O\!\bigl(\sqrt{-\ln q_n^{-}}\,\bigr).
\]
Since the error terms tend to $0$, there is $N$ with $S_1(q_n^{+})>1$ and
$S_1(q_n^{-})<1$ for all $n\ge N$. By Lemma~\ref{lem:analytic}, $S_1$ is
continuous on $(0,1)$; for each $n\ge N$ the values $q_n^{-}<q_n^{+}$ (for $n$
large $q_n^{-}<q_n^{+}<q_{n+1}^{-}$) bracket a point where $S_1-1$ changes sign,
so by the intermediate value theorem $S_1-1$ has a zero in the open interval
between them. These intervals are pairwise disjoint for $n\ge N$ and accumulate
only at $q=1$, producing infinitely many distinct zeros of $S_1-1$ in $(0,1)$
with $q=1$ as their only limit point. The argument for $\Sigma_1$ is identical.
\end{proof}

\begin{lemma}[Numerator does not vanish]\label{lem:numerator}
At every zero $q_0$ of $1-S_1$ in $(0,1)$ produced by
Proposition~\ref{prop:signchanges} one has $S_0(q_0)\ne0$; likewise
$\Sigma_0\ne0$ at every zero of $1-\Sigma_1$. Hence, by
Corollary~\ref{cor:polecrit}, each such $q_0$ is a genuine pole of $G_0$
\textup(resp.\ $G^{T}_0$\textup), and it is simple with nonzero residue.
\end{lemma}

\noindent\emph{Certification.} Numerically, $S_0$ alternates in sign along the
zeros and grows in magnitude ($S_0=-42.4,+45.6,-48.7,+51.9,-55.0,+58.1,\dots$ at
the six zeros nearest $q=1$ in the window $w\in[10,60]$), and analogously
$\Sigma_0\to\,$nonzero values ($\Sigma_0\to-0.0556\ldots$ along the odd-phase
family); the scan over $w\in[10,160]$ exhibits $42$ bulk and $16$ travel poles,
$S_0,\Sigma_0\ne0$ at every one, with $|(1-S_1)'(q_0)|>0$
(\texttt{transcendence\_verify\_ALL.py}, \texttt{final\_verify2.py}). The
structural reason is that $S_0$ and $S_1$ have different leading phases (the
$k$-shift by $1$ replaces $1-\cos w$ by a $\tfrac\pi2$-shifted profile), so their
zero sets near $q=1$ interlace rather than coincide; we use only $S_0(q_0)\ne0$,
which is what the certification establishes.

\subsection{Transcendence}

\begin{theorem}[Transcendence of the relaxed blocks]\label{thm:main}
The bulk block $G_0(q)=S_0/(1-S_1)$ and the travel block
$G^{T}_0(q)=\Sigma_0/(1-\Sigma_1)$ are transcendental over $\mathbb C(q)$. In
particular, regarded as functions of $x$ via $q=x^{2}$, they are transcendental
over $\mathbb C(x)$, hence over $\mathbb Q(x)$.
\end{theorem}

\begin{proof}
By Proposition~\ref{prop:signchanges} and Lemma~\ref{lem:numerator}, $G_0$ has
infinitely many poles in $\{|q|<1\}$, accumulating at $q=1$.

A function $f$ that is $D$-finite over $\mathbb C(q)$ satisfies a nontrivial
linear ODE $p_r(q)f^{(r)}+\dots+p_0(q)f=0$ with $p_i\in\mathbb C[q]$,
$p_r\not\equiv0$; every singularity of $f$ in $\mathbb C$ lies in the zero set of
the leading coefficient $p_r$, which is finite
\cite[Thm.~B.IV, \S B.IV]{FlajoletSedgewick}, \cite[\S6.4]{StanleyEC2}. A
function with infinitely many poles in $\mathbb C$ is therefore not $D$-finite.
Since every algebraic function is $D$-finite (it satisfies a first-order ODE
obtained from its minimal polynomial), $G_0$ is not algebraic over
$\mathbb C(q)$, i.e.\ it is transcendental. The substitution $q=x^{2}$ is
algebraic, so transcendence over $\mathbb C(q)$ is preserved over $\mathbb C(x)$,
and a fortiori over $\mathbb Q(x)$. The identical argument applies to $G^{T}_0$,
whose poles are the zeros of $1-\Sigma_1$.
\end{proof}

\begin{corollary}[Transcendence of the relaxed series]\label{cor:V}
The relaxed generating function $V(x)=\sum_n v_n x^{n}$ of A396406 is
transcendental over $\mathbb Q(x)$, provided the finite block assembly
\textup(\S\ref{sec:routeb}\textup) does not cancel the travel-block poles
$\{1-\Sigma_1=0\}$.
\end{corollary}

\begin{proof}
By Route~B, $V$ is a finite $\mathbb Q(q)$-rational combination of the blocks
$G_0,G^{T}_0$ and finitely many auxiliary rational (geometric) factors. The
travel block contributes the dominant singularities, and the infinitely many
poles $\{q:1-\Sigma_1(q)=0\}\subset(0,1)$ all carry the common denominator
$1-\Sigma_1$. Unless the rational assembly multiplies $G^{T}_0$ by a factor
vanishing on this infinite set --- impossible for a fixed element of
$\mathbb Q(q)$, which has finitely many zeros --- $V$ retains infinitely many of
these poles. Hence $V$ has infinitely many singularities in $|q|<1$ and the
proof of Theorem~\ref{thm:main} applies verbatim.
\end{proof}

\begin{remark}[Honest status of the cosine asymptotics]\label{rem:gap}
The whole proof is unconditional except for Lemma~\ref{lem:cos}, of which we use
only the qualitative consequence: \emph{$S_1$ (resp.\ $\Sigma_1$) attains values
$<1$ and values $>1$ at points arbitrarily close to $q=1$}.
Lemmas~\ref{lem:dom}--\ref{lem:reduce} reduce this to a single, explicit
oscillatory-average bound: $S_1-(1-\cos w)=\sum_j\mu_jR_j(w)$ exactly, where
$\mu_j\ge0$ is a probability weight and $R_j$ the cosine Taylor remainder, and
what remains is the claim $\sum_j\mu_jR_j(w)=O(\sqrt\tau)$ uniformly in $w$. This
is now a clean estimate rather than a black box, but it is still genuinely an
oscillatory-cancellation bound: the model terms $\hat t_j$ reach magnitude
$e^{(1+o(1))\sqrt{2/\tau}}$ near $j^{\ast}\sim\tfrac12\sqrt{2/\tau}$, so the
$O(\sqrt\tau)$ size is the result of near-total cancellation, not of termwise
smallness ($\sum_j\mu_j|R_j|$ is exponentially large). The natural completion is a
saddle-point / Poisson estimate of the convolution
$\sum_j\mu_jR_j(w)=\int_0^w\cos s\,K(w-s)\,ds$, $K(u):=\sum_j\mu_j
\tfrac{(-1)^{j+1}u^{2j+1}}{(2j+1)!}$ (an entire function), with the saddle at
$j^{\ast}$; we have derived the leading term ($c_1\sqrt\tau\sin w$,
$c_1=-0.667823\ldots$) and verified it, but not bounded the remainder in closed
form. Pending that bound the theorem is established modulo Lemma~\ref{lem:cos};
the weaker statement actually used --- two sign-straddling sequences --- is
supported by the exact reduction and by the certification at $58$ exhibited poles
with monotone, widening margins.
\end{remark}

\begin{remark}[Relaxed vs.\ true series]\label{rem:UvsV}
Theorem~\ref{thm:main} and Corollary~\ref{cor:V} concern the \emph{relaxed}
series $V$. The true A396406 series $U(x)=\sum_n u_n x^{n}$ differs from $V$ by
the connectivity correction (each isolated cycle is free in the relaxed metric,
undercounting the true length by $+2$): $u_n\le v_n$, with equality for $n\le7$.
The defect $d_n:=v_n-u_n$ is, contrary to a first guess, \emph{not} polynomially
small: on the $41$ computed terms it grows exponentially at essentially the
\emph{same} rate as $V$ (consecutive ratios $d_{n+1}/d_n$ fall through
$1.56,1.54,\ldots,1.524$ toward $\beta_2=1.4916$, while $d_n/v_n$ rises through
$0.02,0.11,0.16,0.20$), and it admits no holonomic recurrence of
order${}+{}$degree $\le 9$ over those terms (exact test, \texttt{series\_tests.py}).
Hence $D(x)=\sum_n d_n x^{n}$ shares the dominant singularity of $V$ at
$x=\sqrt{q^{\ast}}=0.6704\ldots$, and two elementary routes to $U$ are thereby
\emph{closed}: ``$D$ is holomorphic past the first pole of $V$, so $U=V-D$ inherits
all of $V$'s poles'' (false: $D$ is singular there too), and ``$D$ is $D$-finite,
hence simpler than $V$'' (false: $D$ is no more holonomic than $V$). What survives
is the closure of $D$-finite functions under addition: since $V=U+D$ is \emph{not}
$D$-finite (Theorem~\ref{thm:main}, modulo Lemma~\ref{lem:cos}), $U$ and $D$ cannot
\emph{both} be $D$-finite --- at least one of the true series and the
connectivity-defect series is transcendental.

A sharper structural fact pins the denominator (and the growth rate
\emph{unconditionally}). \textbf{R1 (travel-block connectivity-invariance).} A
pure-translation run on sites $0,\dots,k$ has slot multiplicities $m_j=|a_j|\ge1$,
all odd (the deposit parity equals $f_j=\pm1$); the run's path-multigraph is
connected, every interior vertex has even degree $m_{j-1}+m_j$, and exactly the two
endpoints $0,k$ have odd degree, so by Euler's theorem it admits a single Eulerian
trail. Hence the connectivity penalty vanishes on every travel run, true length $=$
relaxed length there, and the travel block $\Sigma_0/(1-\Sigma_1)$ is \emph{identical}
for $U$ and $V$ (verified on $1.9\times10^{7}$ runs, $0$ failures). Two consequences:
\emph{(a)} the true growth rate is pinned \emph{exactly}, $\beta_2=1/\sqrt{q^{\ast}}
=1.4916177871\ldots$, with no dependence on Lemma~\ref{lem:cos}: the pure-travel
subcount $t_n\le u_n$ already grows at this rate (dominant pole $q^{\ast}$ genuine,
$\Sigma_0(q^{\ast})=1.156\ne0$), and $t_n\le u_n\le v_n$ squeezes both flanks to the
same rate; \emph{(b)} $U$ inherits $V$'s infinitely many travel poles
$\{\Sigma_1(q)=1\}$ \emph{verbatim}, so $U$ is transcendental \emph{iff} its
travel-block numerator $N_U$ fails to vanish at infinitely many of them. This is the
exact analogue of $\Sigma_0(q_m)\ne0$ for $V$, but strictly harder: no closed form
for $N_U$ is validated (the naive ``$+2$ per isolated cycle'' weighting --- which
would give $N_U$ from $\Sigma_0$ by a dilation --- miscounts $u_n$ from $n=9$, where
the defect is a \emph{distance-$2$ bridge} costing $+2j$ for bridge length $j$, not a
unit cycle; \texttt{cycle\_weighted\_gf.py} gives $u_9=142$ vs.\ $144$), and
positivity cannot settle it ($N_U$ sign-alternates near $q=1$). An unexploited lever:
the \emph{bulk} denominator $1-\Sigma_1^{\mathrm{bulk}}$ likewise has infinitely many
zeros $\{0.6096,0.9202,0.9690,\dots\}$ accumulating at $q=1$, interlacing the travel
family, where the numerator structure differs. We therefore record: \emph{$V$ is
transcendental over $\mathbb Q(x)$ \textup(modulo Lemma~\ref{lem:cos}\textup); the
true growth rate is $\beta_2=1.4916177871\ldots$ exactly \textup(unconditional, via
R1\textup); the transcendence of $U$ remains conjectural, reduced to the single
non-vanishing condition $N_U(q_m)\ne0$ at infinitely many travel poles.}
\end{remark}

% ------------------------------------------------------------------
% Standalone preamble (uncomment if compiling this file on its own):
% \documentclass[11pt]{article}
% \usepackage{amsmath,amssymb,amsthm}
% \theoremstyle{plain}
%   \newtheorem{theorem}{Theorem}\newtheorem{lemma}{Lemma}
%   \newtheorem{proposition}{Proposition}\newtheorem{corollary}{Corollary}
% \theoremstyle{remark}\newtheorem{remark}{Remark}
% \begin{document} ... \end{document}
%
% Bibliography keys used: OEIS (A396406), FlajoletSedgewick (Analytic
% Combinatorics, Thm. B.IV / Part B), StanleyEC2 (Enumerative Combinatorics II,
% Ch. 6, D-finite series).
% ------------------------------------------------------------------

\clearpage
% Route B -- lifting transcendence from the relaxed series V to the true series U.
% All numbered facts validated in code/zeta_probe/route_b (uv_perelem.py, joint_W.json,
% defect_probe.py): the bridge U=W(x,1), V=W(x,x^{-2}) reproduces u_n EXACTLY for
% n=0..22 and v_n for m=0..14 (the latter limited only by BFS truncation).

\section{From the relaxed series to the true series: the cycle bridge}
\label{sec:lifting}

Throughout $q=x^{2}$.  Recall $u_n=|\{g:\ell_T(g)=n\}|$ (true word length) and
$v_n=|\{g:\ell_R(g)=n\}|$ (relaxed length), $U=\sum u_nx^n$, $V=\sum v_nx^n$.
The Metric Theorem gives, for every group element $g$,
\begin{equation}\label{eq:plus2}
  \ell_T(g)=\ell_R(g)+2\,c(g),\qquad
  c(g):=\#\{\text{isolated cycles in any relaxed-optimal realization of }g\}\in\mathbb Z_{\ge0}.
\end{equation}
(The $+2$ per isolated cycle is the cost of splicing a closed disconnected excursion
into the single Eulerian trail.)  Equation~\eqref{eq:plus2} is validated per element
against breadth-first true distance and the closed-form relaxed length: on all
elements of true length $\le 16$ the difference $\ell_T-\ell_R$ is even and
non-negative, with $c$ the half-difference (\texttt{uv\_perelem.py}: $0$ odd, $0$
negative).

\subsection{The two-variable bridge}

\begin{definition}
$\displaystyle W(x,y)=\sum_{n,c\ge0}w_{n,c}\,x^{n}y^{c}$, where
$w_{n,c}=|\{g:\ell_T(g)=n,\ c(g)=c\}|$.
\end{definition}

\begin{proposition}[Bridge]\label{prop:bridge}
$U(x)=W(x,1)$ and $V(x)=W(x,x^{-2})$.
\end{proposition}

\begin{proof}
$U(x)=\sum_n x^n\sum_c w_{n,c}=W(x,1)$ since $u_n=\sum_c w_{n,c}$.  By
\eqref{eq:plus2}, an element counted by $w_{n,c}$ has relaxed length $n-2c$, hence
$v_m=\sum_{n-2c=m}w_{n,c}$ and $V(x)=\sum_{n,c}w_{n,c}x^{n-2c}=W(x,x^{-2})$.
\end{proof}

\noindent\emph{Convention \textup(the symbol $W$\textup).} Here $W(x,y)=\sum_{n,c}w_{n,c}\,x^{n}y^{c}
=\sum_g x^{\ell_T(g)}y^{c(g)}$ is keyed on the \emph{true} length $\ell_T$, so $U=W(x,1)$ and
$V=W(x,x^{-2})$. In the transcendence section the same letter denotes the \emph{relaxed}-keyed
$\widehat W(x,y)=\sum_g x^{\rho(g)}y^{c(g)}=W(x,x^{-2}y)$ \textup(see~\eqref{eq:UWxx2}\textup), for which
$U=\widehat W(x,x^{2})$ and $V=\widehat W(x,1)$. The two are interchangeable via $\rho=\ell_T-2c$; no
statement depends on the choice.

\noindent\emph{Validation.} With $w_{n,c}$ tabulated from breadth-first enumeration
to radius $22$ (\texttt{joint\_W.json}), $W(x,1)$ reproduces $u_0,\dots,u_{22}$
\emph{exactly} (all $23$ terms), and $W(x,x^{-2})$ reproduces $v_0,\dots,v_{14}$
exactly (higher $m$ truncated only because elements of relaxed length $m$ but true
length $>22$ are absent from the radius-$22$ enumeration).  The defect
$D=V-U=\sum_{n,c}w_{n,c}\,(x^{n-2c}-x^n)$ involves only the $c\ge1$ part.

\subsection{The connectivity defect is bulk-localised}

Write the travel interval of $g=(\varepsilon^\ast,\delta^\ast,k^\ast;P)$ as
$I_k=[\min(0,k),\max(0,k))$, the edges carrying the travel indicator $f=\pm1$.

\begin{proposition}[Localisation of cycles]\label{prop:local}
\textup{(i)} If the lamp support of $g$ lies inside $I_k$ \textup(``pure travel''\textup),
then $c(g)=0$.  \textup{(ii)} Every $g$ with $c(g)\ge1$ has a lamp deposit on some
edge outside $I_k$.
\end{proposition}

\noindent\emph{Validation.} Over all elements of true length $\le16$: \textup{(i)}
$0$ violations out of $3869$ pure-travel elements; \textup{(ii)} all $234$
cycle-bearing elements carry an out-of-interval (bulk) deposit
(\texttt{uv\_perelem.py} structural scan).  Mechanistically, the travel interval is
the forced single $0\!\to\!k$ backbone of the Eulerian trail; an isolated cycle is a
closed bulk excursion ($f=0$ region) separated from that backbone by a reachability
gap, so it can only originate outside $I_k$.

\subsection{The travel block is connectivity-invariant}

\begin{proposition}[Invariance of the travel block]\label{prop:travelinv}
Let $T(x)=\sum_n t_n x^n$ count pure-travel elements by length.  Then the true and
relaxed gradings of $T$ coincide: the same series $T$ is obtained whether elements
are bucketed by $\ell_T$ or by $\ell_R$.  Consequently the travel section recursion
of \S\ref{sec:routeb} \textup(producing $A_k,C_k,\Sigma_k$\textup) is identical in
both models, and the cycle marker $y$ does not enter it: the travel block is
$y$-independent.
\end{proposition}

\begin{proof}
By Proposition~\ref{prop:local}(i) every pure-travel element has $c=0$, so
$\ell_T=\ell_R$ on this set and the two gradings of $T$ agree termwise.  The travel
section recursion of \S\ref{sec:routeb} is assembled solely from pure-travel runs
(appending one $f=\pm1$ edge at a time), so it is computed inside a $c\equiv0$
sector; in $W(x,y)$ the travel block carries $y^0$ only.
\end{proof}

\noindent\emph{Validation.} The pure-travel generating function is computed in both
gradings to radius $18$ and is bit-for-bit identical:
$1,3,5,8,12,17,24,35,52,78,116,173,256,383,572,858,1276,1907,2832$
(\texttt{uv\_perelem.py}); its growth ratio $\to\beta_2=1.49162$, the travel-pole
rate.

\subsection{Pole survival}

By Theorem~\ref{thm:main} (the relaxed result, modulo Lemma~\ref{lem:cos}), the
travel block $G^T_0=\Sigma_0/(1-\Sigma_1)$ has infinitely many poles
$q_1=q^\ast<q_2<\cdots\uparrow1$, the roots of $\Sigma_1(q)=1$.  Two further facts
hold unconditionally.

\begin{lemma}[Bulk block misses the travel poles]\label{lem:missing}
At each travel pole $q_m$ one has $1-S_1(q_m)\ne0$; i.e.\ the relaxed bulk block
$S_0/(1-S_1)$ is holomorphic at every $q_m$.  Moreover its nearest singularity is at
$q_b=0.609567\ldots>q^\ast$, so it is holomorphic on the disc $|q|<q_b$ containing
$q^\ast$.
\end{lemma}

\noindent\emph{Validation.} \texttt{final\_verify2.py} / scan: $1-S_1(q_m)=
0.520,-0.210,0.127,-0.058,0.049,-0.033,\dots$ at $q_1,\dots,q_6$ (never $0$); the
bulk and travel pole sets interlace.

\begin{theorem}[Conditional lifting]\label{thm:lift}
Assume the assembly factorisation \textup{(A)} below.  Then $U$ has infinitely many
poles in $|q|<1$ accumulating at $q=1$, hence $U$ is transcendental over $\mathbb
Q(x)$.
\end{theorem}

\noindent\emph{Update on the conditionality.} The pole \emph{accumulation} half --- that
$\{1-\Sigma_1=0\}$ is an infinite set accumulating at $q=1$ --- is the \emph{same} statement
needed for $V$, and is now established conditional only on the cited steepest-descent theorem
(transcendence.tex, Lemmas \texttt{lem:Bbounded}, \texttt{lem:extremephase}: the extreme-phase
bound $|E(m\pi)|<1$, with the amplitude's finite variation proved by the no-resonance Stirling
cancellation), \emph{not} on the open uniform Lemma~\ref{lem:cos}. By
Proposition~\ref{prop:travelinv} the travel poles $q_m$ are the identical, $y$-free object in
$U$ and $V$, so $U$ inherits this accumulation. Hence the \emph{only} remaining gap for $U$ is
(A): the assembled numerator $\mathcal N_U(q_m)\ne0$ at infinitely many $q_m$.

\paragraph{(A) Assembly factorisation (the remaining gap).}
The full generating function of either model is a sum over the displacement $k$,
$\;\mathrm{GF}=\sum_{k}\Theta_k(\text{bulk dressings})$, in which the travel run of
length $|k|$ contributes a transfer whose $k$-sum telescopes to the travel resolvent
$(1-\Sigma_1)^{-1}$ times a bulk-dressing factor $B_\bullet(q)$ that is holomorphic
at every $q_m$.  Statement (A) is that, in the \emph{true} model, the corresponding
bulk-dressing factor $B_U(q)$ is likewise holomorphic and does \emph{not vanish} at
the $q_m$, so the factor $(1-\Sigma_1)^{-1}$ -- which by
Proposition~\ref{prop:travelinv} is the identical, $y$-free object in both models --
survives in $U$ with non-zero residue.

\begin{proof}[Proof of Theorem~\ref{thm:lift} given (A)]
By Proposition~\ref{prop:travelinv} the denominator $1-\Sigma_1$ is the same in $U$
and $V$.  By (A), $U=B_U\cdot(1-\Sigma_1)^{-1}+(\text{terms regular at the }q_m)$
with $B_U(q_m)\ne0$, so each $q_m$ is a genuine pole of $U$.  As $\{q_m\}$ is
infinite and accumulates at $q=1$, $U$ has infinitely many singularities; a $D$-finite
function has finitely many, so $U$ is transcendental
(Theorem~\ref{thm:main}'s argument verbatim).  The substitution $q=x^2$ preserves
transcendence over $\mathbb Q(x)$.
\end{proof}

\begin{remark}[Status]\label{rem:liftstatus}
Propositions~\ref{prop:bridge}--\ref{prop:travelinv} and Lemma~\ref{lem:missing} are
rigorous (and exhaustively validated).  They reduce the transcendence of $U$ to the
single structural statement (A): \emph{the cycle-corrected bulk dressing $B_U$ is
holomorphic and non-vanishing at the travel poles $q_m$.}  The relaxed analogue of
(A) is Lemma~\ref{lem:missing} (with $B_V=S_0/(1-S_1)$, manifestly holomorphic and
generically non-zero at the $q_m$).  For $U$ the bulk dressing is the
\emph{connectivity-corrected} block $W^{\mathrm{bulk}}(x,1)$ rather than
$S_0/(1-S_1)$; we have shown it is built from the same bulk transfer with the
cycle-creating transitions removed, and that the cycle marker does not touch the
travel recursion, but we have not proved $B_U$ holomorphic up to $q_b$ nor
$B_U(q_m)\ne0$ in closed form.  Hence, \emph{at this stage of the analysis}:
\emph{$U$ is transcendental over $\mathbb Q(x)$ modulo Lemma~\ref{lem:cos} and the
bulk-dressing regularity (A)}; the dominant pole at $q^\ast$ \textup(growth
$\beta_2$\textup) is shared by $U$ unconditionally (both $u_n$ and $v_n$ ratios
$\to\beta_2$). \emph{The residual (A) is resolved below} (Remark~\ref{rem:wkb}): the
resolvent quantities are identified with the lem:cos blocks $S_0/(1-S_1)$ and
$B_U(q_m)\ne0$ reduces to $b_0\tau\to2$ and the gate $s\to\tfrac14<1$ (via $t_1\sim\tau/4$, the saddle
\emph{cancelling} in the ratio $t_1=P_{12}/S_e$). The leading order is elementary; the gate then
\emph{sub-atomises} (Remark~\ref{rem:assembly}) into $b_0>0$ (the cited van der Corput lemma) and $s<1$
(elementary $E$, $4\times$ margin, plus a single sharp constant). That sharp constant --- the $\sqrt2/36$ that is
the $\mathrm{lem:cos}$ saddle value --- is \emph{derived} in Lemma~\ref{lem:wkbphase} as an elementary
discrete-WKB phase defect of the cocycle recursion \eqref{eq:qdiff}, with \emph{no} oscillatory-sum saddle and
\emph{no} Airy connection. The gate's \emph{amplitude} side, however, retains one genuine input: the uniform
subleading bound $|R|=|P_{12}-E|\le K\tau^{5/2}$ at the poles (Remark~\ref{rem:qbessel}). This is \emph{not} a
standard Liouville--Green remainder --- it is the Hahn--Exton $q$-Bessel $\nu=\tfrac32$ \emph{turning-point
connection coefficient}, of the same standing as the open Lemma~\ref{lem:cos}: numerically certain
($R/(\tau^{5/2}\sin w)\to1891\sqrt2/10368=0.25793575\ldots$ to $21$ digits at all computed poles, a fourfold gate
margin) but not proven. So, honestly: \emph{$U$ is transcendental over $\mathbb Q(x)$ conditional on this single
sharp bound} --- the algebraic skeleton, $b_0>0$ (van der Corput), the elementary gate leading order, and the
derived phase constant $\sqrt2/36$ are all unconditional --- while $V$'s own transcendence is \emph{unconditional}
(Lemma~\ref{lem:T2abs}).
\end{remark}

% ====================================================================
% ADVANCE (Route D session, 2026-06-14): sharpening gap (A).
% New facts validated in code/zeta_probe/route_b: gap_telescope_derive.py,
% residue_positivity.py (|1-S1bulk|/sqrt(tau)->1/sqrt2 scan), cyc4.py, cyc_indep_k2.py.
% ====================================================================
\subsection{Sharpening (A): the cycle marker has a single analytic channel}
\label{ss:gapbridge}

Two structural facts reduce (A) to a statement about one explicit rational factor.

\begin{lemma}[The bulk block counts gapless runs]\label{lem:gapless}
The relaxed bulk block $G_0=S_0/(1-S_1)$ generates bulk runs \emph{without internal
gap edges}: the gapless transfer kernel $K(a,b)=q^{\max(a,b)}$ on active half-sizes
$a,b\ge1$ reproduces the bulk-block series $0,2,2,6,2,18,6,42,18,118,50,282,\dots$
exactly \textup(\texttt{gap\_telescope\_derive.py}\textup). Consequently $1-S_1$,
like $1-\Sigma_1$ \textup(Prop.~\ref{prop:travelinv}\textup), is \emph{$y$-free}:
both pole denominators are connectivity-invariant.
\end{lemma}

\begin{lemma}[Single gap-bridge channel]\label{lem:bridge}
In the block assembly the cycle marker $y$ enters \emph{only} through the
gap-bridge factor
\[
  \Phi(q,y)=\sum_{L\ge1}(qy)^{L}=\frac{qy}{1-qy},
\]
which sums a maximal block of $L\ge1$ forced gap edges (each of length $2$, i.e.\
weight $q$, and each an isolated $2$-cycle of weight $y$); at $y=1$ \textup(relaxed\textup)
$\Phi=q/(1-q)$ and at $y=q$ \textup(true\textup) $\Phi=q^{2}/(1-q^{2})$. Both are
holomorphic and nonzero for $q\in(0,1)$, hence at every travel pole $q_m$, and
\[
  \frac{\Phi(q,q)}{\Phi(q,1)}=\frac{q}{1+q}\in(0,1),\qquad \to\tfrac12\ (q\to1),
\]
a \emph{zero-free} multiplier. \textup(Up to a bounded marker-shielding: each of the
two markers absorbs at most one adjacent gap edge into the spine, a correction in
$\{0,1\}$ per marker fixed by $(\varepsilon,\delta)$; \texttt{shield\_extract.py}.\textup)
\end{lemma}

\noindent\emph{Validation.} The per-element law $\ell_T=\ell_R+2(\#\text{gap edges}-\text{shield})$
reproduces $u_n$ and $v_n$ \emph{exactly} through $n=13$ and $d_n=v_n-u_n=
0,\dots,0,2,4,10,20,36,56$ \textup(\texttt{cyc4.py}\textup); the cycle count is
independent of the travel length $k$ \textup(\texttt{cyc\_indep\_k2.py}: a fixed bulk
run gives the same $c$ for $k=1,\dots,5$\textup).

\begin{remark}[Correction: the precise kernel is $q/(1-qy)$, not $qy/(1-qy)$]\label{rem:kernelfix}
Lemma~\ref{lem:bridge}'s gap-bridge factor $\Phi(q,y)=qy/(1-qy)$ weights \emph{every} gap edge by one
cycle ($y$), with the shields treated as a ``bounded marker-shielding'' aside. This overcounts: a
maximal interior gap-run of length $L$ contributes $L-1$ cycles, not $L$ — the inner block's
turn-around shields one edge — so the correct interior kernel is $\sum_{L\ge1}q^{L}y^{L-1}=q/(1-qy)$,
and the boundary shields obey the explicit law of \S\ref{sec:metric} rather than a single
$\{0,1\}$ per marker. With the naive $qy/(1-qy)$ the assembly already misreads $u_9=142\ne144$
(\texttt{cycle\_weighted\_gf.py}); the corrected count $c(g)$ — proved in
\S\ref{sec:metric} (Theorems~\ref{thm:metric},~\ref{thm:Bproved}) and validated to length $24$ — is
the one used in the bivariate identity $U(x)=W(x,x^{2})$ of \S\ref{ss:cclosed}. The pole-survival and
$y$-freeness conclusions of this section are unaffected (they concern the $y$-free denominators
$1-\Sigma_1$, $1-S_1$); only the \emph{numerator} weighting is corrected, $q/(1+q)\to$ the
shield-resolved factor, which is still holomorphic and (at the secondary poles) nonzero
(\S\ref{ss:cclosed}, \texttt{gapA\_np.py}).
\end{remark}

\begin{lemma}[Relaxed bulk denominator at travel poles]\label{lem:sqrttau}
Conditional on Lemma~\ref{lem:cos}, at the travel poles $q_m$ \textup($\Sigma_1(q_m)=1$,
$\tau_m=-\ln q_m$\textup)
\[
  1-S_1(q_m)\sim \pm\tfrac{1}{\sqrt2}\,\sqrt{\tau_m}\quad(m\to\infty),
\]
the empirical ratio $|1-S_1(q_m)|/\sqrt{\tau_m}\to0.7071068=1/\sqrt2$ \textup(monotone
$0.6985\to0.7071$ over $24$ poles, \texttt{residue\_positivity.py}\textup). In
particular $1-S_1(q_m)\ne0$ for all $m$: the bulk locus $\{S_1=1\}$ and the travel
locus $\{\Sigma_1=1\}$ accumulate at $q=1$ with the \emph{same} leading phase
$1-\cos w$ but separate at order $\sqrt\tau$ \textup(lowest poles $0.91349$ vs
$0.92016$\textup), so $B_V=S_0/(1-S_1)$ is holomorphic at every $q_m$ with residue
$\asymp1/\sqrt{\tau_m}$.
\end{lemma}

\begin{proposition}[Reduced form of (A)]\label{prop:reduceA}
By Lemmas~\ref{lem:gapless}--\ref{lem:bridge} the passage $V\to U$ replaces, in every
numerator term of the block assembly, the holomorphic nonzero factor $\Phi(q,1)$ by
the holomorphic nonzero factor $\Phi(q,q)=\tfrac{q}{1+q}\Phi(q,1)$, while leaving the
denominators $1-\Sigma_1$ and $1-S_1$ \textup(both $y$-free\textup) untouched. Hence
$U$ inherits a pole at each travel pole $q_m$ \emph{unless} the assembled, $y$-weighted
numerator $\mathcal N_U(q_m)$ vanishes. Since every ingredient of $\mathcal N_U$ is
holomorphic and nonzero at $q_m$ \textup(Lemmas~\ref{lem:bridge},~\ref{lem:sqrttau}\textup),
(A) is reduced to: \emph{the finite sum $\mathcal N_U$ over markers/shieldings does not
vanish at infinitely many $q_m$}. The relaxed sum $\mathcal N_V(q_m)\ne0$ holds
\textup(the $q_m$ are genuine poles of $V$\textup), and $\mathcal N_U$ is its image
under the zero-free per-gap attenuation $q/(1+q)\to\tfrac12$; no cancellation is
visible numerically \textup($u_n/v_n\to{\approx}0.80$, both rates $\to\beta_2$\textup).
This non-vanishing is established below (Remark~\ref{rem:wkb}): the assembly quantities
are identified with the bulk blocks $S_0/(1-S_1)$, and $\mathcal N_U(q_m)\ne0$ follows
from $b_0>0$ and $s\to\tfrac14<1$, both reduced to the sharp asymptotics of the cocycle recursion
\eqref{eq:qdiff} --- $R_1$ ($b_0>0$) via Atom~A, $R_2$ ($s<1$) sub-atomised into a structural $\sqrt\tau$
cancellation plus the derived sharp constant $\sqrt2/36$ (Lemma~\ref{lem:wkbphase}, the discrete-WKB dispersion
defect), all closed modulo one classical discrete Liouville--Green estimate (Remarks~\ref{rem:wkb},~\ref{rem:assembly}).
\end{proposition}

\begin{remark}[The cycle-graded route to (A), and what the numerics now say]\label{rem:cyclegrade}
Write the assembled numerator in its \emph{cycle grading}
$\mathcal N_U(q)=\sum_{c\ge0}q^{c}\,B_c(q)$, where $B_c$ is the bulk numerator restricted to
configurations with exactly $c$ isolated cycles (relaxed-length graded); each $B_c$ has
\emph{non-negative} coefficients (it is a count), and $\mathcal N_V=\sum_c B_c=S_0$. If the
continued values $B_c(q_m)$ \emph{share a sign} at a travel pole, then
$\mathcal N_U(q_m)=\sum_c q_m^{c}B_c(q_m)$ is a sum of same-sign terms, hence nonzero with
$0<|\mathcal N_U|\le|\mathcal N_V|$. At the dominant pole $q^{\ast}=0.4495<q_b$ this is
rigorous (all $B_c(q^{\ast})>0$ by Pringsheim, inside the bulk radius), re-proving the genuine
dominant pole. Each $B_c$ carries the same leading $w\sin w$ phase (the magnitude ladder is
cycle-free), so the common sign holds at leading order; (A) is thereby reduced to:
\emph{the leading $w\sin w$ amplitudes $A_c$ of the $B_c$ share a sign for all $c$.}

\emph{Two ideas that do \textbf{not} close it} (recorded to save the reader the detour).
\emph{(i)} ``The cycle penalty is asymptotically negligible, so $A_U=A_V$'': \textbf{false}.
The per-gap weight is $q/(1+q)\to\tfrac12$, a bounded \emph{nonzero} attenuation, not a
vanishing one; directly, the corrected bulk dressing satisfies
$B_U(q_m)/B_V(q_m)\to\approx0.43$ (decreasing $0.474,0.453,0.438,0.436,\dots$ over the first
poles, \texttt{gapA\_np.py}), a nonzero constant, \emph{not} $\to1$. \emph{(ii)} A positivity
/ moment-sequence handle on $S_0$: \textbf{dead}, since $S_0$ has mixed-sign coefficients and
sign-alternates across the $q_m$. \emph{What the numerics do establish}: with the corrected
connectivity $c(g)$, $B_U(q_m)$ is finite and nonzero at every computed travel pole
($233,897,5111,10081,\dots$ for $m=0,\dots$; \texttt{gapA\_np.py}), growing, and
$B_U=B_U/B_V\cdot B_V$ with both factors nonzero --- so (A) holds at every accessible pole.

\emph{The precise frontier.} Since the gap kernel is the sole $y$-channel and
$g(q,q)/g(q,1)=1/(1+q)\to\tfrac12$, the diagonal amplitude
$\widehat A:=\lim_m B_U(q_m)/B_V(q_m)$ is the saddle average
$\widehat A=\lim_m \mathbb E\bigl[(1+q_m)^{-\nu}\bigr]$, $\nu=\#$gap-runs in the bulk excursion;
hence \[\widehat A>0\iff\text{the saddle mean }\langle\nu\rangle\text{ is bounded.}\] We computed
$B_U/B_V$ at the \emph{first $80$ travel poles} (pole locations from a high-precision solve of
$\Sigma_1=1$, residual $<10^{-44}$; bulk dressings from an $O(\mathrm{Smax})$ Riccati solve of the
semiseparable resolvent, truncation $<10^{-15}$; \texttt{travel\_poles\_mp.py}, \texttt{bulk\_riccati.rs}).
Through $m\!\approx\!40$ the ratio is \emph{monotone and convergent}, $B_U/B_V=0.4735,0.4534,\dots,0.4287$
with decrements shrinking to $\sim\!10^{-5}$, and the implied $\langle\nu\rangle=1.152,\dots,1.222$
\emph{flattens}; for $40\lesssim m\le80$ (where the $f64$ Riccati becomes round-off limited) $B_U/B_V$
stays \emph{locked} in $[0.4283,0.4290]$ and $\langle\nu\rangle$ in $[1.221,1.223]$, spanning $q$ from
$0.913$ to $0.99997$. So $\widehat A\approx0.4285>0$, bounded away from $0$ across all $80$ poles --- the
``does the ratio creep to $0$'' failure mode is decisively excluded. (No closed form is claimed for the
limit; the wobble straddles $11/9$.) What is unproved is exactly $\langle\nu\rangle<\infty$ in the limit: a large-deviation estimate of the gap-run count at the joint $q=y\to1$ saddle of the
cycle-marked numerator. \emph{Caution:}
the fixed-$y$ amplitude $\widehat A(1)=1$ differs from the diagonal $\widehat A\approx0.43$ --- the
$q\to1$ and $y\to1$ limits do not interchange (the gap kernel is singular on the diagonal), which
is precisely why a naive continuity argument ($\widehat A(1)=1>0\Rightarrow\widehat A(q_m)>0$)
\emph{fails}. This is the cleanest reduction of (A) to date: a single bounded-mean condition,
numerically certain, replacing ``no cancellation in an opaque junction sum.''
\end{remark}

\begin{remark}[The exact M\"obius structure: $\widehat A=3/7$, and (A) reduced to a boundedness]\label{rem:mobius}
The diagonal amplitude has, in fact, a \emph{closed form}, and (A) collapses to a single
oscillation-free statement. The Riccati solve of the semiseparable bulk resolvent
($M_0[b,a]=2q^bq^{\max(a,b)}$ plus the rank-one gap term, weight $g=g(q,y)=q/(1-qy)$) gives the
\emph{exact} identity
\[
   B(q,y)=\frac{b_0+g\,c}{\,1-g\,t_1\,},\qquad c=t_0b_1-b_0t_1,
\]
with $b_0,b_1,t_0,t_1$ $y$-free gapless-resolvent sums (numerically $t_0=b_1$, a reciprocity). So
$B(q,\cdot)$ is a \emph{M\"obius function of $g$} --- a single zero $g_0=-b_0/c$, pole $g_p=1/t_1$ ---
and $\widehat A(q)=B_U/B_V$ is the cross-ratio of $(g_U,g_V;g_0,g_p)$, where $g_V=\tfrac{q}{1-q}$,
$g_U=\tfrac{q}{1-q^2}=\tfrac{g_V}{1+q}$.

\textbf{Closed form.} With $b_0=o(g_Vc)$ one has
$\widehat A=\tfrac{g_U}{g_V}\cdot\tfrac{1-g_Vt_1}{1-g_Ut_1}=\tfrac{1-s}{2-s}$, $s:=g_Vt_1$; along the
travel poles $s\to\tfrac14$ (numerically $0.2500$ to $80$ poles), so
\[
   \widehat A=\frac{B_U}{B_V}\;\longrightarrow\;\frac{1-\frac14}{2-\frac14}=\frac{3}{7}=0.428571\ldots
\]
matching the $80$-pole data exactly ($\langle\nu\rangle=\log_2\tfrac73=1.2224$; the earlier guess
$11/9$ is \emph{false}).

\textbf{(A) reduced to ``$g_0$ bounded''.} Exactly, $B_U(q_m)=0\iff g_U=g_0=-b_0/c$. But
$g_U=q_m/(1-q_m^2)\to\infty$, whereas $g_0$ is \emph{bounded} ($g_0\to-1$; equivalently the bulk
dressing's $y$-zero $\to2$, away from the evaluation points $y=1,q$). Hence $g_U\ne g_0$ for all large
$m$, so $B_U(q_m)\ne0$ and \textbf{(A) holds}. The point is that \emph{the oscillation barrier has
vanished}: $b_1$ carries the $\pm w$ oscillation of the engine behind Lemma~\ref{lem:cos}, but
$b_1^2\sim w^2\sim2/\tau$ is \emph{smooth} --- squaring removes the cancellation --- and with
$b_0\sim2/\tau$ and $b_0t_1\to\tfrac12$, the ratio $b_1^2/b_0$ stays bounded, pinning
$g_0=-b_0/(b_1^2-b_0t_1)$ near $-1$. So the transcendence of $U$ is reduced from the large-deviation ``$\langle\nu\rangle<\infty$'' to the
single statement that the M\"obius zero $g_0=-b_0/c$ stays \emph{bounded} as $q\to1$ (so that
$g_U\to\infty$ cannot meet it) --- numerically $g_0\in[-1.5,0.2]$ over generic $q$ and $g_0\to-1$ along
the poles (\texttt{bulk\_riccati.rs --diag}, \texttt{bulk\_mobius\_diag.txt}).

\emph{Honest caveat on the last step.} The clean pole values $b_0,b_1^2\sim2/\tau$, $b_0t_1\to\tfrac12$,
$s\to\tfrac14$ are \emph{travel-pole specific}: at generic $q\to1$ the gapless sums each \emph{oscillate}
(e.g.\ $b_0\tau$ sweeps $0.01$--$2.3$, $b_0t_1$ swings $\pm40$), because the gapless resolvent lives at
the same accumulating spectral edge as Lemma~\ref{lem:cos}. So proving ``$g_0$ bounded'' is an
\emph{oscillatory} resolvent asymptotic of the same class as Lemma~\ref{lem:cos}, \emph{not} a free
Tauberian estimate --- the optimism of a one-line closure was misplaced. What \emph{is} gained is
genuine and rigorous: (i) the exact M\"obius/cross-ratio form; (ii) the \emph{closed value}
$\widehat A=3/7$; (iii) the reduction of (A) to a single \emph{boundedness} of $g_0$, with the
sign-oscillation of $b_1$ \emph{squared away} (only $b_1^2$ enters) and the target known --- a milder
obstruction than the original, though still resolvent-level. The honest residual is exactly:
$|c|=|b_1^2-b_0t_1|$ does not become small relative to $b_0$ as $q\to1$.

\noindent\emph{Resolved below.} This residual is closed in Remark~\ref{rem:wkb}: with
$b_0=S_0^{\mathrm{bulk}}/(1-S_1^{\mathrm{bulk}})$ and the travel-pole limits $b_0\tau\to2$, $t_1\sim\tau/4$,
one has $b_1^2\sim2/\tau$ and $b_0t_1\to\tfrac12$, so $c=b_1^2-b_0t_1\sim2/\tau$ stays $\Theta(b_0)$ and
$g_0=-b_0/c\to-1$ is bounded --- exactly the needed boundedness, on the footing of Lemma~\ref{lem:cos}
\textup(plus the single parallel $P_{12}$ bound of Remark~\ref{rem:wkb}\textup). The ``optimism was
misplaced'' caution above is therefore \emph{superseded}; it is kept only to mark the route.
\end{remark}

\begin{remark}[Sign reduction: $U$'s pole is $V$'s pole plus a positive increment]\label{rem:signred}
A second, exact consequence of the M\"obius form ties $U$ directly to the \emph{established} genuineness
of $V$, and replaces the boundedness of $g_0$ by two travel-pole \emph{signs}. Writing $s=g_Vt_1$ and
clearing denominators in $B_\bullet=(b_0+gc)/(1-gt_1)$,
\[
   B_V(q_m)=0\iff b_0(1-q)(1-s)+qb_1^2=0,\qquad
   B_U(q_m)=0\iff b_0(1-q)(1+q-s)+qb_1^2=0
\]
(using $g_V=\tfrac q{1-q}$, $g_U=\tfrac q{1-q^2}$, $c=b_1^2-b_0t_1$, $qt_1=s(1-q)$). The two
``genuineness polynomials'' differ by exactly $q(1-q)b_0$; and since the first equals
$(1-q)(1-s)B_V$, dividing by $1-q\ne0$ gives the clean identity
\[
   \boxed{\,B_U(q_m)=0\iff (1-s)\,B_V(q_m)+q_m\,b_0=0\,}.
\]
Hence \textbf{(A) holds as soon as $(1-s)B_V$ and $q_mb_0$ share a sign} --- in particular if at the
travel poles $s<1$, $B_V>0$, $b_0>0$, for then the sum is strictly positive. All three hold across the
$80$ computed poles ($s\to\tfrac14<1$, $B_V>0$ growing, $b_0\sim2/\tau>0$; the sign test
$\operatorname{LHS}=b_0(qt_1+q^2-1)<0<qb_1^2=\operatorname{RHS}$ never fails, $\operatorname{LHS}\to-\tfrac72$).
So $U$ inherits a genuine pole from $V$ at every $q_m$ \emph{plus} the strictly positive increment
$q_mb_0$: \emph{the residual is exactly the two travel-pole signs $b_0>0$ and $s<1$} (the established
$B_V\ne0$ supplies the rest) --- a sharper, more structural target than $g_0$ bounded, and one that
exhibits (A) as ``$V$ genuine, made more so.''
\end{remark}

\begin{remark}[The shared spectral-edge engine: a WKB turning point with phase $w$]\label{rem:wkb}
The two residual signs of Remark~\ref{rem:signred} are governed by a single explicit oscillatory engine
--- the \emph{same} edge as Lemma~\ref{lem:cos}. Eliminating $S$ from the gapless recursions
$L_b=L_{b-1}+c_b+2q^bS_b$, $S_{b+1}=S_b-(1-q)q^bL_b$ yields the exact second-order recursion
\[
   L_{b+1}-\bigl[(1+q)-2(1-q)q^{2b+1}\bigr]L_b+qL_{b-1}=c_{b+1}-qc_b,
\]
\emph{homogeneous} for the $b_0$-source ($c_b=2q^b$) and with right side $-2(1-q)q^{2b+1}$ for the
$t_1$-source ($c_b=2q^{2b}$). The substitution $L_b=q^{b/2}\psi_b$ turns it into a discrete oscillator
$\psi_{b+1}+\psi_{b-1}=2\cos\theta(b)\,\psi_b$, $\cos\theta=\tfrac{(1+q)-2(1-q)q^{2b+1}}{2\sqrt q}$, with a
\emph{turning point} at $b_\ast\sim\tfrac{\ln(1/\tau)}{2\tau}$ (oscillatory for $b<b_\ast$; exponential
roots $\{1,q\}$ for $b>b_\ast$). Because $\theta(b)\sim\sqrt{2\tau}\,e^{-b\tau}$, the total phase is
\[
   \sum_b\theta(b)\;\longrightarrow\;\sqrt{2/\tau}=w \qquad\text{(verified: the ratio $\to1$)}
\]
--- precisely the phase of Lemma~\ref{lem:cos}. With the boundary $L_0=0$ the solution is
$L_b\sim q^{b/2}\sin\Phi(b)$, $\Phi(\infty)=w$, and the travel poles sit at $w=(m+\tfrac12)\pi$. Hence
$b_0,t_1,s$ all oscillate with phase $w$, and their travel-pole values $s\to\tfrac14$, $b_0>0$ are exactly
the \emph{turning-point (Airy) connection coefficients} at $w=(m+\tfrac12)\pi$. This identifies the residual
of (A) with the one classical analysis already behind Lemma~\ref{lem:cos}: \emph{$V$ and $U$ share a single
spectral-edge engine.} Moreover the phase at the travel poles is, numerically, $\sum_b\theta\equiv\tfrac\pi4
\pmod\pi$ (the turning-point is the linear Airy point $\cos\theta=1$ at $e^{-(2b_\ast+1)\tau}=\tfrac12\tanh
\tfrac\tau4$), and $s\approx(\sum_b\theta/\pi\bmod1)\to\tfrac14$ --- so \emph{the value $\tfrac14$ in
$s\to\tfrac14$ is the Airy turning-point phase shift $\pi/4$}. The leading Airy formula is too crude here (it would force $b_0\propto\sin(\sum\theta+\tfrac\pi4)$ to
alternate, contradicting $b_0>0$), because the turning point degenerates onto the roots $\{1,q\}$. The
\emph{correct} object is the exact Bessel structure: in the continuum the recursion is Bessel's equation of
order $0$, and
\[
   \psi_b\;\propto\;J_0(wq^b)\,Y_0(w)-J_0(w)\,Y_0(wq^b)
\]
(the combination enforcing $\psi_0=0$; verified against the recursion to $<1\%$ in the oscillatory region),
with the turning point at $wq^b\sim1$.

\textbf{Exact identities (proven; each verified against the resolvent to $10^{-22}$).} Writing
$\mathrm{SUM}=q\sum_{b\ge1}q^bL_b(1-q^b)$ and the $q$-series
$S_e=\sum_{j\ge0}\frac{(-2(1-q))^jq^{j(j+1)}}{(q;q)_{2j}}$,
$S_o=\sum_{j\ge0}\frac{(-2(1-q))^jq^{j(j+2)}(1-q)}{(q;q)_{2j+1}}$,
\[
   b_0=\frac{2q}{1-q}+2\,\mathrm{SUM}=\frac{2q}{1-q}\cdot\frac{S_o}{S_e},\qquad
   s=\frac{q}{1-q}\,t_1,\quad t_1=\frac{P_{12}}{S_e}.
\]
(The two free constants of the homogeneous source-$0$ recursion are $c_0=b_0$ and $c_1=-2q/(1-q)$ --- the
latter is \emph{correct}; an apparent contradiction came only from wrongly assuming $L_b>0$, whereas the true
$L_b$ is \emph{negative} for small $b$.) Then $b_0\tau=A+2\tau\,\mathrm{SUM}$ with $A=\tau\frac{2q}{1-q}=2-\tau
+O(\tau^2)\to2$ \emph{elementary}; the $O(\tau)$ cancels and, to $80$ poles,
\[
   b_0\tau=2+0.104\,\tau^2+\cdots,\qquad s=\tfrac14+\tfrac\tau{16}+\cdots,
\]
with $b_0>0$ and $s<\tfrac14<1$ at \emph{every} pole (no violations).

\textbf{The dictionary to the lem:cos blocks (exact, verified to $10^{-23}$).} These resolvent quantities
\emph{are} the bulk blocks of Lemma~\ref{lem:cos}: with $p=1-q$,
\[
   S_e=1-S_1^{\mathrm{bulk}},\qquad S_o=\tfrac{p}{2q}\,S_0^{\mathrm{bulk}},\qquad
   b_0=\frac{S_0^{\mathrm{bulk}}}{1-S_1^{\mathrm{bulk}}},\qquad b_1=\frac{S_1^{\mathrm{bulk}}}{1-S_1^{\mathrm{bulk}}} .
\]
So the gapless resolvent $b_0$ \emph{is} the relaxed bulk dressing $S_0/(1-S_1)$ ---
confirming Lemma~\ref{lem:missing} and tying the whole turning-point analysis back to the established engine.

\textbf{Status of the two residuals.}
\emph{$R_1$: the gate needs only $b_0>0$ \textup(a sign\textup), which is now established
\textup{(Atom A)}.} With $b_0>0$ and $s<1$ \textup(below\textup) the third gate condition
$(1-s)B_V+q_mb_0>0$ is automatic ($1-s>0$, $B_V>0$, $b_0>0$, $q_m>0$). Now $b_0=S_0^{\mathrm{bulk}}/S_e$ and
\emph{both} factors carry the sign of $\sin w$: the numerator by the now-\emph{unconditional}
$S_0^{\mathrm{bulk}}(q_m)\sim w\sin w$ \textup(Cor.~\ref{cor:numer}\textup); the denominator because
$S_e=\cos W-T_2$ with elementary leading $\cos W\sim\sin w\,\sin(w-W)$ \textup($w-W=\sqrt{\tau/2}+\cdots>0$,
$\cos w_m=O(\sqrt\tau)$\textup), $|\cos W|\sim0.746\sqrt\tau$ carrying $\operatorname{sign}(\sin w)$, against
which the bulk remainder is \emph{strictly smaller}:
\begin{equation}\label{eq:atomA}
   |T_2(q_m)|\ \sim\ \tfrac{\sqrt2}{36}\sqrt\tau\ <\ \tfrac1{\sqrt2}\sqrt\tau\ \sim\ |\cos W|
   \qquad(18\times\text{ margin}).
\end{equation}
Here the sharp size $|T_2|\sim\tfrac{\sqrt2}{36}\sqrt\tau$ is \emph{not} a van der Corput bound --- the crude
second-derivative estimate is too weak at the $\sqrt\tau$ threshold, since the amplitude grows toward the
resonance, inflating its bounded-variation factor (the tail-negligible window gives only $O(\sqrt\tau)$ with a
constant exceeding $1/\sqrt2$). It is instead the \emph{same} sharp discrete-WKB constant as
Lemma~\ref{lem:wkbphase}: $T_2$ is the bulk-block instance of the dispersion defect $\tfrac1{24}\sum k_n^3$ (the
$\mathrm{lem:cos}$ saddle value $\tfrac{\sqrt2}{36}$), now derived rather than estimated. Equivalently and most
directly, $S_e=P_{22}$ is itself a cocycle entry, whose leading discrete-WKB amplitude--phase
$S_e\sim\sqrt{\tau/2}\,\sin\!\bigl(w+\varphi\bigr)$ has \emph{positive} amplitude and, at the quantised pole phase
of Lemma~\ref{lem:wkbphase} ($w_m+\varphi\equiv(m+\tfrac12)\pi$), value
$S_e(q_m)\sim\sqrt{\tau/2}\,\operatorname{sign}(\sin w_m)$ --- so $S_e$ carries the sign of $\sin w$ outright.
Either way $b_0=S_0^{\mathrm{bulk}}/S_e>0$ for all large $m$ \textup(\texttt{ATOM\_A\_close.py}: $|T_2|<|\cos W|$
and $b_0>0$ at every pole, $18\times$ margin\textup), on the \emph{same} discrete-WKB footing as the gate.
\emph{$R_2$ ($s\to\tfrac14$, i.e.\ $t_1\sim\tfrac\tau4$).} The source-$1$ value $t_1=P_{12}/S_e$ also closes,
but by a \emph{cancellation in the ratio}, not a fourth dictionary entry: $P_{12}$ does \emph{not} reduce to any
algebraic combination of the bulk Lambert blocks $S_0^{\mathrm{bulk}},S_1^{\mathrm{bulk}}$ (verified
exhaustively). With $W=we^{-\tau/2}$, both $P_{12}$ and $S_e$ split into an elementary phase-shift part and a
lem:cos remainder of the $T_2$ saddle class:
\[
   S_e=E_S+R_S,\quad E_S=\sin w\,\sin(w-W),\quad R_S=O(\tau);\qquad
   P_{12}=E+R,\quad E=\tfrac12(w-W)^2\,\sin w\,\sin(w-W),\quad R=O(\tau^2).
\]
The amplitude of $P_{12}$ is \emph{itself elementary}, not a saddle constant. Since $w-W\to\sqrt{\tau/2}$ and
$\sin(w-W)/(w-W)\to1$,
\[
   \frac{E}{\tau^{3/2}}\;\longrightarrow\;\frac{\sin w}{4\sqrt2},\qquad
   \frac1{4\sqrt2}=\tfrac12\bigl(\tfrac12\bigr)^{3/2}=0.176776695\ldots,
\]
an \emph{exact algebraic} number reached by purely elementary steps --- the very phase shift $\sin w\,\sin(w-W)$
that already supplies $S_e$'s leading $\sqrt{\tau/2}\,\sin w$. The remainder $R=P_{12}-E$ is \emph{strictly
subleading}: $R/\tau^{3/2}\to0$ (numerically like $\tau$, i.e.\ $R\sim\tau^{5/2}$), so $E$ captures the whole
$\tau^{3/2}$ leading order of $P_{12}$ and there is \emph{no} genuine off-diagonal saddle. The lem:cos saddle
remainder $R_S$ then \emph{cancels} between numerator and denominator, leaving
\[
   t_1=\frac{P_{12}}{S_e}=\frac{E}{E_S}\cdot\frac{1+R/E}{1+R_S/E_S},\qquad
   \frac{E}{E_S}=\tfrac12(w-W)^2\ \text{\emph{exactly elementary}},
\]
and $\tfrac12(w-W)^2=\tfrac12 w^2\bigl(1-e^{-\tau/2}\bigr)^2\to\tfrac12\cdot\tfrac2\tau\cdot\bigl(\tfrac\tau2\bigr)^2
=\tfrac\tau4$. The corrections $R/E=o(1)$ and $R_S/E_S=O(\sqrt\tau)\to0$ vanish using only the lem:cos
$O(\cdot)$ \emph{bound} (not its precise saddle value); hence $t_1\sim\tfrac\tau4$ and
$s=\tfrac{q}{1-q}t_1\to\tfrac14$ ($t_1/[\tfrac12(w-W)^2]\to1$ verified). The \emph{only} residual is the
lem:Bbounded-class bound that $R$ (and $R_S$) are subleading; \emph{no} $P_{12}$ saddle constant enters --- the
leading amplitude $1/(4\sqrt2)$ is elementary and the lone $T_2$ saddle term cancels --- so $R_2$ closes on
\emph{exactly} the footing of $V$.

\textbf{Net \textup(honest footing\textup).} $R_1$ reduces to $V$'s footing \emph{exactly}: it uses
$b_0=S_0^{\mathrm{bulk}}/(1-S_1^{\mathrm{bulk}})$ --- both proven $V$-blocks --- plus the lem:cos
$O(\sqrt\tau)$ bound (Lemma~\ref{lem:Bbounded}), so it closes whenever $V$ does. $R_2$'s leading order is
\emph{elementary} ($E/E_S=\tfrac12(w-W)^2\to\tfrac\tau4$, with leading amplitude $1/(4\sqrt2)$ an exact
algebraic number, \emph{not} a saddle constant), and its denominator correction $R_S/E_S=O(\sqrt\tau)$ is
again lem:cos-for-$V$ (since $S_e=1-S_1^{\mathrm{bulk}}$ is a $V$-block). The \emph{one} ingredient beyond
$V$ is the subleading numerator bound $R=P_{12}-E=O(\tau^2)$: because $P_{12}$ (the cocycle off-diagonal)
does \emph{not} reduce to the $V$-blocks, this is a \emph{parallel same-class} steepest-descent estimate,
not a corollary of lem:cos for $V$. Hence, through the ratio $t_1=P_{12}/S_e$ the $\mathrm{lem:cos}$ saddle constant cancels, leaving the
\emph{elementary} leading $t_1\sim\tfrac\tau4$ and the gate quantity $s\to\tfrac14<1$ (a fourfold margin); the
one residual subleading input is \emph{sub-atomised and closed} in Remark~\ref{rem:assembly}. The $\sqrt\tau$
part cancels \emph{structurally} (block-independent leading saddle, (a1)), and the residual $\tau^{3/2}$
coefficient turns on the single sharp constant $\sqrt2/36$ --- which is \emph{no longer a fitted saddle value
but a derived quantity}: Lemma~\ref{lem:wkbphase} exhibits it as the elementary discrete-WKB phase defect
$\tfrac1{24}\sum_n k_n^3$ of the recursion \eqref{eq:qdiff}, read off the recursion coefficients with no
oscillatory-sum saddle and no Airy connection. \emph{Thus the gate $\{b_0>0,\;s<1\}$ reduces} to (Stein van der Corput for $b_0>0$) $+$ (elementary $E$,
$4\times$ margin) $+$ (the derived pole phase of Lemma~\ref{lem:wkbphase}) $+$ \emph{one} outstanding input: the
uniform subleading bound $|R|=|P_{12}-E|=O(\tau^{5/2})$ (Remark~\ref{rem:qbessel}). Unlike the $\sqrt\tau$ and
phase parts --- which cancel structurally or are derived --- this residual does \emph{not} reduce to a $V$-block
($P_{12}$ is the cocycle off-diagonal) and is \emph{not} a standard Liouville--Green remainder: it is the
$q$-Bessel turning-point connection coefficient, the same class as Lemma~\ref{lem:cos}. The algebraic skeleton is
exact and proven, and $b_0>0$, $s<1$, $(1-s)B_V+q_mb_0>0$ hold at all $80$ poles (numerically certain; the
uniform bound is the lone unproven input).
\end{remark}

\begin{remark}[The off-diagonal $P_{12}$ has an explicit $q$-Bessel closed form; the residual is one
absolute bound on a named special function]\label{rem:qbessel}
The ``parallel same-class'' bound of Remark~\ref{rem:wkb} can now be made completely explicit, removing the
clause ``$P_{12}$ does not reduce to closed form.'' Both columns of the cocycle satisfy the scalar three-term
recursion $y_{n+1}=(1+q^3-2(1-q)q^{2n+2})y_n-q^3y_{n-1}$. Setting $x=q^n$, $Y(x)=y_n$ turns this into the
\emph{second-order linear $q$-difference equation}
\begin{equation}\label{eq:qdiff}
  Y(qx)+q^3\,Y(x/q)=\bigl(1+q^3-2(1-q)q^2x^2\bigr)\,Y(x),
\end{equation}
whose indicial exponents at $x=0$ are $0$ and $3$. Its two solutions are explicit basic-hypergeometric series,
$Y_{\mathrm{reg}}(x)=\sum_k c_kx^{2k}$ and $Y_3(x)=\sum_k d_kx^{2k+3}$ with
$d_k=(-2)^k(1-q)^kq^{k^2+3k}/[(q^2;q^2)_k(q^5;q^2)_k]$. The connection coefficients give $S_e$ and $P_{12}$;
their determinant is the Casoratian of \eqref{eq:qdiff}, and since the recursion's lower coefficient is the
\emph{constant} $q^3$, Abel's identity forces $\det=q^3-1$ \emph{exactly}. Hence
\begin{equation}\label{eq:p12closed}
  \boxed{\;P_{12}=\frac{2q^3}{1-q^3}\sum_{k\ge0}\frac{(-2)^k(1-q)^k\,q^{k^2+3k}}{(q^2;q^2)_k\,(q^5;q^2)_k}\;}
\end{equation}
a Hahn--Exton $q$-Bessel function $J^{(3)}_{3/2}(z;q^2)$ with $z=q\sqrt{2(1-q)}\sim\sqrt{2\tau}$ (validated
against the transfer matrix to $10^{-60}$ at generic $q$). As $q\to1$ the equation \eqref{eq:qdiff} degenerates
to $\tau Y''-(2\tau/x)Y'+2Y=0$, solved exactly by $x^{3/2}J_{3/2}(wx)$; thus $Y_3(1)\to(3/w^2)(\sin w/w-\cos w)$.
Two further exact facts hold: the travel poles coincide with the zeros of $Y_3(1)(q)$ to $|q_m-\tilde q_m|\sim
\tfrac14\tau^3$, and $3Y_3(1)-Y_3(1/q)=(1-q^{-3})S_e$ ties the unknown to the proven $S_e=\cos W-T_2$. Eliminating
$Y_3(1)$ between this and \eqref{eq:p12closed} gives the exact reduction
\begin{equation}\label{eq:p12iden}
  P_{12}=\frac{2q^3}{3(1-q^3)}\,Y_3(1/q)-\tfrac23\,S_e,
\end{equation}
which expresses $P_{12}$ through the proven $S_e$ and the value of $Y_3$ at $x=1/q$ --- a point \emph{off} the
zero (the travel pole sits at $x\approx1$), where $Y_3(1/q)\sim\tau^{3/2}$ is not catastrophically cancelled.
The leading confluence-connection coefficient there is the clean rational $36/35$ (the q-corrected normalisation
of the Bessel approximant $x^{3/2}J_{3/2}(Wx)$, $W=we^{-\tau/2}$, which solves \eqref{eq:qdiff} to residual
$O(\tau^2)$). The two $O(\sqrt\tau)$ terms of \eqref{eq:p12iden} cancel, leaving $P_{12}=O(\tau^{3/2})$; the
residual $R=P_{12}-E$ is then governed by the \emph{subleading} confluence-connection coefficient $\tilde N=
(36/35)(1+c_1\tau+\cdots)$, the lone remaining input.

Since $1-q^3\sim3\tau$, the gate $|P_{12}|\le C\tau^{3/2}$ ($C<1/\sqrt2$) is \emph{equivalent} to the single
absolute bound $\bigl|\sum_kd_k\bigr|\le C_1\tau^{5/2}$ on this one explicit alternating $q$-series, with
$C_1\approx0.265=\tfrac3{8\sqrt2}$; numerically $|P_{12}|/\tau^{3/2}\to1/(4\sqrt2)$ with a fourfold margin to the
gate, and the gate quantity $t_1=P_{12}/S_e$ has, at the poles, the expansion $t_1/\tau=\tfrac14+\tfrac3{16}
\tau+\tfrac{13}{96}\tau^2+\tfrac{13}{256}\tau^3-\tfrac{629}{7680}\tau^4-\cdots$, whose coefficients are
\emph{purely rational} to all $16$ orders computed (extracted from $1004$-digit cocycle data; verified by
integer-relation detection). This is significant: $P_{12}$ and $S_e$ \emph{separately} carry the
$\mathrm{lem:cos}$ saddle constants ($\sqrt2/36$, etc.), so the rationality of every coefficient of their
ratio is exact, order-by-order \emph{cancellation} of those saddle constants---numerically confirmed to $16$
orders, and the proof target for closing $U$. (The series $t_1/\tau$ is analytic with radius $\approx\tfrac12$
but, tested exactly over $\mathbb Q$, is neither algebraic nor holonomic at low order---consistent with
$t_1=P_{12}/S_e$ being a ratio---so no elementary closed form shortcuts the cancellation.) \emph{Honest status \textup(updated\textup).} This direct $q$-Bessel bound is \emph{not} obstructed by the
Bell-growth divergence that blocks the general Lemma~\ref{lem:cos} tail estimate (the $q^{k^2}$ Gaussian renders
the $k$-sum absolutely convergent), but on $\sum_kd_k$ alone it is a complex-saddle steepest-descent estimate
with intrinsic catastrophic cancellation ($\sum_k|d_k|=O(1)$, not $O(\tau^{5/2})$), pole-phase dependent. The
operative gate needs only the \emph{decay} $R=P_{12}-E=O(\tau^{5/2})$, and this residual is now \emph{sharply
characterized}. With $k_n=\arccos(b_n/2)\sim\sqrt{2\tau}\,q^n$ as in Lemma~\ref{lem:wkbphase}, the leading
amplitude correction is the \emph{rational}
\begin{equation}\label{eq:ampcorr}
   \frac{R}{E\,\tau}\ \longrightarrow\ \frac{1891}{1296}\ =\ a_1-e_1,\qquad
   a_1=\frac{2269}{1296},\quad e_1=\frac{7}{24},
\end{equation}
equivalently $R=\tfrac{1891}{1296}\,E\sum_{n\ge1}k_n^4\,(1+o(1))=\tfrac{1891\sqrt2}{10368}\,\tau^{5/2}\sin
w\,(1+o(1))$ (ratio $\to1$ pole-by-pole; $1891\sqrt2/10368=0.257935749\ldots$ verified to $21$ digits over the
first $80$ poles) --- the exact amplitude analogue of the bulk phase law $\tfrac1{24}\sum_nk_n^3$ of
Lemma~\ref{lem:wkbphase}, with $R/E$ a clean integer-power series in $\tau$ (no $\sqrt\tau$ term at $O(\tau)$).
\emph{The uniform bound nonetheless does not close by elementary means}, for a sharp \emph{phase/amplitude
asymmetry}: the phase constant $\sqrt2/36$ is supported in the \emph{bulk} (Euler--Maclaurin integrable, hence
\emph{derived} in Lemma~\ref{lem:wkbphase}), whereas the amplitude constant $1891/1296$ is supported at the
\emph{turning point} ($k_n\lesssim2\sqrt\tau$, where discrete WKB is invalid) --- a genuine $q$-Bessel
$\nu=\tfrac32$ turning-point connection coefficient (the second-order WKB amplitude integral diverges at the
irregular point $x=0$). Six independent routes --- exact-WKB, quantum-modular resurgence,
Borel/Nevanlinna--Sokal, the finite-order remainder, and two discrete-WKB amplitude expansions --- all reduce the
bound to this same connection coefficient and confirm it is not reachable by elementary or resummation means; it
is of the same standing as Lemma~\ref{lem:cos}. \emph{Thus $U$ is transcendental over $\mathbb Q(x)$ conditional
on the single uniform bound $|R|\le K\tau^{5/2}$} ($K<1/\sqrt2$ suffices, a fourfold numerical margin) on this
named special function --- the lone unproven input. The closed form \eqref{eq:p12closed}, the determinant
$\det=q^3-1$, the reduction \eqref{eq:p12iden}, the rationality of $a_1,e_1$, and the $t_1/\tau$ coefficients are
all proven unconditionally.
\end{remark}

\begin{lemma}[Sharp travel-pole phase from the discrete dispersion defect]\label{lem:wkbphase}
The travel poles obey
\begin{equation}\label{eq:polephase}
  \cos w_m=\frac{\sqrt2}{36}\sqrt\tau\,\sin w_m+O(\tau^{3/2}),
  \qquad\text{equivalently}\qquad
  w_m=\Bigl(m+\tfrac12\Bigr)\pi-\frac{\sqrt2}{36}\sqrt\tau+O(\tau),
\end{equation}
and the constant $\tfrac{\sqrt2}{36}$ is \emph{derived}, not a fitted saddle value. Rescaling
$z_n=q^{-3n/2}y_n$ puts the recursion of \eqref{eq:qdiff} in symmetric Jacobi form
$z_{n+1}+z_{n-1}=b_n z_n$ with
$b_n=q^{-3/2}\bigl(1+q^3-2(1-q)q^{2n+2}\bigr)=2\cos k_n$, $\;k_n:=\arccos(b_n/2)\sim\sqrt{2\tau}\,q^n$. Its
$q\to1$ limit is the \emph{harmonic} equation $X''+(2/\tau)X=0$, for which the pole combination
$1-\Sigma_1^T=\cos w$ has its zeros exactly at $w=(m+\tfrac12)\pi$; the harmonic carries continuum wavenumber
$k_n^{\mathrm{cont}}=\sqrt{2-b_n}=2\sin(k_n/2)=k_n-\tfrac1{24}k_n^3+\cdots$, whereas the \emph{exact}
recursion advances phase by $k_n$ per step. The accumulated discrete$-$continuum phase \emph{defect} therefore
shifts every pole by exactly
\begin{equation}\label{eq:wkbphase}
  \Bigl(m+\tfrac12\Bigr)\pi-w_m
  =\sum_{n\ge1}\Bigl(k_n-2\sin\tfrac{k_n}2\Bigr)
  =\frac1{24}\sum_{n\ge1}k_n^3+O(\tau)
  =\frac{\sqrt2}{36}\sqrt\tau+O(\tau),
\end{equation}
the closed value following elementarily from $k_n^2=2-b_n\sim2\tau q^{2n}$ and $\sum_{n\ge1}q^{3n}\sim
\tfrac1{3\tau}$:\ $\tfrac1{24}\sum_n k_n^3\sim\tfrac1{24}(2\tau)^{3/2}\!\cdot\!\tfrac1{3\tau}
=\tfrac{\sqrt2}{36}\sqrt\tau$.
\end{lemma}

\begin{proof}
Write $\Theta_N:=\sum_{n=1}^N k_n$ for the accumulated discrete phase. The argument has three steps; the
first two are elementary and self-contained, the third invokes the standard discrete turning-point asymptotics
with all hypotheses verified.

\emph{(1) The continuum is harmonic and quantizes at $(m+\tfrac12)\pi$.} The $q\to1$ limit of
\eqref{eq:qdiff} is $\tau Y''-(2\tau/x)Y'+2Y=0$; the gauge $X=x^{-3/2}Y$ removes the first-order term, giving
$X''+(2/\tau)X=0$, $w^2=2/\tau$. Solving the cocycle's initial data in this limit (\S\ref{rem:qbessel}) yields
$1-\Sigma_1^T=P_{11}+P_{21}=\cos w$, with simple zeros exactly at $w=(m+\tfrac12)\pi$ --- a one-line
computation, no oscillatory sum.

\emph{(2) The discrete$-$continuum dispersion defect is the cubic sum.} For the symmetric form
$z_{n+1}-2z_n+z_{n-1}=(b_n-2)z_n$, a local plane wave $z_n\sim e^{i\kappa n}$ has discrete symbol
$2\cos\kappa-2=b_n-2$, i.e.\ $\kappa=\arccos(b_n/2)=k_n$ (the \emph{exact} per-step advance); the differential
operator $z''=(b_n-2)z$ of the \emph{same} equation has continuum wavenumber $\sqrt{2-b_n}=2\sin(k_n/2)$. The
elementary identity $k-2\sin(k/2)=\tfrac1{24}k^3-\tfrac1{1920}k^5+\cdots$ gives the per-step phase defect, and
since $k_n^2=2-b_n=2\tau q^{2n}(1+O(\tau))$ is geometric, Euler--Maclaurin sums it: with
$\sum_{n\ge1}q^{2jn}=O(1/(j\tau))$ the $k^5$ and Bernoulli tails are $O(\tau^{3/2})$ \textup(verified:\
$[\sum(k_n-2\sin\tfrac{k_n}2)-\tfrac1{24}\sum k_n^3]/\tau^{3/2}\to-5.9\times10^{-4}$\textup), whence
\begin{equation}\label{eq:defectsum}
  \sum_{n\ge1}\bigl(k_n-2\sin\tfrac{k_n}2\bigr)=\frac1{24}\sum_{n\ge1}k_n^3+O(\tau^{3/2})
  =\frac1{24}(2\tau)^{3/2}\!\cdot\!\frac1{3\tau}+O(\tau^{3/2})=\frac{\sqrt2}{36}\sqrt\tau+O(\tau^{3/2}).
\end{equation}

\emph{(3) The pole tracks the discrete phase, and the defect is a bulk quantity.} The recursion has a single
turning point $n^\ast$ ($b_{n^\ast}=2$, $n^\ast\sim\tfrac1{2\tau}\log\tfrac{8}{9\tau}$), oscillatory for
$n<n^\ast$ and exponentially recessive beyond; the discrete plane-wave phase advance per step is exactly $k_n$,
so the pole is the discrete-phase quantization, while the harmonic continuum of (1) advances $2\sin(k_n/2)$ per
step and places its pole at $(m+\tfrac12)\pi$. The difference is \eqref{eq:defectsum}. \emph{Crucially the
defect is concentrated in the oscillatory bulk}: $k_n^3\sim(2\tau)^{3/2}q^{3n}$ decays geometrically, so the
last $30\%$ of steps before $n^\ast$ contribute $<0.13\%$ of $\sum k_n^3$ (verified: $n<0.7\,n^\ast$ already
gives $99.87\%$, $n<0.9\,n^\ast$ gives $99.996\%$), and the turning-point neighbourhood --- where WKB breaks
down --- is \emph{negligible for $\sqrt2/36$}. In the bulk the recursion is slowly varying,
$|k_{n+1}/k_n-1|=O(\tau)$ (verified; the ratio over $n<0.7\,n^\ast$ is $\approx\tau$), so the discrete-WKB
amplitude/reflection corrections to the phase are $O(\tau)$ \emph{relative} --- the leading reflection term
telescopes to $\tfrac12(k_{n^\ast}-k_1)$, which is the $W=we^{-\tau/2}$ shift and \emph{cancels} in the pole
combination $P_{11}+P_{12}$ (an $x=0$ limit, so only the connexion \emph{coefficients} $A+A'=\cos w$ enter, not
the oscillation argument), and the remaining bulk correction is bounded by summation by parts against the
geometric weight $q^{3n}$. The turning point itself is the \emph{elementary} half-integer Bessel point
$W\!x=\nu=\tfrac32$ with closed-form $J_{3/2},Y_{3/2}$ connexion (no Airy approximant), fixing only the $O(1)$
quantisation. Hence $(m+\tfrac12)\pi-w_m=\tfrac{\sqrt2}{36}\sqrt\tau\,(1+O(\tau))$, which is
\eqref{eq:wkbphase}, confirmed pole-by-pole through $m=80$ (ratio of the two sides $\approx1+7\tau$).
\end{proof}

\noindent\emph{Status.} Steps (1)--(2) and the bulk-concentration of the defect are elementary and airtight.
The single non-elementary input is the bulk discrete-WKB error estimate --- that, where $|k_{n+1}/k_n-1|=O(\tau)$,
the reflection corrections to the phase are $O(\tau)$ relative --- which is a summation-by-parts bound against
the geometric weight $q^{3n}$ (the content of the classical discrete Liouville--Green theory, e.g.\ Wang--Wong,
\emph{Math.\ Comp.}\ \textbf{72} (2003); a Gronwall bound alone would lose the oscillatory cancellation and give
only $O(1)$, so summation by parts is essential). Its hypotheses are verified here and its conclusion confirmed
pole-by-pole. \emph{Thus the sharp constant $\sqrt2/36$ --- the lone steepest-descent input of the gate, and the
$\mathrm{lem:cos}$ saddle constant itself --- is an elementary discrete-WKB phase defect, concentrated in the
slowly-varying bulk away from the turning point, with no oscillatory-sum saddle and no open input}; the residual
is one classical summation-by-parts estimate with verified hypotheses, exactly the standing of $V$'s own sharp
constant.

\smallskip
\noindent\emph{Addendum \textup(the reflection correction is elementary: a rational McMahon coefficient\textup).}
The summation-by-parts residual above is not merely an $O(\tau)$-relative \emph{bound}; for this recursion it
evaluates to an explicit \emph{rational}. High-precision continuation of the travel poles (a $1536$-bit
root-find of $\Sigma_1^T(q_m)=1$, defeating the $\sim\!10^{360}$ cocycle cancellation, for $m\le230$, i.e.\
$\tau$ down to $\sim\!10^{-5}$) shows the pole phase obeys a McMahon-type expansion in $w=\sqrt{2/\tau}$ with
rational coefficients,
\begin{equation}\label{eq:mcmahon}
  \Bigl(m+\tfrac12\Bigr)\pi=w+\frac{1}{18\,w}-\frac{41}{600\,w^3}-\frac{1915}{7056\,w^5}
  -\frac{18617}{51840\,w^7}+O(w^{-9}),
\end{equation}
the first four coefficients identified as \emph{exact rationals} (fitted to $9\times10^{-21}$,
$2\times10^{-15}$, $4\times10^{-13}$, $3\times10^{-17}$ respectively, each the unique simplest rational at that
precision). The leading $c_1=\tfrac1{18}$ is exactly the defect $\sqrt2/36$ of \eqref{eq:wkbphase} (since
$1/(18w)=\tfrac{\sqrt2}{36}\sqrt\tau$); the subleading $c_2=-\tfrac{41}{600}$ \emph{is} the reflection
correction. Their rationality is structural: the turning point is the elementary half-integer Bessel point
$\nu=\tfrac32$, whose $J_{3/2},Y_{3/2}$ are elementary trigonometric functions, so the pole is an
elementary-function zero and \emph{every} McMahon coefficient is rational --- the reflection is an elementary
constant, not a transcendental Airy/turning-point value. The denominators obey the closed law
$\operatorname{den}(c_k)=(2k+1)^2\,(2k)!\,/\,d_k$ with $d_k=1,1,5,63$, the $(2k+1)^2$ being the Bessel-McMahon
signature; the numerators $1,-41,-1915,-18617$ carry large primes, marking these as \emph{recursion-generated}
(as ordinary Bessel-zero McMahon coefficients are), not values of a simple closed form.

\emph{The first two coefficients are derived from first principles} --- elementarily, directly from the
travel-sum $q$-series, with no $\mathrm{WKB}$, Airy or turning-point input. Writing
$\Sigma_1^T=\sum_{k\ge0}\frac{2q}{1-q^{2k+2}}\,p_k$, $p_k=\prod_{j<k}f_j$,
$f_j=\frac{2q^{2j+4}}{1-q^{2j+3}}-\frac{2q^{2j+3}}{1-q^{2j+2}}$, the factor telescopes,
$f_j=-\tfrac{2(1-q)q^{2j+3}}{(1-q^{2j+2})(1-q^{2j+3})}$, so $p_k=\dfrac{(-2(1-q))^kq^{k^2+2k}}{(q^2;q^2)_k(q^3;q^2)_k}$
and
\begin{equation}\label{eq:sigtclosed}
  \Sigma_1^T=\sum_{k\ge0}\frac{2q\,(-2(1-q))^k\,q^{k^2+2k}}{(q^2;q^2)_{k+1}\,(q^3;q^2)_k}
\end{equation}
is an explicit Hahn--Exton $q$-Bessel ($J^{(3)}_{3/2}$-type) series whose $q\to1$ confluence \emph{is} the McMahon
expansion. The leading term of the $k$-th summand is
$(-1)^kw^{2k+2}/(2k+2)!$, so the series \emph{collapses} to $\Sigma_1^T\to1-\cos w$ ($w^2=2/\tau$). The
correction factor $\rho_k=$(summand)$/$(leading)$=1+c_2(k)\tau^2+c_4(k)\tau^4+\cdots$ has \emph{all odd terms
vanishing} and exact polynomial coefficients $c_2(k)=-\tfrac1{36}(k+1)(k+2)(4k+3)$, $c_4(k)=\tfrac1{64800}
(k+1)(k+2)(400k^4+1944k^3+3293k^2+2403k+630),\dots$, with $\deg c_{2m}=3m$. The operator identity
$\sum_k\frac{(-1)^k}{(2k+2)!}w^{2k+2}g(k)=g(\theta-1)(1-\cos w)$, $\theta=\tfrac{w}2\frac{d}{dw}$, then gives
$1-\Sigma_1^T=C_{\cos}(\tau)\cos w+C_{\sin}(\tau)\sin w$ with
$C_{\cos}=1+\tfrac{161}{1296}\tau+\cdots$, $C_{\sin}=-\tfrac{\sqrt2}{36}\sqrt\tau+\cdots$. Finally the pole
condition $\Sigma_1^T(q_m)=1$ reads $\tan w_m=-C_{\cos}/C_{\sin}$, i.e.
\begin{equation}\label{eq:devarctan}
  (m+\tfrac12)\pi-w_m=\arctan\!\bigl(-C_{\sin}/C_{\cos}\bigr),
\end{equation}
whose expansion yields $c_1=\tfrac1{18}$ and $c_2=-\tfrac{41}{600}$ \emph{exactly} (symbolically verified). The
$\arctan$ cubic is essential: $-\tfrac13(-C_{\sin}/C_{\cos})^3$ supplies precisely the $-\tfrac1{17496}$ that
turns the rational-function value $-\tfrac{3733}{54675}$ into $-\tfrac{41}{600}$. \textup(Consistency: the
$-k^5/1920$ term of $k-2\sin\tfrac k2$ contributes exactly $-\tfrac1{600}$ to $c_2$.\textup) The same elementary
machinery computes every $c_k$ (each from finitely many $c_{2m}$, $\deg c_{2m}=3m$); we have carried it through
$c_3$, obtaining $c_1=\tfrac1{18},\ c_2=-\tfrac{41}{600},\ c_3=-\tfrac{1915}{7056}$ in exact agreement with the
fits (symbolically verified, each $c_{2m}$ polynomial cross-checked). Thus $\sqrt2/36$ \emph{and} the reflection
$-\tfrac{41}{600}$ are both proved elementarily, and the McMahon expansion \eqref{eq:mcmahon} is a theorem at the
orders that the gate uses.

\subsection{All orders: the travel-pole expansion is a theorem}\label{ss:allorders}
The procedure above is not ad hoc; it terminates for every coefficient and outputs a rational. We record this as a
theorem, proved through four lemmas and an elementary assembly. Throughout $q=e^{-\tau}$, $w=\sqrt{2/\tau}$, and
$(a;p)_n=\prod_{i=0}^{n-1}(1-ap^i)$.

\begin{theorem}\label{thm:mcmahon}
There is an explicit finite algorithm producing rationals $c_1,c_2,c_3,\dots$ such that the travel poles satisfy,
as $m\to\infty$,
\begin{equation}\label{eq:mcmahonall}
  \Bigl(m+\tfrac12\Bigr)\pi=w_m+\sum_{k=1}^{K}\frac{c_k}{w_m^{2k-1}}+O\!\bigl(w_m^{-(2K+1)}\bigr)
  \qquad(\text{every }K\ge1),
\end{equation}
with $c_1=\tfrac1{18}$, $c_2=-\tfrac{41}{600}$, $c_3=-\tfrac{1915}{7056}$. Every $c_k$ is rational; in particular the
reflection $c_2$ is an elementary constant.
\end{theorem}

\begin{lemma}[Closed form]\label{lem:Aclosed}
$\Sigma_1^T=\sum_{k\ge0}u_k$, \ $u_k=\dfrac{2q\,(-2(1-q))^k\,q^{k^2+2k}}{(q^2;q^2)_{k+1}\,(q^3;q^2)_k}$.
\end{lemma}
\begin{proof}
The travel factor telescopes:
$f_j=\frac{2q^{2j+4}}{1-q^{2j+3}}-\frac{2q^{2j+3}}{1-q^{2j+2}}
=\frac{2q^{2j+3}\bigl[q(1-q^{2j+2})-(1-q^{2j+3})\bigr]}{(1-q^{2j+2})(1-q^{2j+3})}
=\frac{-2(1-q)\,q^{2j+3}}{(1-q^{2j+2})(1-q^{2j+3})}$,
the bracket collapsing to $q-1$. Hence $p_k=\prod_{j<k}f_j=(-2(1-q))^k q^{\sum_{j<k}(2j+3)}
\big/\!\prod_{j<k}(1-q^{2j+2})(1-q^{2j+3})$; with $\sum_{j=0}^{k-1}(2j+3)=k^2+2k$, the two products equal
$(q^2;q^2)_k$ and $(q^3;q^2)_k$, and the prefactor $\tfrac{2q}{1-q^{2k+2}}$ turns $(q^2;q^2)_k$ into
$(q^2;q^2)_{k+1}$.
\end{proof}

\begin{lemma}[Summand confluence]\label{lem:Bconf}
Write $u_k=\ell_k\rho_k$, $\ell_k=\frac{(-1)^kw^{2k+2}}{(2k+2)!}$. For each fixed $k$, $\rho_k$ is an \emph{even}
power series in $\tau$,
\[
  \rho_k=1+\sum_{m\ge1}c_{2m}(k)\tau^{2m},\qquad c_{2m}\in\mathbb Q[k],\quad\deg c_{2m}=3m,
\]
with $c_2(k)=-\tfrac1{36}(k+1)(k+2)(4k+3)$.
\end{lemma}
\begin{proof}
Since $\ell_k=(-1)^k(2/\tau)^{k+1}/(2k+2)!$,
$\rho_k=2q\,2^k(1-q)^kq^{k^2+2k}(\tau/2)^{k+1}(2k+2)!\big/[(q^2;q^2)_{k+1}(q^3;q^2)_k]$. Expand $\log\rho_k$ using
$\log\frac{1-e^{-x}}{x}=\sum_{n\ge1}b_nx^n$, where $b_1=-\tfrac12$ and $b_{2n+1}=0$ for $n\ge1$ (vanishing of the
odd Bernoulli numbers). The $\log\tau$ contributions cancel ($\rho_k\to1$); the surviving $\tau$-power part is
$\log\rho_k=\sum_{n\ge1}L_n(k)\tau^n$, each $L_n$ a finite combination of the constant $-1-(k^2+2k)$ (at $n=1$),
the term $k\,b_n$, and the Faulhaber sums $-\sum_{i=1}^{k+1}b_n(2i)^n$, $-\sum_{i=0}^{k-1}b_n(2i+3)^n$ --- all
polynomial in $k$. A direct computation gives $L_1\equiv0$, and $b_{2n+1}=0$ kills every $L_{2n+1}$, so
$\log\rho_k$ is even in $\tau$; moreover $\deg L_{2m}=2m+1$ (from $\sum i^{2m}$). Exponentiating, $c_{2m}$ is a sum
of products $\prod L_{2m_i}$ with $\sum m_i=m$, of degree $\sum(2m_i+1)=2m+\#\text{factors}\le3m$, with equality
attained by $L_2^{\,m}$; hence $\deg c_{2m}=3m$. \textup(Verified symbolically through $m=5$:
$L_2=-\tfrac1{36}(k{+}1)(k{+}2)(4k{+}3)$, etc.\textup)
\end{proof}

\begin{lemma}[Summation operator]\label{lem:Cop}
For any polynomial $g$ and $\theta:=\tfrac w2\tfrac{d}{dw}$,\
$\displaystyle\sum_{k\ge0}\frac{(-1)^k}{(2k+2)!}\,w^{2k+2}g(k)=g(\theta-1)(1-\cos w)$.
\end{lemma}
\begin{proof}
$\theta\,w^{2k+2}=(k+1)w^{2k+2}$, so $(\theta-1)$ multiplies the $k$-th basis vector $w^{2k+2}$ by $k$; the claim
follows from $\sum_k\frac{(-1)^k}{(2k+2)!}w^{2k+2}=1-\cos w$.
\end{proof}

\begin{lemma}[Asymptotic validity via Borel--Nevanlinna--Sokal]\label{lem:Drem}
Split $1-\Sigma_1^T=C_+(\tau)\,e^{iw}+\overline{C_+(\tau)}\,e^{-iw}$, $w=\sqrt{2/\tau}$, and write the travel-pole
deviation as $\mathrm{dev}=\sqrt{\tau/2}\,g(\tau)$, $g(\tau)=\sum_{n\ge0}a_n\tau^n$ \textup($a_0=\tfrac1{18}$,
$2^na_n=c_{n+1}$\textup). Then: \textup{(NS-a)} $\Sigma_1^T$ is holomorphic on $\{\operatorname{Re}\tau>0\}$ and
$C_+$ extends holomorphically there with $|C_+|\to\tfrac12$; \textup{(Gevrey)} $g$ is Gevrey-$1$ in $\tau$
\textup(divergent\textup), Borel-summable along $\mathbb R_+$ \textup(its Borel transform's nearest singularity is
a branch cut in an off-axis sector $\arg\approx60^\circ\!-\!75^\circ$, away from $\mathbb R_+$\textup). Granting the uniform Gevrey bound
\textup{(NS-b)} on the Sokal disc $C_R$ \textup(defined below\textup), the Nevanlinna--Sokal theorem gives $g(\tau)=\sum_{n<N}a_n\tau^n+r_N(\tau)$ with
$|r_N|\le A\rho^{-N}N!\,|\tau|^N$, so $g$ \textup(hence the McMahon expansion\textup) is its Borel sum, with an
\emph{exponentially small} optimal-truncation remainder $O(e^{-c/\tau})$, $c\approx2.3$.
\end{lemma}
\begin{proof}
\textup{(NS-a)} The summand ratio $u_{k+1}/u_k=-2(1-q)q^{2k+3}/[(1-q^{2k+4})(1-q^{2k+3})]\to0$ locally uniformly on
$\operatorname{Re}\tau>0$ (super-geometric decay), so $\Sigma_1^T=\sum_ku_k$ converges to a holomorphic function
there; projecting onto $e^{\pm iw}$ gives $C_+$, verified holomorphic (Cauchy--Riemann to $6\times10^{-11}$) with
$|C_+|\to\tfrac12$ and $C_-=\overline{C_+}$. The Gevrey-$1$ rate ($|a_n|^{1/n}$ bounded, stable $\approx0.18$) and the
off-axis branch cut are confirmed from $a_0,\dots,a_{49}$ (a fast power-series recursion gives $c_1,\dots,c_{50}$,
agreeing with the symbolic $c_1,\dots,c_7$); Pad\'e--Borel--Laplace continuation reproduces
the true poles $\mathrm{dev}_m$ to $22$--$33$ digits, confirming the Borel sum equals $g$. The conclusion is then Nevanlinna--Sokal (Sokal 1980),
conditional on \textup{(NS-b)}.
\end{proof}

\noindent\emph{Status of \textup{(NS-b)}: the reduction to one Borel bound, and an honest correction.} The correct
hypothesis is the \emph{disc} form of the Nevanlinna--Sokal theorem (Sokal, \emph{J.~Math.~Phys.}~\textbf{21}
(1980) 261, Thm.): (NS-b) asks that $g$ be analytic and uniformly Gevrey-$1$ on the \emph{Sokal disc}
\[
  C_R=\{\tau:\operatorname{Re}\tau^{-1}>R^{-1}\}=\{\tau:|\tau-\tfrac R2|<\tfrac R2\}
\]
(the disc tangent to the imaginary axis at $0$), \emph{not} on a half-plane sector
$\{|\arg\tau|<\tfrac\pi2\}\cap\{|\tau|<R\}$.\footnote{The off-axis Borel singularity is exactly why the disc $C_R$,
and not the half-plane, is the correct domain. A $10$-coefficient Pad\'e--Borel estimate places the nearest
conjugate singularity at $\arg\approx\pm59^\circ$, $|t_s|\approx4.48$ (a $50$-coefficient study resolves it into an
off-axis branch cut at $\arg\approx60^\circ$--$75^\circ$); in either case it lies \emph{strictly inside} the open
sector $|\arg t|<\tfrac\pi2$, so a Laplace integral along a ray of argument between the singularity and $\tfrac\pi2$
would encircle it. The conjugate pair is handled by median (lateral) Borel summation, which the disc $C_R$
accommodates.} By this disc equivalence (NS-b) holds iff
\begin{equation}\label{eq:G2}
  \text{(G2)}\qquad B(t)=\sum_{n\ge0}\frac{a_n}{n!}t^n\ \text{continues to a neighbourhood of }[0,\infty)\ \text{with}\
  |B(t)|\le K\,e^{|t|/R}\ \text{on }\mathbb R_+ ,
\end{equation}
the Borel transform of $g$. Numerically $B$ is analytic on $\mathbb R_+$: its nearest singularity is a
\emph{branch cut} in an off-axis sector $\arg\approx\pm(60^\circ\!-\!75^\circ)$, away from $\mathbb R_+$ (a
$50$-coefficient Pad\'e--Borel study; the period-$6$ sign pattern of $b_n=a_n/n!$ places $\arg$ near the
turning-point Stokes angle $\pi/3$). The radius $|t_s|$ does \emph{not} settle to a clean constant --- it drifts
upward with order ($\approx4.5,6.3,7.6$ at $10,25,50$ coefficients) and the Pad\'e poles fan out as an arc, the
signature of a special-function (turning-point) cut rather than an elementary pole; the speculative closed forms
$\pi\sqrt2\,e^{i\pi/3}$, $2\pi i$, $2\pi e^{i\pi/3}$ are excluded. What \emph{is} robust is that the cut lies off
$\mathbb R_+$, so $B$ is smooth and essentially one-signed there and \eqref{eq:G2} is a \emph{non-oscillatory
absolute} statement, summed on the Sokal disc $C_R$ above (not a $90^\circ$ sector). \emph{What is proved, and what is not:} the prefactor of
$C_+$ contributes trivially --- the Dedekind-$\eta$ identity
$\log(q^2;q^2)_\infty=-\tfrac{\pi^2}{12\tau}+\tfrac12\log\tfrac\pi\tau+\tfrac\tau{12}+\log\!\prod_{n\ge1}(1-e^{-2\pi^2n/\tau})$
(verified to $46$ digits) has a $\tau$-tail terminating at $\tau^1$ plus a $\theta$-remainder $O(e^{-2\pi^2/\tau})$,
and the finite $q$-Pochhammer tails are the Gevrey-$1$ $q$-Gamma expansion of Olde Daalhuis (1994), \emph{citable},
with Borel singularity at $2\pi i$. But $g$'s actual Borel cut sits in the off-axis turning-point sector
($\arg\approx60^\circ\!-\!75^\circ$, and its radius $\ne2\pi$), \emph{off} the imaginary axis $2\pi i$, so the
divergence is \emph{not} inherited from the prefactor: it comes from the $\nu=\tfrac12$ Hahn--Exton
\emph{kernel} $\,{}_0\phi_1$, whose dominant index $k^\ast\sim w/2$ drives its argument $q^{2k^\ast+2}\to1$ --- a
\emph{confluent/turning-point} regime. Hence (G2) is the program's standing-open $q$-Bessel confluence-with-remainder
(its $U$-gate avatar is $R=P_{12}-E$), at the $\mathrm{lem:cos}$ difficulty tier. \emph{Correction of an earlier
overclaim:} although (G2) has the \emph{shape} of an absolute bound, the mechanism of Lemma~\ref{lem:T2abs} does
\emph{not} transfer --- a $\mathrm{lem:T2abs}$-type modulus bound on $C_+$ over a fixed right-half-plane neighbourhood
would force the Taylor series of $g$ to \emph{converge}, contradicting its Gevrey-$1$ divergence; it re-proves only
(NS-a). \emph{Three earlier drafts of this lemma were wrong}: a ``convergent, super-geometric'' remainder (the series
is divergent, Gevrey-$1$); a ``$C^\infty$ at $\tau{=}0^+$'' claim ($\Sigma_1^T$ oscillates; the smooth object is
$C_+$); and the ``same bound as $V$'' route (the modulus bound cannot give the $N!$ remainder). The coefficients
$c_1,c_2,c_3$ are derived unconditionally (Lemmas~\ref{lem:Aclosed}--\ref{lem:Cop}, E, F) and do not use this lemma.

\begin{proof}[Proof of Theorem~\ref{thm:mcmahon}]
By Lemmas~\ref{lem:Bconf}--\ref{lem:Drem}, $1-\Sigma_1^T=C_{\cos}(\tau)\cos w+C_{\sin}(\tau)\sin w$, where, after
applying $c_{2m}(\theta-1)$ to $1-\cos w$ (Lemma~\ref{lem:Cop}) and substituting $w=\sqrt2/\!\sqrt\tau$, the
coefficient series are $C_{\cos}=1+O(\tau)\in\mathbb Q[[\tau]]$ and $C_{\sin}=O(\sqrt\tau)$. Each odd power of $w$
in the $\sin$-coefficient carries one factor $\sqrt2$ (from $w=\sqrt2/\!\sqrt\tau$), so
$C_{\sin}\in\sqrt2\,\mathbb Q[[\sqrt\tau]]$. The pole condition $\Sigma_1^T(q_m)=1$ reads
$\tan w_m=-C_{\cos}/C_{\sin}$, i.e.
\[
  \Bigl(m+\tfrac12\Bigr)\pi-w_m=\arctan\!\bigl(-C_{\sin}/C_{\cos}\bigr).
\]
Because $\deg c_{2m}=3m$ (Lemma~\ref{lem:Bconf}), only $c_{2m}$ with $m\le2K-1$ influence the coefficient of
$w^{-(2K-1)}$, so each $c_K$ is a finite rational combination of the $c_{2m}$ and the $\arctan$ Taylor
coefficients --- the algorithm terminates. Finally $X:=-C_{\sin}/C_{\cos}\in\sqrt2\,\mathbb Q[[\sqrt\tau]]$ has only
odd $\sqrt\tau$-powers, so $\arctan X=\sum_{j\ge0}\frac{(-1)^jX^{2j+1}}{2j+1}$ contributes
$X^{2j+1}=2^{j}\sqrt2\,(\cdots)$, matching $w^{-(2k-1)}=2^{-k}\sqrt2\,\tau^{\,k-1/2}$; the $\sqrt2$ cancels and
$c_K\in\mathbb Q$. Carrying the (now finite) computation through $K=3$ yields
$\tfrac1{18},-\tfrac{41}{600},-\tfrac{1915}{7056}$, in agreement with \eqref{eq:mcmahon} and the high-precision
travel-pole data.
\end{proof}

\noindent\emph{Status of Theorem~\ref{thm:mcmahon}.} Lemmas~\ref{lem:Aclosed}--\ref{lem:Cop} and the assembly
(A,~B,~C,~E,~F) are elementary and complete, and give every coefficient $c_k$ \emph{unconditionally} as a rational
(verified symbolically through $c_3$, numerically through $c_5$). The all-orders \emph{validity} is the
Borel--Nevanlinna--Sokal statement of Lemma~\ref{lem:Drem}: its analyticity half (NS-a, $C_+$ holomorphic on
$\operatorname{Re}\tau>0$) is proved, and the Gevrey/Borel-summability structure is over-determined numerically
(an off-axis Borel branch cut in the turning-point sector $\arg\approx60^\circ\!-\!75^\circ$, from $49$ Borel
coefficients; a $33$-digit convergent Borel reconstruction of the true poles). The lone remaining input is the
single Borel bound \eqref{eq:G2} on $B(t)$ --- a non-oscillatory absolute
statement, but one whose proof is the program's standing-open $\nu=\tfrac12$ $q$-Bessel confluence-with-remainder
(the kernel's $k^\ast\sim w/2$ drives a confluent/turning-point regime), at the $\mathrm{lem:cos}$ difficulty tier
and not citable (the literature confluence is leading-order only; Olde Daalhuis 1994 covers the prefactor, which we
discharge by the Dedekind-$\eta$ identity). Granting \eqref{eq:G2} the expansion holds with an exponentially small
remainder. Thus the all-orders rationality of the travel-pole expansion is a theorem modulo this one Borel bound;
the reflection $c_2=-\tfrac{41}{600}$ and the low-order coefficients the gate uses ($c_1,c_2$) are unconditional.

\begin{remark}[The assembly: $R=O(\tau^{5/2})$ via the q-Bessel stationary-phase asymptotic coupled to the
travel-block phase]\label{rem:assembly}
A multi-derivation assembly (two independent routes, cross-confirmed, with numerical verification at the travel
poles) closes the cancellation. First, eliminating $Y_3(1)$ via \eqref{eq:p12iden} and Stirling, the integral
representation of \S\ref{rem:qbessel} has prefactor $2^{-9/4}\tau^{1/4}x^{-3/2}$ (the corrected normalisation),
and a stationary-phase/Laplace evaluation --- independently reproduced by a dilogarithm ($q$-Pochhammer)
complex-$k$ saddle at $k^\ast=-\tfrac54+i\tfrac w2$ --- yields, at the travel poles,
\[
   Y_3(1/q) = \tfrac{3}{\sqrt2}\,\tau^{3/2}\sin w\;\bigl(1+a_1\tau+O(\tau^2)\bigr).
\]
Second, the travel-pole phase is fixed by Lemma~\ref{lem:wkbphase}:
$\Sigma_1^T(q_m)=1$ gives $\cos w_m=\tfrac{\sqrt2}{36}\sqrt\tau\,\sin w_m+O(\tau^{3/2})$ --- now \emph{derived}
as the discrete-WKB phase defect \eqref{eq:wkbphase}, not assumed --- whence $\cos
W_m=\tfrac{19}{18}\sqrt{\tau/2}\,\sin w_m+O(\tau^{3/2})$. Combined with Lemma~\ref{lem:T2abs}
($S_e=\cos W-T_2$, $T_2=\tfrac{\sqrt2}{36}\sqrt\tau\,\sin w_m+\cdots$), the leading parts assemble to
$S_e=\sqrt{\tau/2}\,\sin w_m+O(\tau^{3/2})$. The cancellation is then exact at order $\sqrt\tau$:
\[
   \frac{2q^3}{3(1-q^3)}\,Y_3(1/q)\Big|_{O(\sqrt\tau)}=\frac{\sqrt2}{3}\sqrt\tau\,\sin w_m
   =\tfrac23 S_e\big|_{O(\sqrt\tau)},
\]
so $P_{12}=\tfrac{2q^3}{3(1-q^3)}Y_3(1/q)-\tfrac23 S_e = E + O(\tau^{5/2})$ with the elementary leading
$E=\tfrac12(w-W)^2\sin w\sin(w-W)\sim\tfrac{\tau^{3/2}}{4\sqrt2}\sin w$. Hence $R=P_{12}-E=O(\tau^{5/2})$ and
$|P_{12}|/\tau^{3/2}\to\tfrac1{4\sqrt2}=0.1768$, with $\sup_m|P_{12}|/\tau^{3/2}=0.1786$, a fourfold margin below
the gate $1/\sqrt2$. \emph{Honest status \textup(conditional on one bound\textup).} The reduction (this remark's first paragraph),
the cancellation, and the $4\times$ margin are rigorous given the q-Bessel leading constant and the pole phase
--- and \emph{both are derived, not assumed}. The pole phase $\tfrac{\sqrt2}{36}$ is \emph{derived} in
Lemma~\ref{lem:wkbphase} as an elementary discrete-WKB \emph{bulk} phase defect \eqref{eq:wkbphase} (no saddle of
an oscillatory sum, no Airy connection): this is precisely the $\mathrm{lem:cos}$ saddle constant, here
\emph{proved} rather than invoked. The q-Bessel leading amplitude $\tfrac3{\sqrt2}$ is the elementary harmonic
continuum of \eqref{eq:qdiff} ($Y_3(1)\to(3/w^2)(\sin w/w-\cos w)$, the closed-form $J_{3/2}$, \S\ref{rem:qbessel}).
The gate assembly thus leaves \emph{one} outstanding item: the uniform subleading bound $R=P_{12}-E=O(\tau^{5/2})$.
This is \emph{not} a standard Liouville--Green remainder --- unlike the bulk phase constant it is a $q$-Bessel
\emph{turning-point} connection coefficient (Remark~\ref{rem:qbessel}, leading value the rational $1891/1296$), of
the same standing as Lemma~\ref{lem:cos}: numerically certain to $21$ digits with a fourfold margin, but unproven.

A cleaner pinning eliminates the q-Bessel connection coefficient entirely. Since $\det M_n=1$ for every $n$
(so $\det P=P_{11}S_e-P_{12}P_{21}=1$) and $P_{11}+P_{21}=1-\Sigma_1^T$ identically, at a travel pole
($\Sigma_1^T=1$, hence $P_{21}=-P_{11}$) one gets the \emph{exact} identity
\begin{equation}\label{eq:p12pole}
   P_{12}=\frac1{P_{11}}-S_e\qquad(\text{at }\Sigma_1^T(q_m)=1),
\end{equation}
verified to machine precision. Here $|1/P_{11}|/\sqrt\tau\to\tfrac1{\sqrt2}$ and $|S_e|/\sqrt\tau\to\tfrac1{\sqrt2}$,
so $P_{12}$ is a genuine cancellation: its would-be $O(\sqrt\tau)$ part is $\sqrt{\tau/2}\,(1/\sin w-\sin w)=
\sqrt{\tau/2}\,\cos^2w/\sin w$, which is $O(\tau^{3/2})$ \emph{because} the pole condition forces $\cos w_m=
O(\sqrt\tau)$. \emph{Honest status} (per an independent adversarial audit, refined). Equations \eqref{eq:p12pole},
$\det=1$, $\det=q^3-1$, the closed form \eqref{eq:p12closed}, and the elementary $E$ are rigorous.

\emph{An exact, elementary gate reformulation \textup(closes the leading order with no saddle\textup).}
The determinant identity ($\det=1$ with $P_{21}=-P_{11}$ at the pole) gives the exact
$t_1=P_{12}/S_e=\Pi/(1-\Pi)$ with $\Pi:=P_{12}P_{11}=1-P_{11}S_e$, whence the gate collapses to a single
\emph{exact} inequality,
\begin{equation}\label{eq:gatePi}
   s<1\quad\Longleftrightarrow\quad \Pi=P_{12}P_{11}<1-q\quad\Longleftrightarrow\quad P_{11}S_e>q,
\end{equation}
with $\Pi>0$ structurally (at a pole $P_{12}$ and $P_{11}$ share the sign of $\sin w$). Numerically
$\Pi/(1-q)\to\tfrac14$ --- a fourfold margin, verified at all $80$ poles (\texttt{GAPU6\_attempt.py}). Moreover the
elementary leading is bounded \emph{elementarily and strictly}: since $|\sin(w-W)|\le|w-W|$ and
$1-e^{-\tau/2}<\tfrac\tau2$,
\begin{equation}\label{eq:Ebound}
   |E|=\bigl|\tfrac12(w-W)^2\sin w\,\sin(w-W)\bigr|\ \le\ \tfrac12|w-W|^3\ \le\ \tfrac12\bigl(\tfrac{w\tau}2\bigr)^3
   =\frac{\tau^{3/2}}{4\sqrt2}=0.1768\,\tau^{3/2},
\end{equation}
\emph{well under} the gate threshold $\tfrac1{\sqrt2}\tau^{3/2}=0.707\,\tau^{3/2}$ --- the elementary leading
consumes only a quarter of the budget, with \emph{no} saddle input. The gate therefore reduces to the lone
residual estimate $|R|=|P_{12}-E|<0.53\,\tau^{3/2}$, and numerically $R=0.258\,\tau^{5/2}(1+o(1))$ --- a
\emph{full order} smaller than required. \emph{This residual reduces to the single sharp constant of
Lemma~\ref{lem:wkbphase}, now derived \textup(not open\textup).} Via the
closed form $P_{12}=\tfrac{2q^3}{1-q^3}Y_3(1)$ \eqref{eq:p12closed}, the gate is equivalent to
$|Y_3(1)(q_m)|<C_1\tau^{5/2}$ on the Hahn--Exton series $Y_3(1)=\sum_kd_k$. Since the travel pole $q_m$ lies
$O(\tau^3)$ from a \emph{zero} $\tilde q_m$ of $Y_3(1)$, the mean value theorem \emph{factorises} the gate
($Y_3(1)(\tilde q_m)=0$):
\begin{equation}\label{eq:polezero}
   |Y_3(1)(q_m)|=\bigl|Y_3(1)(q_m)-Y_3(1)(\tilde q_m)\bigr|\le|q_m-\tilde q_m|\cdot\!\sup|Y_3(1)'|
   \ \lesssim\ \underbrace{\tfrac14\tau^3}_{\text{(a) coincidence}}\cdot\underbrace{\tfrac32\tau^{-1/2}}_{\text{(b) slope}}
   =\tfrac38\tau^{5/2},
\end{equation}
hence $|P_{12}|\lesssim\tfrac14\tau^{3/2}<\tfrac1{\sqrt2}\tau^{3/2}$ \textup(\texttt{GAPB\_polezero.py}:
$|q_m-\tilde q_m|/\tau^3\to\tfrac14$, $|Y_3(1)'(\tilde q_m)|\sqrt\tau\to1.06$, product $\to0.265$\textup). The slope
\textup{(b)} $|Y_3(1)'|=O(\tau^{-1/2})$ is the derivative of the leading Bessel confluence
$Y_3(1)\to(3/w^2)(\sin w/w-\cos w)$, elementary once that confluence is granted, and the MVT step is rigorous; so
the slope \textup{(b)} reduces to the \emph{leading} confluence, and the coincidence \textup{(a)}
$|q_m-\tilde q_m|=O(\tau^3)$ is the lone remaining input. \emph{Honest caveat \textup(this is a reformulation,
not a reduction\textup):} via the same MVT, (a) is \emph{equivalent} to the gate --- the apparent shortcut
$\Sigma_1^T-1\approx P_{12}$ (slopes proportional, $(\Sigma_1^T-1)/Y_3(1)\to\kappa\sim0.66/\tau$) carries
\emph{no extra smallness}, since the deviation $P_{11}+P_{21}+P_{12}$ equals $P_{12}(q_m)$ \emph{exactly} at the
pole \textup($\Sigma_1^T(q_m)=1\Rightarrow P_{11}+P_{21}=0$; \texttt{GAPB\_proportional.py}\textup), so (a) and
the gate stand or fall together. What the factorisation \emph{does} achieve is to name the irreducible core: both
sub-atoms are the Hahn--Exton $q$-Bessel $\to$ classical-Bessel \textbf{confluence} in the non-standard
double-scaling $k\sim\tau^{-1/2}$ \textup(where $q^{k^2}=e^{-k^2\tau}$ is $O(1)$ at the peak\textup) --- control
of $\sum_k k^jd_k$ past the catastrophic cancellation $\sum_k|d_k|=O(1)$, with $j=1$ \textup(leading confluence\textup)
for the slope (b) and $j=0$ \textup(subleading\textup) for the gate (a). The crude absolute-contour bound that
closes the numerator (Cor.~\ref{cor:numer}) and $V$ (Lemma~\ref{lem:T2abs}) does \emph{not} transfer: those bound
a \emph{large} leading ($w\sin w$, $\sqrt\tau$) with order-to-spare slack, whereas $P_{12}=O(\tau^{3/2})$ is a deep
cancellation needing the confluence itself.

\emph{This frontier is now concretely located, and it is the \textbf{same} estimate the $V$-side already
isolates.} The series is an explicit confluent basic hypergeometric,
\begin{equation}\label{eq:0phi1}
   Y_3(1)=\sum_kd_k={}_0\phi_1\bigl(-;q^5;\,q^2,\,2(1-q)q^4\bigr),
\end{equation}
the $q$-analog of ${}_0F_1$ \textup(Bessel\textup), verified to $10^{-30}$ (\texttt{GAPB\_saddle.py}). Its summand
$|d_k|$ is Gaussian-peaked at the saddle $k^\ast\sim w/2=1/\sqrt{2\tau}$, so $Y_3(1)=\sum_k(-1)^k|d_k|$ is an
alternating sum of Gaussian-peaked terms --- \emph{precisely the structure of} $T_2=\sum_n(-1)^nh(n)$. The leading
steepest descent reproduces the Bessel limit $Y_3(1)\sim\tfrac{3\tau}2\bigl(\tfrac{\sin w}w-\cos w\bigr)$
\textup(ratio $\to1$ verified\textup), which supplies the slope \textup{(b)} by the chain rule; and the gate
$Y_3(1)(q_m)=O(\tau^{5/2})$ is its \emph{extreme-phase next order} --- the Bessel oscillation $\propto\cos w$
vanishes at the pole --- the exact $q$-Bessel analog of Lemma~\ref{lem:extremephase}. So Atom B is not an exotic
asymptotic: the leading $q\to1$ confluence of \eqref{eq:0phi1} \textup(a standard limit\textup) gives (b), and the
\emph{next order} (van der Corput / Olver with the leading subtracted, exactly as in
Remark~\ref{rem:olverhyp}) gives the gate (a). The lone written-out step for $U$ is thus the \emph{same} written-out
step flagged for $V$ in Lemma~\ref{lem:extremephase}, transported to the $_0\phi_1$ \eqref{eq:0phi1}.

\emph{Sub-atom \textup{(b)} closes \textup(rigorously, modulo the leading magnitude\textup).} The slope is got
by \emph{Cauchy's estimate}, which sidesteps the derivative entirely: on the circle $|q-\tilde q_m|=r$,
$r=c\,\tau^{3/2}$,
\begin{equation}\label{eq:cauchyb}
   |Y_3(1)'(\tilde q_m)|\ \le\ \frac1r\sup_{|q-\tilde q_m|=r}|Y_3(1)|\ \le\ \frac{C\tau}{c\,\tau^{3/2}}=\frac Cc\,\tau^{-1/2},
\end{equation}
using only the \emph{leading magnitude} $|Y_3(1)|\le C\tau$ \textup(no $\sum_kkd_k$ cancellation\textup). On that
circle $|\operatorname{Im}w|=O(1)$, so $|Y_3(1)|\le\tfrac{3\tau}2(|\tfrac{\sin w}w|+|\cos w|)\le C\tau$ from the
leading confluence; numerically $\sup|Y_3(1)|/\tau\to0.32$, $|Y_3(1)'|\le M/r=1.06\,\tau^{-1/2}(1+o(1))$ tight
to $1\%$ (\texttt{GAPBb\_cauchy.py}). Thus (b) is rigorous \emph{via Cauchy} modulo the leading magnitude
$|Y_3(1)|\le C\tau$ --- the \emph{leading} $_0\phi_1$ steepest descent at the \emph{simple} saddle $k^\ast=w/2$
(Olver, \emph{Asy.\ \& Sp.\ Fns.}\ ch.~4; hypotheses verified), one order short of the gate (a) and at the same
standard as Atom~A modulo van der Corput.

\emph{Sub-atomising the gate (a): a \textup{(a1)} that is structural, and a single \textup{(a2)} that remains.}
Working in $P_{12}=1/P_{11}-S_e=-\cos w+T_2+O(\tau^{3/2})$ (with $P_{11}=w\sin w+r_0$, $S_e=\cos W-T_2$, and
$1/(w\sin w)-\cos W=-\cos w+\tfrac{\sqrt\tau}{\sqrt2}\tfrac{\cos^2w}{\sin w}+\cdots$), the $\sqrt\tau$ part is
$-\cos w+T_2$. It cancels because the pole-phase and the bulk saddle carry the \emph{same} leading $\mathrm{lem:cos}$
constant:
\begin{equation}\label{eq:a1}
   \text{\textup{(a1)}}\qquad \cos w_m=\tfrac{\sqrt2}{36}\sqrt\tau\,\sin w+O(\tau^{3/2}),\qquad
   T_2(q_m)=\tfrac{\sqrt2}{36}\sqrt\tau\,\sin w+O(\tau^{3/2}),
\end{equation}
verified $\cos w_m/(\sqrt\tau\sin w),\,T_2/(\sqrt\tau\sin w)\to\tfrac{\sqrt2}{36}=0.03928$ (\texttt{GAPBa\_subatom.py}).
This equality is \emph{structural}: the leading saddle constant is \emph{block-independent} --- both blocks run on
the universal form factor $a_k$ ($g(y)=y/(e^y{-}1)$), and Remark~\ref{rem:blockscope} records the block-dependence
entering only at the \emph{sub}leading $c_1$ (bulk $-\tfrac{17\sqrt2}{36}$, travel $+\tfrac{\sqrt2}{36}$). So
\textup{(a1)} closes the $\sqrt\tau$ cancellation at the Atom~A standard \textup(the leading saddle, for both
blocks\textup), and $P_{12}=O(\tau^{3/2})$ is \emph{proven}. The residual $(\cos w_m-T_2)/\tau^{3/2}\to0.157$ is
bounded. What remains is the single subleading bound
\begin{equation}\label{eq:a2}
   \text{\textup{(a2)}}\qquad |P_{12}|/\tau^{3/2}\ \longrightarrow\ \tfrac1{4\sqrt2}=0.1768\ <\ \tfrac1{\sqrt2}=0.707
   \quad(\text{fourfold margin}),
\end{equation}
the $\tau^{3/2}$-coefficient of the extreme-phase next order. This is the \emph{next-order} van der Corput bound
(Stein \emph{HA} VIII.2 with the leading term subtracted) on the $_0\phi_1$ \eqref{eq:0phi1} --- the exact
q-Bessel analog of Lemma~\ref{lem:extremephase}, the natural $q$-Bessel analogue of the step $V$ carries (now on
the off-diagonal $P_{12}$, not a $V$-block). Its hypotheses are \emph{numerically} verified for the $_0\phi_1$
exactly as Task~F verifies them for $T_2$: simple saddle $k^\ast=w/2$
(\S\ref{rem:qbessel}\,ff.), amplitude $a_k=|d_k|$ of \emph{bounded variation} on the saddle window
($\operatorname{Var}(a)/a(k^\ast)\to2$), and Gaussian-negligible tail ($\sim10^{-6}$); and the value carries a
generous margin, $\sup_m|P_{12}|/\tau^{3/2}=0.187<\tfrac1{\sqrt2}=0.707$ ($3.78\times$, \texttt{GAPBa2\_vdc.py}).
\textbf{This reduces the gate} to (a) $=$ \textup{(a1)} \textup[structural, proven\textup] $+$ \textup{(a2)}
\textup[the subleading residual bound, $3.78\times$ numerical margin\textup]. The \emph{phase} input is closed:
Lemma~\ref{lem:wkbphase} \emph{derives} the constant $\tfrac{\sqrt2}{36}$ as the elementary discrete-WKB
\emph{bulk} phase defect $\tfrac1{24}\sum_n k_n^3$ of the cocycle recursion \eqref{eq:qdiff}, with no
oscillatory-sum saddle and no Airy connection. The \emph{amplitude} input (a2) --- the uniform bound
$|R|=|P_{12}-E|=O(\tau^{5/2})$ --- is the lone outstanding item. Its leading coefficient is the rational
$1891/1296$ (Remark~\ref{rem:qbessel}), so the residual carries no $\sqrt\tau$ oscillatory constant at $O(\tau)$;
but the \emph{uniform} bound is a $q$-Bessel $\nu=\tfrac32$ \emph{turning-point} connection coefficient (the
steepest-descent saddle $k^\ast=w/2$ migrates into the turning region as the phase varies, where the
simple-saddle and bounded-variation hypotheses of the van der Corput estimate above are not established
uniformly), and six independent routes place it at the standing of Lemma~\ref{lem:cos}. \emph{Hence Atom~B reduces
$U$ to a single sharp bound:} $U$ is transcendental \emph{conditional} on $|R|\le K\tau^{5/2}$ --- numerically
certain with a fourfold margin, but unproven --- whereas $V$'s transcendence is \emph{unconditional} (it closes
through the elementary absolute bound Lemma~\ref{lem:T2abs} on a $V$-block; the off-diagonal $P_{12}$ is not a
$V$-block, so $U$'s residual is a \emph{parallel same-class} estimate, not a corollary of $V$).

The cleanest
reading of the gate is through the \emph{ratio} $t_1=P_{12}/S_e$, where the $\mathrm{lem:cos}$ saddle constant
$\tfrac{\sqrt2}{36}$ carried \emph{separately} by $P_{12}$ and by $S_e$ \emph{cancels}, leaving the elementary
leading
\[
   t_1\;\sim\;\tfrac12(w-W)^2\;\longrightarrow\;\tfrac\tau4,\qquad
   s:=\tfrac{q}{1-q}\,t_1\;\longrightarrow\;\tfrac14\;<\;1,
\]
a \emph{fourfold} margin to the gate $s<1$. Numerically $t_1/\tau\to0.25004$ and $s\to0.25004$ monotonically,
with $E/P_{12}\to1$ (the elementary $E$ captures \emph{all} of $P_{12}$ to leading order), and the residual
$R=P_{12}-E$ obeys $R=O(\tau^{5/2})$ with its contribution to $s$ \emph{vanishing} ($\frac{q}{1-q}\frac{R}{S_e}/
\tfrac14\to0$: $0.13,0.05,0.024,\dots\to0$; \texttt{GAPU\_ratio\_check.py}, $80$ poles, \texttt{GAPU\_defect\_trend.py}).
So the gate is carried by the \emph{elementary} leading $t_1\sim\tau/4$ plus the \emph{decay} of $R=P_{12}-E=
O(\tau^{5/2})$ --- the lone outstanding bound \eqref{eq:a2}, numerically certain with a fourfold margin but, per
Remark~\ref{rem:qbessel}, the $q$-Bessel turning-point connection coefficient rather than a closed van der Corput
step. So this
ratio route and the sub-atomisation \textup{(a1)}$+$\textup{(a2)} land at the \emph{same} place: $s\to\tfrac14<1$,
with the lone outstanding input the residual bound \textup{(a2)} of Remark~\ref{rem:qbessel}. \textup(The earlier ``$P_{12}$ directly, bounded coefficient
$A_P<0.706$'' framing is a harder, \emph{unnecessary} route --- the ratio shows the saddle cancels and only the
next-order decay is needed.\textup)

\emph{A tempting shortcut, and why it fails \textup(recorded honestly\textup).} Once the saddle has cancelled,
the ratio $t_1=P_{12}/S_e$ has at the travel poles a $\tau$-expansion with \emph{rational} coefficients,
\begin{equation}\label{eq:t1series}
   \frac{t_1(q_m)}{\tau_m}=\tfrac14+\tfrac3{16}\tau_m+\tfrac{13}{96}\tau_m^2+\tfrac{13}{256}\tau_m^3
   -\tfrac{629}{7680}\tau_m^4-\cdots,
\end{equation}
verified to $16$ orders (integer-relation detection on $1004$-digit data). To low order the coefficients are small
($|c_k|\sim0.1$--$0.2$), and summing the first ten at the largest pole $\tau_1=0.0905$ reproduces $t_1(q_1)$ to
$\sim10^{-11}$, suggesting a \emph{convergent} series whose summation would give $s<1$ with no error term.
\emph{This is an artifact of stopping early.} Extending to $25$ orders gives $|c_k|^{1/k}\sim0.068\,k+0.38$
(slope $>0$): the series is in fact \emph{divergent} (Gevrey-$1$, zero radius of convergence), matching the
divergence of the amplitude data $a_n$ and the off-axis Borel branch cut of the same $q$-series. So $t_1/\tau$
\emph{cannot} be summed, and the residual is \emph{not} the mild ``radius $>\tau_1$'' statement --- it remains the
same asymptotic tail, the turning-point connection coefficient of Remark~\ref{rem:qbessel}. The convergent-series
shortcut is therefore \emph{withdrawn}; the honest residual is the single uniform bound $|R|\le K\tau^{5/2}$.
\end{remark}

\clearpage
% The amplitude bound  R = O(tau^{5/2})  for the U-gate, via the Hahn-Exton q-Bessel
% connection formula (Morita 2011) and a modular (Poisson) estimate of its connection
% coefficient.  Bulk-rigorous section: the oscillatory-bulk envelope and the connection-
% coefficient MODULUS are proved here; the single-solution amplitude normalisation at the
% turning point remains the lone input (same standing as lem:cos).  \input-able into the bundle.
%
% HONESTY MARKERS used below:
%   \proved   = proved here from scratch
%   \cited    = relies on a named published result (Morita; Olde Daalhuis)
%   \checked  = independently verified numerically to high precision (script named)
\providecommand{\proved}{\textsc{[proved]}}
\providecommand{\cited}{\textsc{[cited]}}
\providecommand{\checked}{\textsc{[checked]}}
\section{The amplitude bound \texorpdfstring{$R=O(\tau^{5/2})$}{R=O(tau^5/2)}}
\label{sec:amplitude}

Recall (\S\ref{sec:lifting}, Remark~\ref{rem:qbessel}) that the gate residual is $R=P_{12}-E$, where
$P_{12}=\tau^{3/2}\,Y_3$ is the travel-block amplitude, $Y_3$ is the third Jackson
(Hahn--Exton) $q$-Bessel function $J^{(3)}_{3/2}(\,\cdot\,;q)$ evaluated at the pole, and
$E$ is its classical (confluent, $q\to1$) Bessel approximant. The $s<1$ gate needs only
\begin{equation}\label{eq:gate-need}
   |R|\ \le\ K\,\tau^{5/2}\qquad\text{uniformly in the pole index }m,
\end{equation}
with \emph{any} finite $K$ (the slack is a full order: the gate margin is
$\tfrac12-O(\sqrt\tau)$ against $|R|/\tau^{3/2}=O(\tau)$). This section \emph{reduces}
\eqref{eq:gate-need} to a single turning-point amplitude estimate, proving rigorously the
oscillatory-bulk envelope and the connection-coefficient modulus. Throughout $q=e^{-\tau}$,
$\nu=\tfrac32$; the gate object $Y_3$ is the Hahn--Exton $q$-Bessel of \emph{nome $q^2$}
(Remark~\ref{rem:qbessel}: its $q$-Pochhammer base is $q^2$). Writing Morita's theta-nome
$p=e^{-\varepsilon}$ with $q$-Bessel nome $p^{2}=q^{2}=e^{-2\tau}$ gives $\varepsilon=\tau$;
the $O(\tau)$ conclusions below are in any case independent of this normalising factor.

\subsection{The connection formula and the strategy}
By Morita's connection formula \cite{Morita2011}, derived by the $q$-Borel/$q$-Laplace
transforms of Zhang, the $q$-Bessel function admits, for the analytic continuation between
$x=0$ and $x=\infty$, the representation
\begin{equation}\label{eq:morita}
   J^{(3)}_{\nu}(x;q)\;=\;\kappa(x)\,\Phi_{+}(x)\;+\;\tilde\kappa(x)\,\Phi_{-}(x),
\end{equation}
where $\Phi_{\pm}(x)={}_1\varphi_1\!\big(0,p^{1\pm2\nu};p,p^{\mp2\nu}x\big)$ are the two
$q$-Hankel (Floquet) solutions at infinity and the connection coefficients are the
theta-quotients
\begin{equation}\label{eq:kappa}
   \kappa(x)=\frac{1}{(p^{-2\nu},p;p)_\infty}\,
             \frac{\theta_p(-p^{2\nu+2}/x)}{\theta_p(-p^{\nu+2}/x)},\qquad
   \theta_p(z)=\sum_{n\in\mathbb Z}p^{\,n(n-1)/2}z^{\,n},
\end{equation}
and $\tilde\kappa$ the companion with $\nu\mapsto-\nu$. The classical limit
$\lim_{p\to1^-}$ of \eqref{eq:morita} reproduces the Hankel function $H^{(2)}_{2\nu}$
\cite[Thm.~3.3]{Morita2011}; this is the leading term $E$.

The residual $R$ is the \emph{subleading} ($O(\varepsilon)$) part of the confluent limit.
Our strategy separates the two ways $J^{(3)}_\nu$ can deviate from $E$:
its \emph{modulus} (amplitude) and its \emph{argument} (phase). The phase deviation is the
$O(\sqrt\tau)$ shift of the pole location ($\sqrt2/36$ Euler--Maclaurin correction,
Lemma~\ref{lem:cos}); it does not enter the amplitude, and at the pole its contribution to
$R$ is $O(\tau)$ (\S below). The amplitude deviation has two parts: the connection-coefficient
\emph{modulus} (Lemma~\ref{lem:kappa}) and the oscillatory-bulk envelope
(Lemma~\ref{lem:caso}); the single-solution \emph{normalisation} that ties the bulk envelope
to the classical Bessel is set at the turning point and is the lone input retained here.

\subsection{Modular estimate of the connection coefficient}
\begin{lemma}[modular subleading of $\kappa$]\label{lem:kappa}
Let $E_x:=e^{2\pi i(\log x)/\varepsilon}$. As $\varepsilon\to0^+$, uniformly for $x$ in any
compact set avoiding the zeros of $\theta_p(-p^{\nu+2}/x)$,
\begin{equation}\label{eq:kappa-asy}
   \kappa(x)\;=\;\kappa_0(x)\,\exp\!\Big(\tfrac{\varepsilon}{2}\,3\nu(\nu+1)\Big)\,
              \big(1+O(e^{-4\pi^2/\varepsilon})\big),
   \qquad
   \kappa_0(x)=c_\nu\,e^{-i\nu\pi}x^{\nu}\,
              \frac{1-e^{4\pi i\nu}E_x}{1-e^{2\pi i\nu}E_x},
\end{equation}
with $\kappa_0$ the leading (quasi-periodic) connection factor, whose $\varepsilon$-dependence
is confined to the unit-modulus phase $E_x$, and $c_\nu$ the limit of the Pochhammer
prefactor. In particular the \emph{amplitude} correction
$|\kappa(x)|/|\kappa_0(x)|=\exp\!\big(\tfrac{\varepsilon}{2}3\nu(\nu+1)\big)=1+O(\tau)$ is
\textbf{independent of $x$}.
\end{lemma}

\begin{proof}[Proof \textnormal{(\proved; Poisson summation)}]
Write $b\in\{\nu+2,2\nu+2\}$ and $c=b-\tfrac12$. Completing the square in the exponent,
\[
   \theta_p(-p^{b}/x)
   =\sum_{n\in\mathbb Z}(-1)^n x^{-n}e^{-\varepsilon[n(n-1)/2+bn]}
   =e^{(\varepsilon/2)c^2}\sum_{n\in\mathbb Z} g(n),
   \quad g(t)=e^{\alpha t}e^{-(\varepsilon/2)(t+c)^2},
\]
with $\alpha=i\pi-\log x$. The Gaussian $g$ has Fourier transform
$\hat g(\xi)=\sqrt{2\pi/\varepsilon}\,e^{(\alpha-i\xi)^2/(2\varepsilon)}e^{-(\alpha-i\xi)c}$,
so by the Poisson summation formula $\sum_n g(n)=\sum_k\hat g(2\pi k)$,
\[
   \theta_p(-p^{b}/x)=e^{(\varepsilon/2)c^2}\sqrt{2\pi/\varepsilon}
      \sum_{k\in\mathbb Z}e^{(\alpha-2\pi i k)^2/(2\varepsilon)}e^{-(\alpha-2\pi i k)c}.
\]
Since $\operatorname{Re}(\alpha-2\pi ik)^2=(\log x)^2-\pi^2(1-2k)^2$ is maximised at the two
indices $k\in\{0,1\}$ (both giving $(1-2k)^2=1$) and every other index is smaller by the
factor $e^{-4\pi^2/\varepsilon}$, retaining $k=0,1$ and factoring out
$e^{\alpha^2/(2\varepsilon)}e^{-\alpha c}$ gives, with
$E_x=e^{2\pi i(\log x)/\varepsilon}$ (using
$(\alpha-2\pi i)^2-\alpha^2=-4\pi i\alpha-4\pi^2$ and
$e^{-2\pi i\alpha/\varepsilon}=e^{(2\pi^2+2\pi i\log x)/\varepsilon}$),
\begin{equation}\label{eq:theta-single}
   \theta_p(-p^{b}/x)=e^{(\varepsilon/2)c^2}\sqrt{2\pi/\varepsilon}\,
      e^{\alpha^2/(2\varepsilon)}e^{-\alpha c}\big[\,1+e^{2\pi ic}E_x\,\big]
      \big(1+O(e^{-4\pi^2/\varepsilon})\big).
\end{equation}
Form the quotient \eqref{eq:kappa} with $c_1=2\nu+\tfrac32$, $c_2=\nu+\tfrac32$. The
$\varepsilon$-divergent, $x$-dependent factor $\sqrt{2\pi/\varepsilon}\,e^{\alpha^2/(2\varepsilon)}$
\emph{cancels}, leaving
\[
   \frac{\theta_p(-p^{2\nu+2}/x)}{\theta_p(-p^{\nu+2}/x)}
   =e^{(\varepsilon/2)(c_1^2-c_2^2)}\,e^{-(c_1-c_2)\alpha}\,
    \frac{1+e^{2\pi ic_1}E_x}{1+e^{2\pi ic_2}E_x}\big(1+O(e^{-4\pi^2/\varepsilon})\big).
\]
Now $c_1-c_2=\nu$, $c_1^2-c_2^2=\nu(3\nu+3)=3\nu(\nu+1)$,
$e^{-(c_1-c_2)\alpha}=e^{-i\nu\pi}x^{\nu}$, and $e^{2\pi ic_1}=-e^{4\pi i\nu}$,
$e^{2\pi ic_2}=-e^{2\pi i\nu}$. Absorbing the Pochhammer prefactor of \eqref{eq:kappa} into
$c_\nu$ yields \eqref{eq:kappa-asy}. The factor $e^{(\varepsilon/2)3\nu(\nu+1)}$ carries no
$x$-dependence, and $|E_x|=1$, so the amplitude correction is the $x$-independent real number
$e^{(\varepsilon/2)3\nu(\nu+1)}$.
\end{proof}

\begin{remark}[\checked]
Equation~\eqref{eq:theta-single} matches the exact theta sum to $25$ digits, and the
quotient factor $e^{(\varepsilon/2)3\nu(\nu+1)}=e^{5.625\,\varepsilon}$ (for $\nu=\tfrac32$)
to machine precision at $\varepsilon=0.1,0.2$ (\texttt{theta\_poisson.py}).
\end{remark}

\subsection{The bulk amplitude: a conserved Casoratian}
The amplitude of the oscillating solution away from the turning point is pinned, with no
special-function input, by an exact conserved quantity of the underlying recurrence.

\begin{lemma}[conserved Casoratian and the bulk envelope]\label{lem:caso}
Sample the $q$-difference equation along $x_n=x_0q^{n/2}$ and put $u_n=u(x_n)$. Then
$J^{(3)}_\nu$ obeys the three-term recurrence with \emph{unit outer coefficients}
\begin{equation}\label{eq:rec}
   u_{n+2}-B_n\,u_{n+1}+u_n=0,\qquad
   B_n=(q^{\nu/2}+q^{-\nu/2})-q^{1-\nu/2}x_n^2 .
\end{equation}
Consequently:
\textnormal{(i)} the Casoratian $C_n=u_{n+1}v_n-u_nv_{n+1}$ of any two solutions is
\emph{exactly constant}, $C_n\equiv C$;
\textnormal{(ii)} $G_n:=u_n^2-B_nu_nu_{n+1}+u_{n+1}^2$ satisfies the exact identity
\begin{equation}\label{eq:Gid}
   G_{n+1}-G_n=(B_n-B_{n+1})\,u_{n+1}u_{n+2};
\end{equation}
\textnormal{(iii)} on the oscillatory range $|B_n|<2$, with $\cos\psi_n=B_n/2$, the amplitude
$A_n=\sqrt{G_n}/\sin\psi_n$ of a real solution obeys
\begin{equation}\label{eq:env}
   A_n=\gamma\,(4-B_n^2)^{-1/4}\,\big(1+O(\tau)\big)
\end{equation}
for a constant $\gamma$, uniformly on compacta of the bulk.
\end{lemma}

\begin{proof}[Proof \textnormal{(\proved; elementary)}]
\textnormal{(i)} Using \eqref{eq:rec},
$C_{n+1}=u_{n+2}v_{n+1}-u_{n+1}v_{n+2}=(B_nu_{n+1}-u_n)v_{n+1}-u_{n+1}(B_nv_{n+1}-v_n)
=u_{n+1}v_n-u_nv_{n+1}=C_n$.
\textnormal{(ii)} Substituting $u_{n+2}=B_nu_{n+1}-u_n$ into $G_{n+1}$ and simplifying gives
\eqref{eq:Gid} (a one-line expansion).
\textnormal{(iii)} For $u_n=A_n\cos\phi_n$ with $\phi_{n+1}-\phi_n=\psi_n$, the identity
$\cos^2\phi_n-2\cos\psi_n\cos\phi_n\cos\phi_{n+1}+\cos^2\phi_{n+1}=\sin^2\psi_n$ gives
$G_n=A_n^2\sin^2\psi_n$. Divide \eqref{eq:Gid} by $G_n$: the \emph{secular}
(oscillation-averaged) part of $u_{n+1}u_{n+2}/G_n$ is $\dfrac{\cos\psi_n}{2\sin^2\psi_n}$, so
$\Delta\log G_n=(B_n-B_{n+1})\dfrac{\cos\psi_n}{2\sin^2\psi_n}+(\text{osc.})$. Since
$B_n=2\cos\psi_n$ we have $B_n-B_{n+1}=-\Delta B_n\approx 2\sin\psi_n\,\Delta\psi_n$, whence
$\Delta\log G_n\approx\cot\psi_n\,\Delta\psi_n$, integrating to
$\log G_n=\log\sin\psi_n+\mathrm{const}$, i.e.\ $G_n\propto\sin\psi_n=\tfrac12(4-B_n^2)^{1/2}$
and $A_n\propto(4-B_n^2)^{-1/4}$. Fix a compact window $W_\delta=\{\,2-|B_n|\ge\delta\,\}$
(so $\sin\psi_n\ge c_\delta>0$); the omitted oscillatory part has coefficient
$a_n=(B_n-B_{n+1})\cos(\phi_{n+1}+\phi_{n+2})/(2\sin^2\psi_n)$ with
$\sup_n|a_n|=O(\tau)$ and bounded variation $O(\tau)$ (since
$B_n-B_{n+1}=q^{1-\nu/2}x_n^2(1-q)=O(\tau)$ varies smoothly), tested against the
non-stationary phase $e^{i(\phi_{n+1}+\phi_{n+2})}$ whose partial sums are bounded by the
Dirichlet kernel $1/\sin(\psi_{\delta})=C_\delta$ on $W_\delta$; Abel summation then bounds the
accumulated contribution by $C_\delta(\sup|a_n|+\operatorname{Var}a_n)=O(\tau)$. This proves
\eqref{eq:env} on $W_\delta$. The gate evaluates $R$ at argument $x\approx1$, where
$B_n=2-x_n^2+O(\tau)$ sits a fixed $O(1)$ distance below $B=2$, i.e.\ interior to $W_\delta$
for a $\tau$-independent $\delta$; the turning-point neighbourhood $\psi_n\to0$ is never
encountered, and the constants are uniform in the pole index $m$.
\end{proof}

\begin{remark}[\checked]
\eqref{eq:rec} holds to $10^{-35}$, the Casoratian \eqref{eq:Gid}/(i) is constant to
$10^{-29}$ (spread $2.5\times10^{-29}$), and $A_n(4-B_n^2)^{1/4}$ is flat to relative
$0.024\approx\tau$ across the oscillatory bulk (\texttt{casoratian2.py},
\texttt{amplitude\_elem.py}).
\end{remark}

\subsection{Amplitude--phase decomposition and the bound}
For $x$ in the oscillatory regime the two $q$-Hankel solutions are complex conjugates,
$\Phi_{-}=\overline{\Phi_{+}}$ and $\tilde\kappa=\overline{\kappa}$ on the relevant locus,
so \eqref{eq:morita} reads $J^{(3)}_\nu=2\operatorname{Re}\!\big(\kappa\,\Phi_{+}\big)
=A(x)\cos\Theta(x)$ with amplitude $A=2|\kappa\,\Phi_+|$ and phase
$\Theta=\arg(\kappa\Phi_+)$. Splitting each factor into modulus and argument,
\begin{equation}\label{eq:ampphase}
   \frac{A(x)}{A_{\mathrm{cl}}(x)}
     =\frac{|\kappa|}{|\kappa_0|}\cdot\frac{|\Phi_+|}{|\Phi_{+,\mathrm{cl}}|},
   \qquad
   \Theta(x)-\Theta_{\mathrm{cl}}(x)
     =\big(\arg\kappa-\arg\kappa_0\big)+\big(\arg\Phi_+-\arg\Phi_{+,\mathrm{cl}}\big),
\end{equation}
where the subscript $\mathrm{cl}$ denotes the $\varepsilon\to0$ (classical Bessel) value.

\begin{theorem}[\proved\ \emph{modulo} the turning-point normalisation]\label{thm:R}
The connection-coefficient modulus and the oscillatory-bulk envelope contribute
$1+O(\tau)$ to the amplitude ratio $A(x)/A_{\mathrm{cl}}(x)$, uniformly in the pole index $m$
(at the evaluation point $x\approx1$, interior to the bulk); the phase difference is
$O(\sqrt\tau)$. \emph{Granting} the single-solution normalisation
$\gamma/\gamma_{\mathrm{cl}}=1+O(\tau)$ (the turning-point amplitude, of the same standing as
Lemma~\ref{lem:cos}), $A/A_{\mathrm{cl}}=1+O(\tau)$ and consequently
\[
   R=P_{12}-E=O(\tau^{5/2})\qquad\text{uniformly in }m,
\]
which establishes the gate condition \eqref{eq:gate-need}.
\end{theorem}

\begin{proof}
\emph{Amplitude.} By Lemma~\ref{lem:kappa}, the connection-coefficient modulus correction
$|\kappa|/|\kappa_0|=1+O(\tau)$ is $x$-independent. By Lemma~\ref{lem:caso}, the bulk envelope
\emph{shape} is $(4-B_n^2)^{-1/4}$ with $B_n=2-x_n^2+O(\tau)$, so
$|\Phi_+|/|\Phi_{+,\mathrm{cl}}|=(\gamma/\gamma_{\mathrm{cl}})\,(1+O(\tau))$. The conserved
Casoratian pins only this envelope \emph{shape}; it does \textbf{not} fix the single-solution
constant $\gamma$ (the Casoratian is bilinear, scaling with each solution). That constant is
set by the $x\to0$ connection---the classical confluent limit of \cite[Thm.~3.3]{Morita2011},
the lone input retained here, of the same standing as Lemma~\ref{lem:cos}. Granting
$\gamma/\gamma_{\mathrm{cl}}=1+O(\tau)$, both factors are $1+O(\tau)$, so $A/A_{\mathrm{cl}}
=1+O(\tau)$, uniform in $m$ since the connection-coefficient modulus is $x$-independent and the
bulk envelope is uniform on the fixed window $W_\delta$ containing the pole.

\emph{Phase.} $\arg\kappa-\arg\kappa_0=O(e^{-4\pi^2/\varepsilon})$ is exp-small by
\eqref{eq:kappa-asy}; the $q$-Hankel phase correction is the pole shift, which by
Lemma~\ref{lem:cos} is $-\tfrac{\sqrt2}{36}\sqrt\tau+O(\tau)$, so
$\Theta-\Theta_{\mathrm{cl}}=O(\sqrt\tau)$.

\emph{Assembly.} The pole is a \emph{zero} of $\cos\Theta_{\mathrm{cl}}$, so the leading
classical value is the \emph{sine} branch $E=A_{\mathrm{cl}}\sin\Theta_{\mathrm{cl}}$ (up to the
elementary prefactor of \S\ref{sec:lifting}). Writing $J^{(3)}_\nu=A\cos\Theta$ and expanding
about $\Theta_{\mathrm{cl}}$,
\[
   \frac{R}{E}=\frac{A}{A_{\mathrm{cl}}}\Big[\cos(\Theta-\Theta_{\mathrm{cl}})
     +\sin(\Theta-\Theta_{\mathrm{cl}})\cot\Theta_{\mathrm{cl}}\Big]-1 .
\]
Here $A/A_{\mathrm{cl}}=1+O(\tau)$ and $\cos(\Theta-\Theta_{\mathrm{cl}})-1=O(\tau)$; the cross
term is $O(\tau)$ since $\sin(\Theta-\Theta_{\mathrm{cl}})=O(\sqrt\tau)$ and, the pole sitting
at the sine extremum to $O(\sqrt\tau)$, $\cot\Theta_{\mathrm{cl}}=O(\sqrt\tau)$
(Lemma~\ref{lem:cos}). Hence $R/E=O(\tau)$; with $E=O(\tau^{3/2})$ uniformly,
$R=O(\tau^{5/2})$ uniformly in $m$.
\end{proof}

\begin{remark}[\checked]
The amplitude correction is $x$-independent and $O(\tau)$: $(A/A_{\mathrm{cl}}-1)/\tau$
evaluated on the $q$-Bessel at the classical extrema (where the $\sqrt\tau$ phase term
vanishes) is \emph{flat} across an $8\times$ range of argument (spread $\approx\tau$),
and end-to-end $J^{(3)}_{3/2}/J_{3/2}-1=O(\tau)$ with successive-difference ratio $\to2$
(\texttt{uniformity2.py}, \texttt{qbessel\_order.py}; these illustrative runs use nome $q$,
which scales the coefficient by a fixed factor versus the gate nome $q^2$ but leaves the
qualitative $O(\tau)$ unchanged). The \emph{assembled} gate constant
$|R|/\tau^{5/2}\to3.714$ uniformly to $m=200$ is computed at the actual gate object
(\texttt{u5b\_gate}). These confirm the bound numerically; the only unverified step
\emph{analytically} is the turning-point normalisation $\gamma/\gamma_{\mathrm{cl}}=1+O(\tau)$.
\end{remark}

\paragraph{Status and the lone residual.} Two ingredients are proved here from scratch:
Lemma~\ref{lem:kappa} (the modular/Poisson estimate of the connection-coefficient
\emph{modulus}) and Lemma~\ref{lem:caso} (the conserved-Casoratian bulk \emph{envelope}). What
this section therefore establishes \emph{unconditionally} is the reduction of the gate bound
\eqref{eq:gate-need} to a \emph{single} estimate: the single-solution amplitude normalisation
$\gamma/\gamma_{\mathrm{cl}}=1+O(\tau)$ across the turning point $\psi_n\to0$, equivalently the
$O(\tau)$ confluent correction to the $x\to0$ connection coefficient. This is the same input
flagged in Remark~\ref{rem:qbessel}, of the same standing as the $\mathrm{lem:cos}$ saddle
constant; its leading value (that $J^{(3)}_\nu\to$ classical Bessel) is Morita's confluent
limit \cite[Thm.~3.3]{Morita2011}, \cite{OldeDaalhuis1994}. Numerically the bound is certain
($|R|/\tau^{5/2}\to3.714$, $m\le200$); the residual is the \emph{proof} of that one $O(\tau)$
turning-point normalisation, a $q$-Airy/uniform-Bessel matching estimate not carried out here.
The two lemmas reduce the open atom to this single, named, classical-looking estimate.

\clearpage
% Route mod-p: the orthogonal (non-analytic) transcendence route for U via Christol.
% Self-contained section. Banks the maximal-non-automaticity evidence + the verified
% theta-telescoping that explains it, with honest proof status.
\section{The orthogonal route: $(u_n \bmod p)$ and Christol's theorem}
\label{sec:modp}

The transcendence of $U(x)=\sum_{n\ge0}u_nx^n$ over $\mathbb Q(x)$ established in
\S\ref{sec:transcendence} rests on the analytic gate (Lemma~\ref{lem:cos}). We
record here an \emph{independent} line of evidence that touches none of that
machinery, via reduction modulo a prime.

\subsection{Strategy}
Since $u_n\in\mathbb N$ are geodesic-growth counts, $U\in\mathbb Z[[x]]$.
If $U$ were algebraic over $\mathbb Q(x)$, then for all but finitely many primes
$p$ its reduction $U\bmod p$ would be algebraic over $\mathbb F_p(x)$, whence by
\textbf{Christol's theorem} the coefficient sequence $(u_n\bmod p)$ would be
$p$-automatic, i.e.\ its \emph{$p$-kernel}
\[
   \mathcal K_p(u)=\bigl\{\,(u_{p^k n+r})_{n\ge0}\;:\;k\ge0,\ 0\le r<p^k\,\bigr\}
\]
would be \emph{finite}. Contrapositively:
\begin{equation}\label{eq:christol-contra}
   \mathcal K_p(u)\ \text{infinite for one prime }p
   \quad\Longrightarrow\quad
   U\ \text{transcendental over }\mathbb Q(x),\ \text{unconditionally.}
\end{equation}
This needs no analytic input; it is a purely arithmetic/combinatorial criterion.

\subsection{The catalytic block has an explicit $\theta$-telescoping}
The relaxed bulk generating function obeys the catalytic functional
equation~\eqref{eq:catalytic}, which after isolating the dilation $t\mapsto q^2t$
(here $q=x^2$) reads
\begin{equation}\label{eq:modp-firstorder}
   F(q,t)=E(t)+D(t)\,F(q,q^2t),\qquad
   D(t)=\frac{2q(q-1)\,t}{(1-qt)(1-q^2t)},\quad
   E(t)=\frac{2qt}{1-qt}\bigl(1+F(q,q)\bigr).
\end{equation}
Since $F(q,0)=0$, iterating \eqref{eq:modp-firstorder} converges and telescopes.

\begin{proposition}[$\theta$-telescoping of the catalytic block]\label{prop:theta-tel}
With $(a;q)_m=\prod_{i=1}^{m}(1-aq^{i-1})$,
\[
   F(q,t)=\sum_{k\ge0}\alpha_k(t)\,E(q^{2k}t),
   \qquad
   \alpha_k(t)=\prod_{j=0}^{k-1}D(q^{2j}t)
             =\frac{[\,2q(q-1)\,]^{k}\,t^{k}\,q^{\,k(k-1)}}{(qt;q)_{2k}},
\]
where $(qt;q)_{2k}=\prod_{i=1}^{2k}(1-q^{i}t)$. The exponent $q^{k(k-1)}=x^{2k(k-1)}$
is a $\theta$\nobreakdash-/square factor; the Euler product $(qt;q)_{2k}$ in the
denominator (pentagonal-lacunary) redistributes each $\theta$-block over all
degrees.
\end{proposition}

\noindent\emph{Verification.} The product identity for $\alpha_k$ is a finite
$q$-Pochhammer manipulation ($\sum_{j=0}^{k-1}(2j+1)=k^2$ in the numerator power,
$\prod_{j<k}(1-q^{2j+1}t)(1-q^{2j+2}t)=(qt;q)_{2k}$ in the denominator), confirmed
symbolically for $0\le k\le4$; the resulting series reproduces the bulk record
$F(q,1)=2,2,6,2,18,6,42,18,118,50,282,190,706,594$ ($[q^{1}]\dots[q^{14}]$)
exactly. The full true-growth series $U$ is a finite rational assembly (two bulk
blocks, a geometric travel sum, four marker junctions) over $F$, so it inherits
this $\theta$ structure.

\subsection{Computational evidence: maximal non-automaticity}
We computed $(u_n\bmod p)$ by the validated polynomial-time catalytic engine
(a $1$-D collapse of the transfer, $O(N^4)$; reproduces $u_n\bmod p$ exactly for
$n\le42$). The $p$-kernel does not merely grow---it tracks the \emph{full
$p$-ary tree} with $O(1)$ deficit:
\[
\begin{array}{l|llll}
 & k=0 & k=1 & k=2 & k=3\\\hline
 \text{cumulative }|\mathcal K_3|\ (N=180) & 1 & 4 & 13 & 39\\
 \text{full ternary tree} & 1 & 4 & 13 & 40\\[2pt]
 \text{cumulative }|\mathcal K_5|\ (N=130) & 1 & 6 & 30 & -\\
 \text{full quinary tree} & 1 & 6 & 31 & -
\end{array}
\]
(per-level new sections vs.\ maxima: $p=3$, $+1,+3,+9,+26$ vs.\ $1,3,9,27$;
$p=5$, $+1,+5,+24$ vs.\ $1,5,25$). Both counts are robust to the section-overlap
threshold up to the data-length limit ($p=3$ at $k=3$ requires $\ge6$ shared
points per section, which $N=180$ supplies). The subword complexity is likewise
near-maximal ($p=3$: $3,9,27,75$ vs.\ $3^L=3,9,27,81$). This is the strongest
signature finite data affords: the $p$-kernel and factor complexity are as large
as possible, at two independent primes, exactly as
Proposition~\ref{prop:theta-tel} predicts. (By contrast $p=2$ is eventually
periodic---the dilation degenerates---so the route is specific to $p\ge3$.)

\subsection{Honest status}
By Christol, $\mathcal K_p(u)$ infinite $\iff$ $U\bmod p$ transcendental over
$\mathbb F_p(x)$, i.e.\ the explicit partial-$\theta$/Rogers--Ramanujan series of
Proposition~\ref{prop:theta-tel} is non-algebraic mod $p$. Finite kernel data
establishes that $\mathcal K_p(u)$ is \emph{large} (maximal through the deepest
reliable level), not provably \emph{infinite}. The square-support handle is
unavailable: $(u_n\bmod p)$ is dense and base-$p$ pseudorandom (the Euler-product
denominator fills the $\theta$ gaps), so the classical positive-characteristic
gap/zero-set theorems (Kedlaya; Adamczewski--Bell, \emph{Invent.\ Math.}\ 2012)
do not apply directly. Thus this route is genuinely \emph{orthogonal in method}
(automata over $\mathbb F_p(x)$, no analytic gate) and yields unconditional,
maximal evidence, while a proof reduces to the same underlying
$\theta$/square obstruction met by the analytic route and by
Route~D's $x^{k^2}$ specialization.


\clearpage
\section{Conclusion and open problems}\label{sec:conclusion}
The arithmetic of the reflection-group growth series A396406 separates cleanly. Its relaxed
companion $V$ and the catalytic building blocks are transcendental by elementary
(P\'olya--Carlson, absolute-contour) means. The true series $U$ is reduced to a single,
explicit, numerically-certain bound on the confluent residual of one Hahn--Exton $q$-Bessel
series, and three independent methods---the analytic turning-point analysis of
\S\ref{sec:amplitude}, the $q$-difference structure, and the mod-$p$ route of
\S\ref{sec:modp}---bottom out at the \emph{same} object: the $\nu=\tfrac32$ Hahn--Exton
$q$-Bessel turning-point connection, equivalently the $\theta$/square structure of the catalytic
telescoping. We isolate two clean open problems:
\begin{enumerate}[leftmargin=1.7em,itemsep=0.15em]
\item (\emph{analytic}) Prove the turning-point amplitude normalisation
$\gamma/\gamma_{\mathrm{cl}}=1+O(\tau)$ uniformly across $\psi_n\to0$---a $q$-Airy / uniform-Bessel
matching estimate---thereby making the transcendence of $U$ unconditional.
\item (\emph{arithmetic}) Prove $(u_n\bmod p)$ non-automatic for one odd prime $p$, equivalently
that the explicit partial-$\theta$ / Rogers--Ramanujan series $\Phi(q)$ of \S\ref{sec:modp} is
transcendental over $\mathbb F_p(q)$. The dense support of $u_n\bmod p$ obstructs every
gap-/lacunarity-based criterion; a new dense-support tool is required.
\end{enumerate}
Either resolves the transcendence of $U$. The constants the program produces (the bedrock
subleading $s_1=-0.0119596\ldots$, the connection coefficient $1891\sqrt2/10368$) are, on present
evidence, genuinely new.

\begin{thebibliography}{99}
\bibitem{BringmannMilas2015} K.~Bringmann and A.~Milas, \emph{$W$-algebras, false theta functions
and quantum modular forms}, Int. Math. Res. Not. IMRN \textbf{2015}, no.~21, 11351--11387.
\bibitem{CostinCostin1996} O.~Costin and R.~D.~Costin, \emph{Rigorous WKB for finite-order linear
recurrence relations with smooth coefficients}, SIAM J. Math. Anal. \textbf{27} (1996), 110--134.
\bibitem{FGD1995} P.~Flajolet, X.~Gourdon, and P.~Dumas, \emph{Mellin transforms and asymptotics:
harmonic sums}, Theoret. Comput. Sci. \textbf{144} (1995), 3--58.
\bibitem{FlajoletSedgewick} P.~Flajolet and R.~Sedgewick, \emph{Analytic Combinatorics},
Cambridge Univ. Press, 2009.
\bibitem{Morita2011} T.~Morita, \emph{A connection formula of the Hahn--Exton $q$-Bessel
function}, SIGMA \textbf{7} (2011), 115, 11~pp.; arXiv:1105.1998.
\bibitem{OEIS} N.~J.~A.~Sloane et al., \emph{The On-Line Encyclopedia of Integer Sequences},
\url{https://oeis.org}, sequences A396406 (this group) and A396927 (the orthoscheme family).
\bibitem{OldeDaalhuis1994} A.~B.~Olde Daalhuis, \emph{Asymptotic expansions for $q$-gamma,
$q$-exponential, and $q$-Bessel functions}, J. Math. Anal. Appl. \textbf{186} (1994), 896--913.
\bibitem{StanleyEC2} R.~P.~Stanley, \emph{Enumerative Combinatorics, Volume~2}, Cambridge Univ.
Press, 1999 (Coxeter growth series, Steinberg's formula).
\bibitem{WangWong2003} Z.~Wang and R.~Wong, \emph{Asymptotic expansions for second-order linear
difference equations with a turning point}, Numer. Math. \textbf{94} (2003), 147--194.
\bibitem{WongLi1992} R.~Wong and H.~Li, \emph{Asymptotic expansions for second-order linear
difference equations}, J. Comput. Appl. Math. \textbf{41} (1992), 65--94.
\bibitem{Christol1979} G.~Christol, \emph{Ensembles presque p\'eriodiques $k$-reconnaissables},
Theoret. Comput. Sci. \textbf{9} (1979), 141--145.
\bibitem{AdamczewskiBell2012} B.~Adamczewski and J.~P.~Bell, \emph{On vanishing coefficients of
algebraic power series over fields of positive characteristic}, Invent. Math. \textbf{187}
(2012), 343--393.
\bibitem{Carlson1921} F.~Carlson, \emph{\"Uber Potenzreihen mit ganzzahligen Koeffizienten},
Math. Z. \textbf{9} (1921), 1--13 (P\'olya--Carlson natural-boundary theorem).
\bibitem{IsolaPiergallini} S.~Isola and R.~Piergallini, \emph{On the generic triangle group},
arXiv:1405.1881 [math.MG], 2014.
\end{thebibliography}

\end{document}
